
\documentclass[10pt,letterpaper]{article}
\usepackage[T1]{fontenc}
\usepackage[utf8]{inputenc}
\usepackage{lmodern}
\usepackage{amsmath,amssymb,mathtools}
\usepackage{geometry}
\usepackage{enumitem}
\usepackage{tabularx}
\usepackage{tikz-cd}
\geometry{margin=0.86in}
\setlength{\parindent}{1.2em}
\setlength{\parskip}{0.12em}
\emergencystretch=3em
\setlist{leftmargin=2.1em,itemsep=0.08em,topsep=0.2em}
\newcommand{\C}{\mathbb C}
\newcommand{\Q}{\mathbb Q}
\newcommand{\Z}{\mathbb Z}
\newcommand{\F}{\mathbb F}
\newcommand{\Ql}{\mathbb Q_\ell}
\newcommand{\Zl}{\mathbb Z_\ell}
\newcommand{\A}{\mathbb A}
\newcommand{\G}{\mathbb G}
\newcommand{\Gal}{\operatorname{Gal}}
\newcommand{\GL}{\operatorname{GL}}
\newcommand{\Hom}{\operatorname{Hom}}
\newcommand{\Ind}{\operatorname{Ind}}
\newcommand{\Res}{\operatorname{Res}}
\newcommand{\Inf}{\operatorname{Inf}}
\newcommand{\detm}{\operatorname{det}}
\newcommand{\dimv}{\operatorname{dim}}
\newcommand{\Tr}{\operatorname{Tr}}
\newcommand{\Fr}{\operatorname{Frob}}
\newcommand{\rk}{\operatorname{rk}}
\newcommand{\ppcm}{\operatorname{ppcm}}
\newcommand{\Spec}{\operatorname{Spec}}
\newcommand{\iso}{\xrightarrow{\sim}}
\newcommand{\into}{\hookrightarrow}
\newcommand{\onto}{\twoheadrightarrow}

\providecommand{\Ker}{\operatorname{Ker}}
\providecommand{\coker}{\operatorname{coker}}
\providecommand{\Stab}{\operatorname{Stab}}
\providecommand{\Char}{\operatorname{Hom}}
\providecommand{\Frac}{\operatorname{Frac}}
\providecommand{\ord}{\operatorname{ord}}
\providecommand{\im}{\operatorname{im}}


\providecommand{\OO}{\mathcal O}
\providecommand{\mm}{\mathfrak m}
\providecommand{\ab}{\mathrm{ab}}
\providecommand{\geom}{\mathrm{geom}}
\providecommand{\cD}{\mathcal D}
\providecommand{\cF}{\mathcal F}
\providecommand{\Frob}{\operatorname{Frob}}
\providecommand{\vol}{\operatorname{vol}}


\providecommand{\R}{\mathbb R}
\providecommand{\OO}{\mathcal O}
\providecommand{\mm}{\mathfrak m}
\providecommand{\ab}{\mathrm{ab}}
\providecommand{\geom}{\mathrm{geom}}
\providecommand{\cD}{\mathcal D}
\providecommand{\cF}{\mathcal F}
\providecommand{\Frob}{\operatorname{Frob}}
\providecommand{\vol}{\operatorname{vol}}


\providecommand{\bbA}{\mathbb A}
\providecommand{\OO}{\mathcal O}
\providecommand{\mm}{\mathfrak m}
\providecommand{\ab}{\mathrm{ab}}
\providecommand{\TrL}{\operatorname{Tr}}
\providecommand{\Nm}{\operatorname{N}}
\providecommand{\sw}{\operatorname{sw}}
\providecommand{\cond}{\operatorname{cond}}
\providecommand{\vol}{\operatorname{vol}}

\providecommand{\spc}{\operatorname{sp}}

\providecommand{\bbA}{\mathbb A}
\providecommand{\Ga}{\mathbb G_a}
\providecommand{\Ql}{\mathbb Q_\ell}
\providecommand{\Zl}{\mathbb Z_\ell}
\providecommand{\OO}{\mathcal O}
\providecommand{\Spec}{\operatorname{Spec}}
\providecommand{\Gal}{\operatorname{Gal}}
\providecommand{\Hom}{\operatorname{Hom}}
\providecommand{\GL}{\operatorname{GL}}
\providecommand{\Aut}{\operatorname{Aut}}
\providecommand{\End}{\operatorname{End}}
\providecommand{\im}{\operatorname{im}}
\providecommand{\ss}{\mathrm{ss}}
\providecommand{\spc}{\operatorname{sp}}
\providecommand{\Lie}{\operatorname{Lie}}


\providecommand{\bbA}{\mathbb A}
\providecommand{\OO}{\mathcal O}
\providecommand{\LLL}{\mathcal L}
\providecommand{\Lie}{\operatorname{Lie}}
\providecommand{\im}{\operatorname{im}}
\providecommand{\ad}{\operatorname{ad}}
\providecommand{\Gr}{\operatorname{Gr}}
\providecommand{\Gal}{\operatorname{Gal}}
\providecommand{\GL}{\operatorname{GL}}
\providecommand{\Hom}{\operatorname{Hom}}
\providecommand{\Spec}{\operatorname{Spec}}
\providecommand{\Tr}{\operatorname{Tr}}
\providecommand{\Nm}{\operatorname{N}}
\providecommand{\Frob}{\operatorname{Frob}}
\providecommand{\ss}{\mathrm{ss}}
\providecommand{\Fr}{\operatorname{F}}


\providecommand{\bbA}{\mathbb A}
\providecommand{\bbF}{\mathbb F}
\providecommand{\bbZ}{\mathbb Z}
\providecommand{\bbQ}{\mathbb Q}
\providecommand{\OO}{\mathcal O}
\providecommand{\cG}{\mathcal G}
\providecommand{\cD}{\mathcal D}
\providecommand{\cH}{\mathcal H}
\providecommand{\Dparf}{D_{\mathrm{parf}}}
\providecommand{\Rr}{R\Gamma}
\providecommand{\Spec}{\operatorname{Spec}}
\providecommand{\Gal}{\operatorname{Gal}}
\providecommand{\Hom}{\operatorname{Hom}}
\providecommand{\End}{\operatorname{End}}
\providecommand{\Tr}{\operatorname{Tr}}
\providecommand{\Sym}{\operatorname{Sym}}
\providecommand{\detm}{\operatorname{det}}
\providecommand{\rank}{\operatorname{rank}}
\providecommand{\spc}{\operatorname{sp}}
\providecommand{\Fr}{\operatorname{F}}
\providecommand{\RHom}{\operatorname{RHom}}
\providecommand{\Coind}{\operatorname{Coind}}
\providecommand{\cent}{\operatorname{Cent}}


\providecommand{\RGamma}{\mathrm R\Gamma}
\providecommand{\bbA}{\mathbb A}
\providecommand{\OO}{\mathcal O}
\providecommand{\Gal}{\operatorname{Gal}}
\providecommand{\Tr}{\operatorname{Tr}}
\providecommand{\Ind}{\operatorname{Ind}}
\providecommand{\Ker}{\operatorname{Ker}}
\newcommand{\auditmark}[1]{\textsuperscript{\normalfont\scriptsize[E#1]}}
\begin{document}

\begin{center}
{\Large\bfseries LES CONSTANTES DES ÉQUATIONS FONCTIONNELLES DES FONCTIONS \(L\)}\\[2em]
{\large Par P. Deligne}\\[6em]
International Summer School on Modular Functions, Antwerp 1972\\[2em]
{\small Vérification du témoin et mises à jour éditoriales de cette révision :
OpenAI Codex, GPT-6 Sol, effort Ultra. Les transcriptions, traductions et
corrections antérieures sont documentées dans le dossier de provenance ;
aucune relecture humaine n'est revendiquée ici.}
\end{center}

\newpage

\begin{center}
\bfseries Sommaire
\end{center}

\noindent Les numéros renvoient à la pagination du texte source (Del-1--Del-97).
\medskip

\noindent\begin{tabularx}{\textwidth}{@{}p{0.04\textwidth}Xr@{}}
--& Introduction & 3\\
--& \S1. Représentations virtuelles d'un groupe fini & 7\\
--& \S2. Groupes de Weil & 21\\
--& \S3. Rappels sur les fonctions \(L\) & 26\\
--& \S4. Existence des constantes locales & 35\\
--& \S5. Formulaire & 48\\
--& \S6. Réduction modulo \(\ell\) des constantes locales & 54\\
--& \S7. Fonctions \(L\) modulaires & 58\\
--& \S8. Représentations \(\ell\)-adiques des groupes de Weil locaux non archimédiens; structures de Hodge et groupes de Weil locaux archimédiens & 66\\
--& \S9. Systèmes compatibles de représentations \(\ell\)-adiques & 74\\
--& \S10. La théorie de Grothendieck & 81\\
--& \S11. Appendice. Le calcul de \(\varepsilon\) modulo les racines de l'unité & 93\\
& Bibliographie & 96
\end{tabularx}

\newpage

\section*{Introduction}

Les résultats essentiels de cet article sont les suivants.

\paragraph{(A)} On donne une démonstration simple du théorème de Langlands [9], déjà prouvé au signe près par Dwork [6], qui permet d'écrire la constante des équations fonctionnelles des fonctions \(L\) d'Artin, ou de Weil [19], comme produit de constantes locales.

Le théorème à démontrer est local, et la démonstration de Langlands avait l'élégance d'être purement locale. Pour l'essentiel, il s'agit de démontrer une identité multiplicative entre sommes de Gauss un peu généralisées chaque fois qu'on a une relation \(\Z\)-linéaire entre caractères de groupes de Galois locaux induits par des caractères de représentations de dimension un de sous-groupes. Langlands construisait un ensemble générateur de telles relations, et ramenait les identités à démontrer à des identités entre sommes de Gauss qu'il démontrait en s'appuyant sur le théorème de Stickelberger.

La démonstration présentée ici est très proche de la démonstration [5] de l'identité de Hasse--Davenport (cf. aussi Weil [18]). Elle utilise un argument global (4.12), et le fait que les constantes locales à définir sont de nature essentiellement triviale à l'infini, dans le cas non ramifié, et dans certains cas très ramifiés (4.14).

\paragraph{(B)} Soient \(K\) un corps de fonctions et
\[
\rho:\Gal(\overline K/K)\longrightarrow \GL(n,\Q_\ell)
\]
une représentation \(\ell\)-adique du groupe de Galois de \(K\), presque partout non ramifiée. Grothendieck a montré que le produit eulérien étendu aux places de \(K\), avec \(p^{-s}\) remplacé par une indéterminée \(t\),
\[
Z(\rho,t)=\prod_v\det\bigl(1-F_vt^{\deg(v)},(\Q_\ell^n)^{\rho(I_v)}\bigr)^{-1}
\in\Q_\ell[[t]],
\]
est une fraction rationnelle, et qu'une équation fonctionnelle relie \(Z(\rho,t)\) à
\[
Z(\rho^\vee,t):
\qquad
Z(\rho^\vee,p^{-1}t^{-1})=\varepsilon\, Z(\rho,t),
\]
la constante \(\varepsilon\) étant un monôme \(at^b\).\auditmark{19} Cette théorie est résumée au §10. Le résultat (A) suggère une formule pour exprimer la constante \(\varepsilon\) comme produit de constantes locales. Nous montrons que cette formule est vraie lorsque \(\rho\) appartient à un système compatible infini de représentations \(\ell\)-adiques (9.3 et 9.9). D'après Weil [16] et Jacquet--Langlands [8], cette formule, appliquée au système compatible de représentations \(\ell\)-adiques défini par une courbe elliptique sur \(K\), permet d'associer une forme modulaire pour \(\GL(2,K)\) à toute courbe elliptique sur \(K\).

Il serait très intéressant de savoir si la formule obtenue reste valable pour une quelconque représentation \(\ell\)-adique.

Voici le contenu des divers paragraphes, et leur interdépendance.

Le début du §1 rassemble quelques résultats sur les représentations induites des groupes finis qui nous seront utiles.

La fin du §1, de 1.11 à la fin, ne sert pas dans le reste de l'article. Pour \(G\) un groupe résoluble, on y décrit un ensemble \(R\) de relations linéaires sur \(\Z\) entre les caractères de \(G\) induits par des caractères de dimension un de sous-groupes, et on montre que toute telle relation est combinaison linéaire d'éléments de \(R\). Le résultat, et sa démonstration, sont contenus implicitement dans [9].

Les §§2 et 3 servent essentiellement à fixer les notations. Au §2, on rappelle la structure du groupe de Galois d'un corps local [10], on définit les groupes de Weil, variante des groupes de Galois, et on donne quelques énoncés de la théorie du corps de classes. Le lecteur prendra garde que, dans tout cet article, l'isomorphisme de la théorie du corps de classes local transforme uniformisante en inverse de la substitution de Frobenius; la raison est donnée en 3.6. Le §3 est un rappel de la thèse de Tate, et de Artin [1].

Le §4 contient la démonstration du résultat annoncé en (A) ci-dessus. Le formulaire du §5 donne une liste de propriétés fondamentales des constantes locales construites au §4, et quelques identités entre sommes de Gauss qui en résultent.

Les §§6, 7 et 8 (qui ne dépend que du §2) sont préliminaires à la démonstration de (B), donnée au §9. Dans les énoncés des §§3 et 4, nous avions été très attentifs à n'utiliser que des formules rationnelles, ne contenant ni passage à la limite ni racine carrée, fût-ce d'un entier positif. Nous en récoltons le fruit aux §§6 et 7. Au §6, pour \(K\) un corps local non archimédien et \(V\) une représentation modulaire de \(W(\overline K/K)\) sur un corps \(\Lambda\) de caractéristique \(\ell\ne p\), nous définissons, par transport de structure pour \(\ell=0\) et par réduction modulo \(\ell\) pour \(\ell\ne0\), une constante locale
\[
\varepsilon(V,\psi,dx)\in \Lambda^*.
\]
Au §7, pour \(K\) global de caractéristique \(p\) et \(V\) une représentation modulaire de \(W(\overline K/K)\), nous prouvons par réduction modulo \(\ell\) une équation fonctionnelle pour la fonction \(L\in \Lambda(t)\) correspondante.

Au §8, nous donnons des classes d'isomorphie de représentations \(\ell\)-adiques du groupe de Weil d'un corps local une description purement algébrique, indépendante de la topologie de \(\Q_\ell\). Ceci sert à définir la compatibilité de représentation \(\ell\) et \(\ell'\)-adique. La notion introduite est plus fine que celle de Serre [12].

Le §9 donne la démonstration de (B) (9.3 et 9.9). L'idée est que pour démontrer une identité, il suffit de la démontrer modulo \(\ell\) pour une infinité de \(\ell\). Nous donnons aussi, pour les corps de fonctions, une réponse partielle à une question de Serre en montrant (9.8) que deux représentations \(\ell\) et \(\ell'\)-adiques donnant lieu au même polynôme caractéristique de Frobenius, supposé dans \(\Q[t]\), en presque toutes les places, vérifient une compatibilité aux places résiduelles.

Le §10 résume la théorie de Grothendieck des fonctions \(L\) (cf. [7] et SGA 5), qui nous a servi de façon cruciale au §9. En 10.11, je tente d'expliquer le sel de sa méthode. Les résultats farfelus 10.13 résultent aussitôt de ce qui précède. J'ignore leur signification.

\section*{\S1. Représentations virtuelles d'un groupe fini}

\paragraph{1.1.} Soit \(K\) un corps commutatif. Le cas le plus important pour nous sera celui où \(K\) est algébriquement clos de caractéristique \(0\), par exemple \(K=\C\). Soit \(G\) un groupe. On notera \(R_K(G)\) le groupe de Grothendieck de la catégorie des \(K[G]\)-modules de dimension finie sur \(K\). Il s'identifie au \(\Z\)-module libre de base l'ensemble des classes d'isomorphie de représentations linéaires \(K\)-irréductibles de \(G\) sur \(K\). C'est un anneau commutatif à unité, pour le produit tensoriel sur \(K\), et même un \(\lambda\)-anneau au sens de Grothendieck. Ses éléments sont les représentations virtuelles de \(G\) sur \(K\).

Pour définir un homomorphisme \(f\) de \(R_K(G)\) dans un groupe commutatif \(X\), il suffit:
\begin{enumerate}[label=\alph*)]
\item de définir \(f(V)\), pour \(V\) une représentation de \(G\) sur \(K\);
\item de vérifier que pour toute suite exacte de représentations
\[
0\longrightarrow V'\longrightarrow V\longrightarrow V''\longrightarrow0
\]
on a \(f(V)=f(V')+f(V'')\).
\end{enumerate}
Appliquant cette règle, on définit la dimension d'une représentation virtuelle
\[
\dim:R_K(G)\longrightarrow\Z,
\]
et son déterminant
\[
\det:R_K(G)\longrightarrow\Hom(G,K^*).
\]

Si \(H\) est un sous-groupe d'indice fini de \(G\), on définit l'induction
\[
\Ind_H^G:R_K(H)\longrightarrow R_K(G)
\]
comme correspondant à l'extension des scalaires
\[
V\longmapsto K[G]\otimes_{K[H]}V,
\]
foncteur exact; il passe donc aux représentations virtuelles. Pour \(f:G\to H\), on définit la restriction
\[
\Res_G^H:R_K(H)\longrightarrow R_K(G)
\]
comme restriction des scalaires de \(K[H]\) à \(K[G]\). Pour \(f\) surjectif, on parle parfois plutôt d'inflation.

Pour tout groupe \(G\), on note \(G^{\mathrm{ab}}\) le plus grand quotient abélien de \(G\). Un \(K\)-quasi-caractère \(\chi\) de \(G\) est un homomorphisme
\[
\chi:G\longrightarrow K^*;
\]
on l'identifie au \(K\)-quasi-caractère de \(G^{\mathrm{ab}}\) par lequel il se factorise, et on note \([\chi]\) la représentation de dimension \(1\) correspondante. Pour \(G\) fini, nous emploierons indifféremment les mots caractère et quasi-caractère. Caractère ne signifiera jamais caractère d'une représentation de dimension \(>1\).

The following proposition is due to P. Gallagher [20].

\paragraph{Proposition 1.2.} Soient \(G\) un groupe, \(H\) un sous-groupe d'indice fini et
\[
t:G^{\mathrm{ab}}\longrightarrow H^{\mathrm{ab}}
\]
le transfert. Pour \(\rho\) une représentation virtuelle de dimension \(0\) de \(H\) et \(x\in G^{\mathrm{ab}}\), on a
\[
\det(\Ind_H^G(\rho))(x)=\det(\rho)(t(x)).
\]
On montrera, à peine plus généralement, que si \(\varepsilon\) est le déterminant de la représentation de permutation de \(G\) sur \(G/H\),
\[
\varepsilon:G\longrightarrow \mathfrak S_{G/H}\longrightarrow \{\pm1\},
\]
on a, pour \(\rho\) de dimension quelconque,
\[
\det(\Ind_H^G(\rho))(x)=\varepsilon^{\dim\rho}\det(\rho)(t(x)).
\]

C'est assez évident en termes topologiques. Soient \(B_G\) le classifiant de \(G\) et \(B_H\) celui de \(H\), vu comme revêtement à \([G:H]\) feuillets de \(B_G\) : \(f:B_H\to B_G\). Les représentations de \(G\) s'identifient aux systèmes locaux de \(K\)-vectoriels de dimension finie sur \(B_G\), et \(\Ind_H^G\) au foncteur \(f_*\). Le transfert \(H_1(B_G)\to H_1(B_H)\) est transposé aux morphismes trace \(f_!:H^1(B_H,X)\to H^1(B_G,X)\), où \(X\) est un groupe abélien.\auditmark{25} Pour \(X=K^*\), \(f_!\) correspond à la norme : si \(\eta\in H^1(B_H,K^*)\) est la classe d'un système local \(L\) de rang \(1\) de \(K\)-vectoriels, \(f_!(\eta)\) est la classe de \(N(L)\), système local sur \(B_G\) qui, localement sur \(B_G\), est le produit tensoriel des images réciproques de \(L\) par les \([G:H]\) sections disjointes.

Posons
\[
\det V=\bigwedge^{\dim V}V
\]
pour \(V\) un système local, et soit \(\varepsilon\) le système local \(\det f_*K\). La proposition 1.2 résulte de l'existence d'un isomorphisme canonique, pour \(V\) système local sur \(B_H\),
\[
\det f_*V\simeq \varepsilon^{\otimes\dim V}\otimes N(\det V).
\]
La construction de cet isomorphisme est locale sur \(B_G\); elle se ramène à la construction suivante.

\paragraph{Construction 1.3.} Pour \((V_i)_{i\in I}\) une famille finie de \(K\)-vectoriels de dimension \(d\), et \(\varepsilon=\det(K^I)\), on a canoniquement
\[
\det\Bigl(\bigoplus_{i\in I}V_i\Bigr)\simeq
\varepsilon^{\otimes d}\otimes\bigotimes_{i\in I}\det(V_i).
\]
La flèche est, pour \(i_1,\ldots,i_N\) la suite des éléments de \(I\),
\[
\begin{aligned}
&(e_{i_1,1}\wedge\cdots\wedge e_{i_1,d})\wedge\cdots\wedge
(e_{i_N,1}\wedge\cdots\wedge e_{i_N,d}) \\
&\quad\longmapsto
(i_1\wedge\cdots\wedge i_N)^{\otimes d}\otimes
(e_{i_1,1}\wedge\cdots\wedge e_{i_1,d})\otimes\cdots\otimes
(e_{i_N,1}\wedge\cdots\wedge e_{i_N,d}).
\end{aligned}
\]

\paragraph{Parenthèse 1.4.} Plutôt qu'un classifiant \(BG\), on eût plus intrinsèquement pu prendre, dans le raisonnement précédent, le topos classifiant \(B_G\) de Grothendieck (SGA 4, IV, 2.4, p. 315, pour \(E=(\mathrm{Ens})\)).

Supposons \(G\) fini. Nous dirons que \(K\) est assez gros s'il contient toutes les racines de l'unité d'ordre divisant le ppcm des ordres des éléments de \(G\). Un groupe \(H\) sera dit \(p\)-élémentaire s'il est produit d'un \(p\)-groupe par un groupe cyclique; élémentaire s'il est \(p\)-élémentaire pour un nombre premier \(p\) convenable. Nous aurons besoin de la variante suivante du théorème de Brauer sur les caractères induits (qu'on obtient en omettant partout ``de dimension zéro'').

\paragraph{Proposition 1.5.} Si \(K\) est assez gros, toute représentation virtuelle de dimension \(0\) de \(G\) est somme de représentations virtuelles
\[
\Ind_H^G(\chi),
\]
avec \(\chi\) virtuel de dimension \(0\), \(H\) élémentaire et \(\chi\) combinaison linéaire de représentations de dimension \(1\) de \(H\).

D'après le théorème de Brauer appliqué à la représentation triviale \(1_G\) de \(G\), il existe une décomposition
\[
1_G=\sum_H \Ind_H^G(z_H),
\qquad H\text{ élémentaire}.
\]
Pour \(y\) représentation virtuelle de dimension \(0\) de \(G\), on a alors
\[
y=\sum_H y\cdot \Ind_H^G(z_H)
  =\sum_H\Ind_H^G(\Res_H^G(y)\cdot z_H),
\]
avec \(\dim(\Res_H^G(y)\cdot z_H)=0\). Ceci nous ramène à supposer \(G\) élémentaire.

Par linéarité, on peut supposer \(y\) de la forme
\[
\rho-\dim(\rho)\,1_G,
\]
avec \(\rho\) irréductible, donc, puisque \(G\) est nilpotent, induite par un caractère \(\chi\) d'un sous-groupe \(H\) de \(G\) contenant le centre \(Z\). On a alors
\[
\rho-\dim(\rho)\,1_G
=
\Ind_H^G(\chi-1_H)+\bigl(\Ind_H^G(1_H)-[G:H]1_G\bigr).
\]
Le premier terme est du type voulu, et le second est l'inflation d'une représentation virtuelle de dimension \(0\) de \(G/Z\). On conclut par récurrence sur l'ordre de \(G\), en notant que \(|G/Z|<|G|\) si \(G\ne\{e\}\).

\paragraph{Scholie 1.6.} Un homomorphisme
\[
f:R_K(G)\longrightarrow X
\]
est connu quand on connaît sa valeur en \(1_G\) et en les
\[
\Ind_H^G(\chi-1_H),
\]
pour \(H\) élémentaire et \(\chi\) un caractère de \(H\).

\paragraph{1.7.} Soit \(A\) un anneau de valuation discrète d'inégale caractéristique \(p\), de corps des fractions \(K\) et de corps résiduel \(k=A/\mathfrak m\). Rappelons la définition de l'homomorphisme de décomposition ([11], III, 2.2)
\[
d:R_K(G)\longrightarrow R_k(G):
\]
pour \(V\) une représentation de \(G\) sur \(K\) et \(E\) un réseau \(G\)-invariant,
\[
d([V])=[E\otimes_Ak].
\]




\paragraph{Proposition 1.8.} On suppose \(K\) assez gros. Le noyau de l'homomorphisme de décomposition \(d\) est alors engendré par les
\[
\Ind_H^G\bigl([\chi']-[\chi'']\bigr)
\]
pour \(H\) un sous-groupe élémentaire et \(\chi',\chi''\in\Hom(H,K^*)\) congrus modulo \(\mathfrak m\).

L'homomorphisme d'anneaux \(d\) commutant à l'induction, le raisonnement de 1.5 nous ramène au cas où \(G\) est élémentaire, donc produit d'un \(p\)-groupe \(P\) par un groupe \(H\) d'ordre premier à \(p\). Les corps \(K\) et \(k\) étant assez gros, on a
\[
R_K(G)=R_K(P)\otimes R_K(H),
\]
et de même pour les anneaux de représentations sur \(k\). Pour un groupe d'ordre premier à \(p\), \(d\) est bijectif. On a donc\auditmark{26}
\[
\operatorname{Ker}\bigl(R_K(G)\to R_k(G)\bigr)
=\operatorname{Ker}\bigl(R_K(P)\to R_k(P)\bigr)\otimes R_K(H).
\]
Appliquant le théorème de Brauer à \(H\), on voit qu'il suffit de prouver 1.8 pour un \(p\)-groupe. Ce cas est trivial, car tout caractère \(\chi\in\Hom(P,K^*)\) d'un \(p\)-groupe est congru à \(1\) modulo \(\mathfrak m\).

\paragraph{1.9.} Dans la fin de ce numéro, on suppose que \(K=\C\), on ne considère que des groupes finis, et on abrège \(R_\C(G)\) en \(R(G)\). Pour \(G\) commutatif, on note \(G^*\) le groupe des caractères de \(G\).

Soit \(R^+(G)\) le \(\Z\)-module libre de base l'ensemble des classes de \(G\)-conjugaison de couples \((H,\chi)\), où \(H\) est un sous-groupe de \(G\) et \(\chi\) un caractère de \(H\). On définit dans \(R^+(G)\) une multiplication par la règle
\[
\tag{1.9.1}
(A,\alpha)\,(B,\beta)=
\sum_{(x,y)\in G\backslash(G/A\times G/B)}
\bigl(A^x\cap B^y,\,
(\alpha^x|_{A^x\cap B^y})(\beta^y|_{A^x\cap B^y})\bigr),
\]
où \((x,y)\) parcourt des représentants des orbites de \(G\) dans \(G/A\times G/B\), \(A^x=xAx^{-1}\), et \(\alpha^x\) est le caractère \(a\mapsto \alpha(x^{-1}ax)\) de \(A^x\). Pour cette multiplication, \(R^+(G)\) est un anneau commutatif d'unité \((G,1)\), et on définit un homomorphisme d'anneaux
\[
\tag{1.9.2}
b:R^+(G)\longrightarrow R(G),\qquad (H,\chi)\longmapsto \Ind_H^G(\chi).
\]
Cet homomorphisme est surjectif d'après le théorème de Brauer. Pour \(\chi\) un caractère de \(G\), la multiplication par \((G,\chi)\) se note \(\cdot\chi\) et s'appelle torsion par \(\chi\):
\[
\tag{1.9.3}
(H,\varphi)\cdot\chi=(H,\varphi\cdot\chi|_H).
\]
Pour \(G'\subset G\), on définit l'induction
\[
\tag{1.9.4}
\Ind_{G'}^G:R^+(G')\longrightarrow R^+(G),\qquad (H,\chi)\longmapsto(H,\chi).
\]
Pour \(u:G\to G'\) un morphisme, on définit la restriction
\[
\tag{1.9.5}
\Res_u:R^+(G')\longrightarrow R^+(G),\qquad (H,\chi)\longmapsto
\sum_{s\in u(G)\backslash G'/H}
\bigl(u^{-1}(sHs^{-1}),\,\chi^s\circ u\bigr),
\]
où \(s\) parcourt un ensemble de représentants des doubles classes. On dispose de diagrammes commutatifs
\[
\begin{tikzcd}[column sep=3.0em,row sep=2.0em]
R^+(G') \arrow[r,"\Ind"] \arrow[d,"b"'] & R^+(G) \arrow[d,"b"]\\
R(G') \arrow[r,"\Ind"'] & R(G)
\end{tikzcd}
\qquad
\begin{tikzcd}[column sep=3.0em,row sep=2.0em]
R^+(G') \arrow[r,"\Res"] \arrow[d,"b"'] & R^+(G) \arrow[d,"b"]\\
R(G') \arrow[r,"\Res"'] & R(G),
\end{tikzcd}
\]
pour le second cf. [11], p. 75. On a la formule, pour \(R\in R^+(G)\),
\[
\tag{1.9.6}
R\cdot(H,\chi)=\Ind_H^G\bigl(\Res_H^G(R)\cdot\chi\bigr).
\]
On parle d'inflation lorsque \(u\) est surjectif :
\[
\tag{1.9.7}
\operatorname{Infl}:R^+(G')\longrightarrow R^+(G),\qquad
(H,\chi)\longmapsto\bigl(u^{-1}(H),\chi\circ u\bigr).
\]

\paragraph{Notation 1.9.8.} Pour toute fonction additive de représentations \(F\), on note de même la fonction qu'elle définit sur \(R(G)\), et le composé \(F\circ b\). Ainsi
\[
\dim(v)=\dim(b(v)),\quad\ldots
\]

\paragraph{Variantes 1.10.} (i) On note \(R^+_0(G)\) le sous-groupe de \(R^+(G)\) engendré par les \((H,\chi)-(H,1)\). Pour \((H,\chi)\) parcourant la base canonique de \(R^+(G)\), à l'exception des \((H,1)\), ces éléments forment une base de \(R^+_0(G)\). On vérifie que \(R^+_0(G)\) est un idéal de \(R^+(G)\). Soit \(R^0(G)\subset R(G)\) l'ensemble des représentations virtuelles de dimension \(0\) de \(G\). D'après 1.5, l'application naturelle, induite par (1.9.2), de \(R^+_0(G)\) dans \(R^0(G)\) est surjective.

(ii) Pour \(G\) un groupe topologique, fini ou non, on note \(R^+(G)\) et \(R^+_0(G)\) les groupes analogues aux précédents obtenus en se limitant aux sous-groupes fermés \(H\) d'indice fini et aux quasi-caractères continus. Les fonctorialités 1.9 se généralisent mutatis mutandis.

La fin de ce paragraphe ne servira plus par la suite. Nous y exposons des résultats tirés de Langlands [9].

\paragraph{Lemme 1.11.} Supposons que \(G=H\cdot C\), avec \(C\) commutatif distingué, et soit \(\chi\) un caractère de \(H\). Pour tout \(\mu\in C^*\), soit \(G_\mu\subset G\) le fixateur de \(\mu\) par conjugaison. Si
\[
\chi|_{H\cap C}=\mu|_{H\cap C},
\]
on note \(\{\chi,\mu\}\) le caractère de \(G_\mu=(H\cap G_\mu)\cdot C\) qui prolonge \(\chi|_{H\cap G_\mu}\) et \(\mu\). On a
\[
\Ind_H^G(\chi)=
\sum_{\substack{\mu\in C^*/G\\ \chi|_{H\cap C}=\mu|_{H\cap C}}}
\Ind_{G_\mu}^G\{\chi,\mu\}.
\]
La vérification est laissée au lecteur.

\paragraph{1.12.} Un élément \(R\) de \(\Ker(R^+(G)\to R(G))\) sera dit de type I, II ou III s'il existe un quotient \(A\) d'un sous-groupe \(B\) de \(G\) tel que \(R\) s'obtienne à partir d'un élément \(S\) de l'un des types respectifs suivants de \(\Ker(R^+(A)\to R(A))\), par inflation de \(A\) à \(B\), torsion par un caractère de \(B\), et induction de \(B\) à \(G\). On désigne par \(\ell\) un nombre premier.

\emph{I.} \(A\simeq\Z/\ell\Z\), et
\[
 S=(e,[1])-\sum_{\chi\in A^*}(A,\chi).
\]

\emph{II.} \(A\) est une extension centrale de \((\Z/\ell\Z)^2\) par un groupe abélien \(Z\). Soient \(H_i\) (\(i=1,2\)) l'image réciproque dans \(A\) des premier et second facteurs \(\Z/\ell\Z\), et \(\chi_i\) un caractère de \(H_i\). On suppose que \(\chi_1|_Z=\chi_2|_Z\), et que ce caractère de \(Z\) est non trivial sur un commutateur.
\[
\begin{tikzcd}[column sep=3.2em,row sep=2.7em]
Z \arrow[r,hook] \arrow[d,hook] & H_1 \arrow[d,hook]\\
H_2 \arrow[r,hook] & A
\end{tikzcd}
\]
La relation \(S\) est
\[
S=(H_1,\chi_1)-(H_2,\chi_2).
\]

\emph{III.} \(A\) est un produit semi-direct \(H\cdot C\), \(\chi\) est un caractère de \(H\), et le sous-groupe distingué \(C\) est minimal parmi les sous-groupes commutatifs distingués non triviaux de \(A\). Pour tout caractère \(\mu\) de \(C\), \(A_\mu\) est le fixateur de \(\mu\), et \(\{\chi,\mu\}\) est le caractère de \(A_\mu\) qui prolonge \(\chi|_{H\cap A_\mu}\) et \(\mu\).\auditmark{27} La relation \(S\) exprime
\[
\Ind_H^A(\chi)=
\sum_{\mu\in C^*/A}
\Ind_{A_\mu}^A(\{\chi,\mu\}),
\]
où \(\mu\) parcourt un ensemble de représentants des orbites de \(A\) dans \(C^*\).

Ces relations sont des cas particuliers de 1.11: pour I, \(H=\{e\}\) et \(C=A\); pour II, \(H=H_1\) et \(C=H_2\); pour III, \(H\) et \(C\) sont ceux indiqués.

Toute relation déduite par induction, inflation ou torsion d'une relation de type I, respectivement II ou III, est encore du même type.

\paragraph{Théorème 1.13 (Langlands [9], §18).} Si \(G\) est nilpotent, le noyau de
\[
b:R^+(G)\longrightarrow R(G)
\]
est engendré, comme \(\Z\)-module, par les relations de type I et II.

\paragraph{Lemme 1.13.1.} Si \(G\) est commutatif, \(\Ker(b)\) est engendré par les relations de type I.

Il faut montrer que, pour tout sous-groupe \(H\) de \(G\) et tout \(\chi\in H^*\), la relation
\[
\tag{1}
\Ind_H^G(\chi)=
\sum_{\substack{\varphi\in G^*\\ \varphi|_H=\chi}}[\varphi]
\]
est conséquence (combinaison linéaire) de relations de type I. Celles-ci sont les suivantes : pour \(H'\subset H''\subset G\), avec \([H'':H']\) premier, et \(\chi\) un caractère de \(H'\), toujours induit par un caractère de \(G\),
\[
\Ind_{H'}^G(\chi)=
\sum_{\substack{\varphi\in (H'')^*\\ \varphi|_{H'}=\chi}}
\Ind_{H''}^G(\varphi),
\]
et le fait que (1) en résulte se voit en prenant une suite de Jordan--Hölder de \(G/H\).

\paragraph{Lemme 1.13.2.} Soient \(C\) un sous-groupe commutatif distingué de \(G\), \(T\) un ensemble de représentants de \(C^*/G\), \(\mu\in T\), et \(G_\mu\) le fixateur de \(\mu\). Soit \((H_i)\) une famille de sous-groupes de \(G\) contenant \(C\), et \(\chi_i\) un caractère de \(H_i\), avec \(\chi_i|_C\in T\). Si la relation
\[
\sum_i n_i\,\Ind_{H_i}^G(\chi_i)=0
\]
est vraie dans \(R(G)\), alors la relation
\[
\sum_{\chi_i|_C=\mu}
n_i\,\Ind_{H_i}^{G_\mu}(\chi_i)=0
\]
est vraie dans \(R(G_\mu)\). La vérification est laissée au lecteur.

Soit \(Z\) le centre de \(G\), et prouvons 1.13 par récurrence sur l'ordre de \(G/Z\). D'après 1.13.1, on peut supposer \(G\) non commutatif. Soit donc \(C\) un sous-groupe commutatif distingué contenant \(Z\), tel que \([C:Z]=\ell\) soit premier, et d'image centrale dans \(G/Z\).

Soit une relation
\[
\tag{1}
\sum_i n_i\,\Ind_{H_i}^G(\chi_i)=0.
\]
Elle est conséquence de relations induites par les relations
\[
\tag{$a_i$}
\Ind_{H_i}^{H_iZ}(\chi_i)=
\sum_{\substack{\chi_i'\in(H_iZ)^*\\ \chi_i'|_{H_i}=\chi_i}}[\chi_i']
\]
et d'une relation (2) analogue à (1), avec \(H_i\supset Z\). Les relations \((a_i)\) sont combinaisons linéaires de relations de type I; on peut le déduire de 1.13.1 appliqué à \(H_iZ/\Ker(\chi_i)\).

La relation (2) est conséquence de relations induites par les relations 1.11 pour \(HC,H,C,\chi\), avec \(H\supset Z\),
\[
\tag{b}
\Ind_H^{HC}(\chi)=
\sum_{\substack{\mu\in C^*/H\\ \chi|_{H\cap C}=\mu|_{H\cap C}}}
\Ind_{(HC)_\mu}^{HC}\{\chi,\mu\},
\]
et d'une relation (3) analogue à (1), avec cette fois \(H_i\supset C\). Appliquons 1.13.2 à (3), pour l'exprimer comme somme de relations induites par des relations
\[
\tag{4}
\sum_i n_i\,\Ind_{H_i}^{G_\mu}(\chi_i)=0,
\]
avec \(H_i\supset C\) et \(\chi_i|_C=\mu\). L'image de \(C\) dans \(G_\mu/\Ker(\mu)\) est centrale, car \(\mu\), invariant par \(G_\mu\), est injectif sur \(C/\Ker(\mu)\). L'hypothèse de récurrence s'applique donc à \(G_\mu\), et les relations (4) sont du type voulu.

Prouvons que les relations (b) le sont. Si \(HC\ne G\), cela résulte de l'hypothèse de récurrence appliquée à \(HC\). Si \(HC=G\), \(H\) est distingué, car il contient \(Z\) et \(C\) est central modulo \(Z\). Soit \(A\) l'intersection des noyaux des conjugués de \(\chi\). La relation (b) est l'inflation d'une relation analogue pour \(G/A\). Si l'hypothèse de récurrence s'applique à \(G/A\), on a gagné. Sinon, on se retrouve dans la même situation qu'avant, avec \(A=\{e\}\), et on conclut par le lemme suivant.

\paragraph{Lemme 1.13.3.} Avec les notations précédentes, si \(A=\{e\}\), la relation (b) est de la forme II.

Le groupe \(H\) est commutatif, car des caractères séparent ses éléments. Le groupe \(Z\) est cyclique, car un seul caractère sépare ses éléments. Pour \(c\in C\) et \(h\in H\), posons
\[
u(c,h)=chc^{-1}h^{-1}.
\]
C'est un élément de \(Z\), car \(C\) est central modulo \(Z\). Il dépend biadditivement de \(c\) et \(h\), et définit
\[
\bar u:C/Z\otimes H/Z\longrightarrow Z.
\]
Soit \(c\) engendrant \(C/Z\). Puisque \(Z\) est le centre et que \(G\ne Z\), \(\bar u(c,\ ) :H/Z\to Z\) est injectif, d'image cyclique tuée par \(\ell\) et non triviale. Dès lors \(|C/Z|=|H/Z|=\ell\), et \(G\) est extension centrale de \(C/Z\times H/Z\) par \(Z\). Le caractère \(\chi\) est distinct d'au moins un de ses conjugués, sinon \(H\) serait central; il est donc non trivial sur un commutateur, et on est dans la situation II.

\paragraph{Théorème 1.14 (Langlands [9]).} Si \(G\) est résoluble, le noyau de
\[
b:R^+(G)\longrightarrow R(G)
\]
est engendré par les relations de type I, II et III.

\paragraph{Lemme 1.14.1.} Si \(G\) est résoluble, le sous-\(\Z\)-module \(N(G)\) de \(R^+(G)\) engendré par les relations de type I, II et III est un idéal.

Prouvons 1.14 et 1.14.1 par une récurrence simultanée sur l'ordre de \(G\).

(A) Le théorème 1.14 pour les sous-groupes propres de \(G\) implique 1.14.1 pour \(G\). Soient \(R\in R^+(G)\) de type I, II ou III, \(H\) un sous-groupe et \(\chi\) un caractère de \(G\).\auditmark{28} Prouvons que
\[
R\cdot(H,\chi)\in N(G).
\]
Si \(H=G\), cela résulte de ce que \(N(G)\) est stable par torsion par un caractère de \(G\). Si \(H\ne G\), on a, par (1.9.6),
\[
R\cdot(H,\chi)=\Ind_H^G(\Res_H^G(R)\cdot\chi).
\]
Vu l'hypothèse, \(\Res_H^G(R)\in N(H)\), et la torsion puis l'induction donnent un élément de \(N(G)\).

(B) Le lemme 1.14.1 pour \(G\), plus le théorème 1.14 pour les groupes d'ordre plus petit, impliquent 1.14 pour \(G\). Soit \(Z\) le centre de \(G\). D'après 1.13, on peut supposer que \(G\) n'est pas nilpotent. Distinguons deux cas.

(B1) \(Z\ne\{e\}\). D'après le théorème de Brauer dans \(G/Z\), il existe des sous-groupes nilpotents \(H_i\supset Z\) de \(G\), et des caractères \(\chi_i\) de \(H_i/Z\), tels que
\[
1_{G/Z}=\sum_i n_i\,\Ind_{H_i/Z}^{G/Z}(\chi_i).
\]
L'hypothèse de récurrence s'applique à \(G/Z\), de sorte que
\[
S=(G,1)-\sum_i n_i(H_i,\chi_i)\in N(G).
\]
Soit \(R\in\Ker(b)\). On a, par (1.9.6),
\[
R=R\cdot S+\sum_i n_i\,\Ind_{H_i}^G\bigl(\Res_{H_i}^G(R)\cdot\chi_i\bigr).
\]
Puisque \(S\in N(G)\), on a \(R\cdot S\in N(G)\). On conclut en notant que, d'après 1.13 ou l'hypothèse de récurrence, \(\Res_{H_i}^G(R)\in N(H_i)\).

(B2) \(Z=\{e\}\). Soit \(C\) un sous-groupe commutatif distingué non trivial minimal et, pour \(\mu\in C^*\), soit \(G_\mu\) son fixateur. Procédons comme dans la démonstration de 1.13. On trouve que toute relation \(R\in\Ker(b)\) est combinaison linéaire de relations:
\begin{enumerate}[label=(\alph*)]
\item induites par des relations 1.11 pour \(HC,H,C,\chi\), où \(H\subset G\), \(\chi\) est un caractère de \(H\), et \(H\not\supset C\);
\item induites par des relations exprimant une identité
\[
\sum_i n_i\,\Ind_{H_i}^{G_\mu}(\chi_i)=0,
\]
avec \(H_i\supset C\) et \(\chi_i|_C=\mu\), où \(\mu\) est un caractère de \(C\).\auditmark{29}
\end{enumerate}
Puisque \(Z=\{e\}\), on a \(|G_\mu/\Ker(\mu)|<|G|\), et l'hypothèse de récurrence s'applique aux dernières. Considérons une relation (a). Si \(HC\ne G\), on applique l'hypothèse de récurrence. Si \(HC=G\), \(H\cap C\) est distingué, car distingué dans \(H\) et dans \(C\). Par minimalité de \(C\), \(H\cap C=\{e\}\): la relation est du type III.

\paragraph{Remarque 1.15.} Il est vraisemblablement possible d'extraire de [9] la description d'un système générateur de
\[
\Ker(R^+_0(G)\longrightarrow R^0(G))
\]
pour des groupes convenables \(G\). Je n'ai pas eu le courage de m'y essayer.

\section*{2. Groupes de Weil}

\paragraph{2.1. Racines de l'unité.} Soit \(\overline K\) un corps séparablement clos d'exposant caractéristique \(p\). Pour tout entier \(n\), on note \(\Z/n(1)(\overline K)\), ou simplement \(\Z/n(1)\), le groupe des racines \(n\)-ièmes de l'unité de \(\overline K\). Pour \(n\mid m\), on dispose d'une application de transition
\[
\Z/m(1)\longrightarrow\Z/n(1),\qquad x\longmapsto x^{m/n};
\]
on pose
\[
\widehat\Z(1)=\varprojlim_n \Z/n(1).
\]
Pour \(\ell\) un nombre premier, on pose de même
\[
\Z_\ell(1)=\varprojlim_r\Z/\ell^r(1).
\]
Les applications évidentes \(\widehat\Z(1)\to\Z_\ell(1)\) induisent un isomorphisme
\[
\widehat\Z(1)\iso\prod_{\ell\ne p}\Z_\ell(1).
\]
Pour tout \(\widehat\Z\)-module \(V\), on pose \(V(1)=V\otimes_{\widehat\Z}\widehat\Z(1)\). Pour \(V\) un \(\Z_\ell\)-module, on a
\[
V(1)\simeq V\otimes_{\Z_\ell}\Z_\ell(1).
\]
Pour le \(\widehat\Z\)-module \(\Q/\Z\), \(\Q/\Z(1)\) est canoniquement isomorphe au groupe des racines de l'unité de \(\overline K\) : pour \(n/m\in\Q\) et \(x\in\widehat\Z(1)\), d'image \(\bar x\) dans \(\Z/m(1)\), l'image de \(\frac{n}{m}\otimes x\) dans \(\Q/\Z(1)\) est identifiée à \(\bar x^{n}\).\auditmark{30}

Pour \(\ell\ne p\), \(\Z_\ell(1)\) est libre de rang \(1\) sur \(\Z_\ell\). Pour \(i\in\Z\), on note \(\Z_\ell(i)\) sa \(i\)-ème puissance tensorielle; pour \(i\le 0\), \(\Z_\ell(i)\) est le dual de \(\Z_\ell(-i)\). Pour tout \(\Z_\ell\)-module \(V\), on pose encore \(V(i)=V\otimes_{\Z_\ell}\Z_\ell(i)\).

\paragraph{2.2. Notations. 2.2.1.} Soit \(K\) un corps local. On note \(\|\ \|\) la valeur absolue normalisée de \(K\). Si \(dx\) est une mesure de Haar,
\[
\int f(ax)\,dx=\|a\|^{-1}\int f(x)\,dx,
\]
c'est-à-dire \(d(ax)=\|a\|\,dx\).




Pour \(s\in\C\), on note \(\omega_s\) le quasi-caractère
\[
x\longmapsto \|x\|^s
\]
de \(K^*\) dans \(\C^*\).

Si \(K\) est non archimédien, c'est-à-dire \(K\ne\R,\C\), on note \(\mathcal O\) l'anneau de la valuation \(v\) de \(K\), \(\pi\) une uniformisante, \(k\) le corps résiduel \(\mathcal O/(\pi)\), \(p\) la caractéristique de \(k\), et on pose
\[
d=[k:\F_p],\qquad q=p^d=\#k.
\]
On a donc
\[
\|x\|=q^{-v(x)}.
\]

\paragraph{2.2.2.} Supposons \(K\) non archimédien. On désigne par \(\overline K\) une clôture séparable de \(K\), et on note \(\overline{\mathcal O}\) la clôture intégrale de \(\mathcal O\) dans \(\overline K\), \(\overline k\) le corps résiduel de \(\overline{\mathcal O}\), donc une clôture algébrique de \(k\), et \(v\) la valuation de \(\overline K\), à valeurs dans \(\Q\), qui prolonge celle de \(K\).

On dispose d'une filtration en trois crans du groupe de Galois
\[
\Gal(\overline K/K)\supset I\supset P
\]
de quotients successifs
\[
\begin{array}{rcl}
\Gal(\overline K/K)/I&\iso&\Gal(\overline k/k)\iso\widehat\Z,\\[0.35em]
I/P&\iso&\widehat\Z(1)(\overline k),\\[0.35em]
P&:&\text{un pro-\(p\)-groupe.}
\end{array}
\]
On l'obtient comme suit; pour les démonstrations, on renvoie à [10].

\begin{enumerate}[label=\alph*)]
\item Par transport de structure, \(\Gal(\overline K/K)\) agit sur \(\overline{\mathcal O}\) et \(\overline k\). Le groupe d'inertie \(I\) est le noyau de la flèche de réduction
\[
v':\Gal(\overline K/K)\longrightarrow \Gal(\overline k/k).
\]
Le quotient \(\Gal(\overline k/k)\) est canoniquement isomorphe à \(\widehat\Z\), de générateur topologique la substitution de Frobenius
\[
\varphi:x\longmapsto x^q.
\]

\item Le groupe d'inertie sauvage \(P\) est l'unique pro-\(p\)-groupe de Sylow de \(I\). Le morphisme équivariant
\[
t:I\longrightarrow \widehat\Z(1)(\overline k),
\]
de noyau \(P\), est caractérisé par la congruence, modulo l'idéal maximal de \(\overline{\mathcal O}\),
\[
\frac{\sigma(x)}{x}\equiv t(\sigma)\,v(x)
\qquad(\sigma\in I,\ x\in\overline K).
\]
Ici, \(t(\sigma)\,v(x)\in\Q/\Z(1)(\overline k)\simeq\overline k^*\).
Pour \(\ell\ne p\) premier, on note \(t_\ell\) l'application composée
\[
I\longrightarrow\widehat\Z(1)(\overline k)\longrightarrow\Z_\ell(1)(\overline k).
\]
\end{enumerate}

\paragraph{2.2.3.} Soit \(k\) un corps fini à \(q\) éléments, de clôture algébrique \(\overline k\). On note \(W(\overline k/k)\) le sous-groupe de \(\Gal(\overline k/k)\simeq\widehat\Z\) engendré par
\[
\varphi:x\longmapsto x^q.
\]
On a canoniquement
\[
W(\overline k/k)\simeq\Z
\]
engendré par \(\varphi\). Le Frobenius géométrique \(F\in W(\overline k/k)\) est par définition \(\varphi^{-1}\).

\paragraph{2.2.4.} Soient \(K,\overline K\) comme en 2.2.2. Le groupe de Weil \(W(\overline K/K)\), ou groupe de Weil absolu de \(K\), est le sous-groupe de \(\Gal(\overline K/K)\) formé des \(\sigma\) tels que \(v'(\sigma)\) soit une puissance entière du Frobenius \(\varphi\). On le munit de la topologie induisant la topologie naturelle de \(I\), et pour laquelle \(I\) est ouvert :
\[
\begin{tikzcd}[column sep=1.55em,row sep=2.1em]
0 \arrow[r] & I \arrow[r] \arrow[d,equal] &
W(\overline K/K) \arrow[r] \arrow[d,hook] &
W(\overline k/k)\simeq\Z \arrow[r] \arrow[d,hook] & 0\\
0 \arrow[r] & I \arrow[r] &
\Gal(\overline K/K) \arrow[r] &
\Gal(\overline k/k)\simeq\widehat\Z \arrow[r] & 0 .
\end{tikzcd}
\]
Le groupe \(\Gal(\overline K/K)\) est le complété de \(W(\overline K/K)\) pour la topologie des sous-groupes ouverts d'indice fini; pour \(L\) une extension finie de \(K\) dans \(\overline K\), le sous-groupe correspondant de \(W(\overline K/K)\) s'identifie à \(W(\overline K/L)\). On appellera parfois Frobenius géométriques les éléments de \(\Gal(\overline K/K)\) ou \(W(\overline K/K)\) d'image \(F\).

\paragraph{2.2.5.} Supposons \(K\) archimédien, de clôture algébrique \(\overline K\). On définit \(W(\overline K/K)\) comme l'extension suivante de \(\Gal(\overline K/K)\) par \(\overline K^*\):
\begin{enumerate}[label=\arabic*)]
\item \(K\simeq\R\): \(W(\overline K/K)\) est engendré par \(\overline K^*\) et un élément \(F\), avec
\[
F^2=-1,\qquad FzF^{-1}=\overline z\quad(z\in\overline K^*).
\]
\par
\item \(K\simeq\C\): \(W(\overline K/K)=\overline K^*\).
\end{enumerate}

\paragraph{2.3.} La théorie du corps de classe local fournit un isomorphisme entre \(W(\overline K/K)^{\ab}\) et \(K^*\).

Supposons \(K\) non archimédien. Pour une raison expliquée en 3.6, nous normaliserons alors cet isomorphisme, c'est-à-dire choisirons lui ou son inverse, de telle sorte que les Frobenius géométriques correspondent aux uniformisantes:
\[
\begin{tikzcd}[column sep=2.7em,row sep=1.9em]
P \arrow[r,hook] \arrow[d] & I \arrow[r,hook] \arrow[d] &
W(\overline K/K) \arrow[r] \arrow[d] &
W(\overline k/k)\simeq\Z \arrow[d,"{-1}"]\\
1+(\pi) \arrow[r,hook] & \mathcal O^* \arrow[r,hook] &
K^* \arrow[r,"v"] & \Z .
\end{tikzcd}
\]
Une représentation de \(W(\overline K/K)\) est non ramifiée, respectivement modérément ramifiée, si elle est triviale sur \(I\), respectivement sur \(P\). De même, un quasi-caractère \(\chi\) de \(K^*\) est non ramifié, respectivement modérément ramifié, s'il est trivial sur \(\mathcal O^*\), respectivement sur \(1+(\pi)\). Pour \(\ell\) premier à \(p\), l'action de \(\Gal(\overline K/K)\) sur \(\Z_\ell(1)\) est non ramifiée, décrite par le quasi-caractère
\[
\omega_1=\|x\|
\]
de \(K^*\).

Pour \(K\) complexe, l'isomorphisme de la théorie du corps de classe local est celui évident sur 2.2.5; pour \(K\) réel, c'est celui qui rend vrais 2.3.1 et 2.3.2 ci-dessous.

Soit \(L\subset\overline K\) une extension finie de \(K\), de sorte que \(W(\overline K/L)\subset W(\overline K/K)\). Les diagrammes suivants, où \(N\) est la norme et \(t\) le transfert, sont commutatifs ([10], XIII, §4):
\[
\tag{2.3.1}
\begin{tikzcd}[column sep=3.3em,row sep=2.2em]
W(\overline K/L) \arrow[r] \arrow[d] &
W(\overline K/L)^{\ab} \arrow[r,"\sim"] \arrow[d] &
L^* \arrow[d,"N_{L/K}"]\\
W(\overline K/K) \arrow[r] &
W(\overline K/K)^{\ab} \arrow[r,"\sim"] &
K^* ,
\end{tikzcd}
\]
\[
\tag{2.3.2}
\begin{tikzcd}[column sep=3.5em,row sep=2.3em]
W(\overline K/K)^{\ab} \arrow[r,"\sim"] \arrow[d,"t"'] &
K^* \arrow[d]\\
W(\overline K/L)^{\ab} \arrow[r,"\sim"] &
L^* .
\end{tikzcd}
\]

\paragraph{2.4.} Soient \(K\) un corps de fonctions d'une variable sur \(\F_p\) et \(k\) son corps des constantes, c'est-à-dire la fermeture algébrique de \(\F_p\) dans \(K\). Soient \(\overline K\) une clôture séparable de \(K\), \(\overline k\) la clôture algébrique de \(k\) dans \(\overline K\), et \(\Gal^0(\overline K/K)\) le noyau de l'application de restriction de \(\Gal(\overline K/K)\) dans \(\Gal(\overline k/k)\). Le groupe de Weil global absolu \(W(\overline K/K)\) est, en analogie avec 2.2.4, défini par le diagramme
\[
\begin{tikzcd}[column sep=1.55em,row sep=2.1em]
0 \arrow[r] & \Gal^0(\overline K/K) \arrow[r] \arrow[d,equal] &
W(\overline K/K) \arrow[r] \arrow[d,hook] &
W(\overline k/k)=\Z \arrow[r] \arrow[d,hook] & 0\\
0 \arrow[r] & \Gal^0(\overline K/K) \arrow[r] &
\Gal(\overline K/K) \arrow[r] &
\Gal(\overline k/k)=\widehat\Z \arrow[r] & 0 .
\end{tikzcd}
\]
Soient \(v\) une place de \(K\), et supposons choisie une place de \(\overline K\) au-dessus de \(v\). On dispose alors d'un diagramme commutatif
\[
\begin{tikzcd}[column sep=1.55em,row sep=2.1em]
0 \arrow[r] & I_v \arrow[r] \arrow[d,hook] &
W(\overline K_v/K_v) \arrow[r] \arrow[d,hook] &
W(\overline k_v/k_v)=\Z \arrow[r] \arrow[d,"{[k_v:k]}"] & 0\\
0 \arrow[r] & \Gal^0(\overline K/K) \arrow[r] &
W(\overline K/K) \arrow[r] &
W(\overline k/k)=\Z \arrow[r] & 0 .
\end{tikzcd}
\]
La théorie du corps de classe global fournit un isomorphisme
\[
W(\overline K/K)^{\ab}\iso\A_K^*/K^*
\]
où \(\A_K^*/K^*\) est le groupe des classes d'idèles, rendant commutatifs les diagrammes
\[
\begin{tikzcd}[column sep=3.0em,row sep=2.0em]
W(\overline K_v/K_v)^{\ab} \arrow[r] \arrow[d] &
W(\overline K/K)^{\ab} \arrow[d]\\
K_v^* \arrow[r] & \A_K^*/K^* .
\end{tikzcd}
\]

\paragraph{2.5.} Pour la définition des groupes de Weil globaux dans le cas des corps de nombres, on renvoie à Weil [19] ou à [2], XIV. Nous ne nous en servirons qu'incidemment.

\section*{3. Rappels sur les fonctions \(L\)}

\paragraph{3.1.} Soit \(K\) un corps local et reprenons les notations 2.2.1. On désignera par \(dx\) une mesure de Haar sur \(K\), par \(d^*x\) une mesure de Haar sur \(K^*\), et par \(\psi\) un caractère additif non trivial de \(K\).

\paragraph{3.2.} Pour un quasi-caractère, c'est-à-dire un homomorphisme continu,
\[
\chi:K^*\longrightarrow\C^*,
\]
on définit comme suit \(L(\chi)\in\C\cup\{\infty\}\).

\[
\tag{3.2.1}
K\simeq\R.
\]
On pose
\[
\Gamma_\R(s)=\pi^{-s/2}\Gamma(s/2),
\]
et, pour \(x\) le plongement de \(K\) dans \(\C\) et \(N=0\) ou \(1\),
\[
L(x^{-N}\omega_s)=\Gamma_\R(s).
\]

\[
\tag{3.2.2}
K\simeq\C.
\]
On pose
\[
\Gamma_\C(s)=2(2\pi)^{-s}\Gamma(s),
\]
et, pour \(z\) un plongement de \(K\) dans \(\C\) et \(N\ge0\),
\[
L(z^{-N}\omega_s)=\Gamma_\C(s).
\]

\[
\tag{3.2.3}
K\text{ non archimédien.}
\]
On pose
\[
L(\chi)=1\quad(\chi\text{ ramifié}),
\qquad
L(\chi)=\frac1{1-\chi(\pi)}\quad(\chi\text{ non ramifié}).
\]

\paragraph{3.3.} \(\psi\) et \(dx\) étant donnés, on pose, pour \(f\) une fonction \(C^\infty\) à support compact sur \(K\),
\[
\widehat f(y)=\int_K f(x)\psi(xy)\,dx.
\]
On définit \(\varepsilon(\chi,\psi,dx)\in\C^*\) par l'équation fonctionnelle locale de Tate, indépendante de \(f\),
\[
\tag{3.3.1}
\frac{\displaystyle\int_{K^*}\widehat f(x)\,\omega_1\chi^{-1}(x)\,d^*x}
     {L(\omega_1\chi^{-1})}
=
\varepsilon(\chi,\psi,dx)\,
\frac{\displaystyle\int_{K^*}f(x)\chi(x)\,d^*x}
     {L(\chi)}.
\]
On prend la même mesure de Haar \(d^*x\) dans les deux membres. Les deux membres sont définis par prolongement analytique en \(\chi\) dans les familles \(\chi\omega_s\), \(s\in\C\). Dans ces familles, \(\varepsilon\) est de la forme \(Ae^{Bs}\).

On a
\[
\tag{3.3.2}
\varepsilon(\chi,\psi,a\,dx)=a\,\varepsilon(\chi,\psi,dx),
\]
\[
\tag{3.3.3}
\varepsilon(\chi,\psi(ax),dx)=\chi(a)\|a\|^{-1}\varepsilon(\chi,\psi,dx).
\]

\paragraph{3.4.} Pour \(K\) non archimédien, nous poserons:
\[
n(\psi)=\text{ le plus grand entier }n\text{ tel que }\psi|_{\pi^{-n}\mathcal O}=1;
\]
\[
a(\chi)=
\begin{cases}
0,&\chi\text{ non ramifié},\\
\text{le plus petit }m\text{ tel que }\chi|_{1+(\pi)^m}=1,&\chi\text{ ramifié};
\end{cases}
\]
\[
sw(\chi)=
\begin{cases}
0,&\chi\text{ non ou modérément ramifié},\\
a(\chi)-1,&\chi\text{ ramifié};
\end{cases}
\]
et \(\gamma\) désignera un élément de \(K^*\) de valuation
\[
v(\gamma)=n(\psi)+sw(\chi)+1.
\]
La constante locale \(\varepsilon\) est donnée par (3.3.2), (3.3.3), et les formules suivantes.

\[
\tag{3.4.1}
K\simeq\R.
\]
Soient \(x\) le plongement de \(K\) dans \(\C\) et \(N=0\) ou \(1\). Pour \(\psi=\exp(2\pi ix)\) et \(dx\) la mesure de Lebesgue,
\[
\varepsilon(x^{-N}\omega_s,\psi,dx)=i^N.
\]

\[
\tag{3.4.2}
K\simeq\C.
\]
Soient \(z\) un plongement de \(K\) dans \(\C\) et \(N\ge0\). Pour
\[
\psi=\exp(2\pi i\,\Tr_{\C/\R}(z))
\]
et
\[
dx=|dz\wedge d\overline z|=2\,da\,db\quad(z=a+bi),
\]
on a
\[
\varepsilon(z^{-N}\omega_s,\psi,dx)=i^N.
\]

\paragraph{3.4.3.} \(K\) non archimédien. Si \(\chi\) est non ramifié,
\[
\int_{\mathcal O}dx=1,\qquad n(\psi)=0,
\]
alors
\[
\tag{3.4.3.1}
\varepsilon(\chi,\psi,dx)=1.
\]
Pour \(\chi\) ramifié, on a
\[
\tag{3.4.3.2}
\varepsilon(\chi,\psi,dx)=
\int_{K^*}\chi^{-1}(x)\psi(x)\,dx
=
\sum_n\int_{v(x)=n}\chi^{-1}(x)\psi(x)\,dx
=
\int_{\gamma^{-1}\mathcal O^*}\chi^{-1}(x)\psi(x)\,dx.
\]

On verra qu'il est naturel d'unifier ces deux formules de la façon suivante. Si on définit \(\varepsilon_0(\chi,\psi,dx)\) par
\[
\varepsilon(\chi,\psi,dx)=\varepsilon_0(\chi,\psi,dx)
\quad(\chi\text{ ramifié}),
\]
et
\[
\tag{3.4.3.3}
\varepsilon(\chi,\psi,dx)=\varepsilon_0(\chi,\psi,dx)\,(-\chi(\pi))^{-1}
\quad(\chi\text{ non ramifié}),
\]
on a
\[
\tag{3.4.3.4}
\varepsilon_0(\chi,\psi,dx)=
\int_{\gamma^{-1}\mathcal O^*}\chi^{-1}(x)\psi(x)\,dx.
\]
De (3.4.3.3) et (3.4.3.4), on déduit que, pour \(\omega\) non ramifié,
\[
\tag{3.4.3.5}
\varepsilon_0(\chi\omega,\psi,dx)
=
\omega(\pi)^{n(\psi)+sw(\chi)+1}\varepsilon_0(\chi,\psi,dx),
\]
et
\[
\varepsilon(\chi\omega,\psi,dx)
=
\omega(\pi)^{n(\psi)+a(\chi)}\varepsilon(\chi,\psi,dx).
\]

\paragraph{3.5.} Supposons \(K\) non archimédien, et reprenons les notations 2.2.1. Soit \(V\) un espace vectoriel de dimension finie sur un corps \(E\) de caractéristique \(0\), ou sur \(\C\). Les représentations
\[
\rho:W(\overline K/K)\longrightarrow\GL(V)
\]
seront toujours supposées continues: un sous-groupe ouvert de \(I\) agit trivialement. Pour \(V\) l'espace d'une représentation et \(V^I\) le sous-espace des invariants sous l'inertie, on pose
\[
\tag{3.5.1}
Z(V;t)=\det(1-Ft^d,V^I)^{-1},
\]
\[
\tag{3.5.2}
L(V)=Z(V;1),\qquad L(V)\in E^*\cup\{\infty\}.
\]
Notons
\[
\omega^t:W(\overline K/K)\longrightarrow E((t))^*
\]
le caractère non ramifié tel que \(\omega^t(F)=t^d\). En un sens évident, on a \(\omega_s=\omega^{p^{-s}}\). On a
\[
\tag{3.5.3}
Z(V,t)=L(V\otimes\omega^t).
\]

\paragraph{3.6.} C'est pour que (3.2.3) et (3.5.2) soient compatibles que nous avons normalisé l'isomorphisme du corps de classe local de telle sorte que les Frobenius géométriques correspondent aux uniformisantes. Tant (3.2.3) que (3.5.2) sont très naturels: 3.2.3 s'impose du point de vue analytique (cf. 3.3.1), et 3.5.2 suit la tradition des géomètres; c'est le Frobenius géométrique qui tend à agir sur la cohomologie des variétés algébriques avec pour valeurs propres des entiers algébriques.

\paragraph{3.7.} Les représentations continues complexes irréductibles de \(W(\overline K/K)\) sont, pour \(K\) complexe, les quasi-caractères. Pour \(K\) réel, elles sont, outre les quasi-caractères de \(K^*\), les représentations induites par les quasi-caractères \(\chi\) de \(\overline K^*\) tels que \(\chi\ne\chi\circ\sigma\), pour \(\sigma\) la conjugaison complexe de \(\overline K\).

Pour toute représentation complexe \(V\), exprimée, dans le groupe de Grothendieck des représentations de \(W(\overline K/K)\), comme somme de représentations simples, on définit \(L(V)\) comme suit:
\[
\tag{3.7.1}
K\simeq\R,\quad
V=\sum_i[\chi_i]+\sum_j\Ind_{\overline K^*}^{W}(\chi_j),
\]
\[
L(V)=\prod_i L_K(\chi_i)\prod_j L_{\overline K}(\chi_j).
\]
\[
\tag{3.7.2}
K\simeq\C,\quad V=\sum_i[\chi_i],
\qquad
L(V)=\prod_iL(\chi_i).
\]

\paragraph{Proposition 3.8.} Soient \(K\) un corps local et \(\overline K\) une clôture algébrique de \(K\).
\begin{enumerate}[label=(\roman*)]
\item Pour toute suite exacte
\[
0\longrightarrow V'\longrightarrow V\longrightarrow V''\longrightarrow0
\]
de représentations complexes de \(W(\overline K/K)\), on a
\[
\tag{3.8.1}
L(V)=L(V')L(V'').
\]
\item Soient \(L\) une extension finie séparable de \(K\) dans \(\overline K\), et \(V_L\) une représentation complexe de \(W(\overline K/L)\). Soit \(V_K\) la représentation induite de \(W(\overline K/K)\). On a
\[
\tag{3.8.2}
L(V_K)=L(V_L).
\]
\end{enumerate}

L'assertion (i) est triviale. Dans le cas archimédien, l'assertion (ii) résulte de la formule de duplication
\[
\Gamma_\C(s)=\Gamma_\R(s)\Gamma_\R(s+1),
\]
correspondant à
\[
\Ind([\omega_s])=[\omega_s]+[\chi^{-1}\omega_{s+1}].
\]

Prouvons (ii) pour \(K\) non archimédien. Notant \(k\) et \(\ell\) les corps résiduels de \(K\) et \(L\), on a
\[
V_K^{I_K}=
\left(\Ind_{W(\overline K/L)}^{W(\overline K/K)}V_L\right)^{I_K}
\simeq
\Ind_{W(\overline\ell/\ell)}^{W(\overline k/k)}(V_L^{I_L})
\]
comme \(W(\overline k/k)\)-représentations, les deux membres s'identifiant à
\[
\Hom_{W(\overline K/L)}(W(\overline k/k),V_L),
\]
fonctions covariantes de \(W(\overline k/k)\) dans \(V_L\), \(W(\overline K/L)\) agissant sur \(W(\overline k/k)\) par translation. Cf. le diagramme commutatif




\[
\begin{tikzcd}[column sep=3.0em,row sep=2.0em]
W(\overline K/L)\arrow[r]\arrow[d]&W(\overline\ell/\ell)\arrow[d]\\
W(\overline K/K)\arrow[r]&W(\overline k/k).
\end{tikzcd}
\]
Ceci ramène (3.8.2) au lemme bien connu suivant.

\paragraph{Lemme 3.9.} Soient \(F\) un endomorphisme de \(V\), \(n\) un entier et \(F_n\) l'endomorphisme de
\[
V_n=V\otimes \C^n
\]
tel que, pour une base \(e_0,\ldots,e_{n-1}\) de \(\C^n\),
\[
F_n(v\otimes e_i)=v\otimes e_{i+1}\quad(0\le i<n-1),
\qquad
F_n(v\otimes e_{n-1})=F(v)\otimes e_0.
\]
On a
\[
\det(1-F_nt)=\det(1-Ft^n).
\]
Par le principe de prolongement des identités algébriques, on peut supposer \(F\) diagonal, dans une base convenable de \(V\); on se ramène aussitôt au cas où \(\dim(V)=1\), et où \(F\) est la multiplication par un nombre \(a^n\), \(a\ne0\). Dans la base \(f_i=a^{-i}(v\otimes e_i)\), on a alors \(F_n(f_i)=af_{i+1}\), pour \(i\in\Z/n\Z\). Les vecteurs propres de \(F_n\) sont les
\[
\sum_i \chi(i)f_i
\]
où \(\chi\) parcourt les caractères de \(\Z/n\Z\), et l'identité des déterminants en résulte.

\paragraph{3.10.} Soit \(K\) un corps global, c'est-à-dire un \(\mathfrak A\)-corps au sens de [15]. Pour toute place \(v\) de \(K\), on notera \(K_v\) le corps local complété de \(K\) en \(v\). Les objets définis plus haut pour tout corps local seront, pour \(K_v\), notés avec un indice \(v\). On note \(\mathbb A\) l'anneau des adèles de \(K\), \(\|\ \|\) la valeur absolue vérifiant
\[
d(ax)=\|a\|\,dx
\]
pour une mesure de Haar \(dx\) sur \(\mathbb A\), et \(\omega_s\) le quasi-caractère \(\|x\|^s\) de \(\mathbb A^*\). Notant \(x_v\) la composante de \(x\in\mathbb A\) dans \(K_v\), on a donc
\[
\|x\|=\prod_v\|x_v\|_v.
\]
On désignera par \(dx\) l'unique mesure de Haar sur \(\mathbb A\) telle que
\[
\int_{\mathbb A/K}dx=1
\]
(mesure de Tamagawa), par \(d^*x\) une mesure de Haar sur \(\mathbb A^*\), et par \(\psi\) un caractère non trivial de \(\mathbb A/K\), de composantes locales \(\psi_v\). Pour \(f\) \(C^\infty\) à support compact sur \(\mathbb A\), on pose
\[
\widehat f(y)=\int_{\mathbb A}f(x)\psi(xy)\,dx.
\]



\paragraph{3.11.} L'équation fonctionnelle globale de Tate [14] s'écrit, pour \(\chi\) un quasi-caractère de \(\mathbb A^*/K^*\),
\[
\tag{3.11.1}
\int_{\mathbb A^*}\widehat f(x)\chi'(x)\,d^*x
=
\int_{\mathbb A^*}f(x)\chi(x)\,d^*x,
\qquad
\chi'=\omega_1\chi^{-1},
\]
les intégrales étant définies par prolongement analytique dans les familles \(\chi\omega_s\). Posons
\[
\tag{3.11.2}
L(\chi)=\prod_v L(\chi_v)
\]
par prolongement analytique, et, pour \(dx=\otimes_v dx_v\),
\[
\tag{3.11.3}
\varepsilon(\chi)=\prod_v\varepsilon(\chi_v,\psi_v,dx_v).
\]
Le produit (3.11.3) est indépendant du choix de \(\psi\) et de celui de la décomposition de \(dx\), comme on le déduit de (3.3.2) et (3.3.3). Si, pour presque tout \(v\),
\[
\int_{\mathcal O_v}dx_v=1,
\]
presque tous ses facteurs sont égaux à \(1\). Les formules (3.3.1) et (3.11.1) fournissent l'équation fonctionnelle globale des fonctions \(L\) de Hecke
\[
\tag{3.11.4}
\varepsilon(\chi)L(\chi')=L(\chi).
\]



\paragraph{3.12.} Les résultats énoncés ci-dessous sont démontrés dans [1] et [19]. Soient \(K\) un corps global, \(\overline K\) une clôture algébrique de \(K\), et \(V\) une représentation complexe de \(\Gal(\overline K/K)\), toujours supposée continue: l'action de \(\Gal(\overline K/K)\) se factorise par un groupe de Galois fini \(\Gal(K'/K)\). Pour chaque place \(v\) de \(K\), soit encore \(v\) une place de \(\overline K\) au-dessus, et soit \(V_v\) la représentation déduite de \(V\) par restriction au sous-groupe
\[
W(\overline K_v/K_v)
\]
du groupe de décomposition \(\Gal(\overline K_v/K_v)\).

\begin{enumerate}[label=(\Alph*)]
\item Pour \(\chi\) un quasi-caractère de \(\mathbb A^*/K^*\), de composantes locales \(\chi_v\), le produit infini
\[
L(V,\chi)=\prod_v L(V_v\otimes\chi_v)
\]
peut être défini par prolongement analytique en \(\chi\) dans les familles \(\chi\omega_s\). Il converge pour \(\Re(s)\) grand et se prolonge en une fonction méromorphe de \(s\).

\item Notons \(V^*\) le dual de \(V\). On a
\[
L(V,\chi)=\varepsilon(V,\chi)\,L(V^*,\chi'),
\qquad \chi'=\omega_1\chi^{-1},
\]
avec \(\varepsilon(V,\chi)\in\C^*\), de la forme \(Ae^{Bs}\) pour \(\chi\) variant dans les familles \(\chi\omega_s\).

\item Soient \(L\) une extension finie de \(K\) dans \(\overline K\), \(N_{L/K}\) la norme, \(V_L\) une représentation de \(\Gal(\overline K/L)\), et
\[
V_K=\Ind_{\Gal(\overline K/L)}^{\Gal(\overline K/K)}(V_L).
\]
Pour \(\chi\) un quasi-caractère de \(\mathbb A_K^*/K^*\), on a
\[
\varepsilon(V_K,\chi)=\varepsilon(V_L,\chi\circ N_{L/K}).
\]
\end{enumerate}
Les assertions (A) et (B) se démontrent par réduction au cas \(\dim V=1\), via le théorème de Brauer. L'assertion (C) résulte de (B) et de ce que
\[
L(V_K,\chi)=L(V_L,\chi\circ N_{L/K})
\]
(conséquence de 2.10).

\begin{enumerate}[label=(\Alph*),start=4]
\item On a aussi
\[
\varepsilon(V'+V'',\chi)=\varepsilon(V',\chi)\,\varepsilon(V'',\chi),
\]
ce qui permet de définir \(\varepsilon(V,\chi)\) pour \(V\) seulement une représentation virtuelle.
\end{enumerate}



\section*{4. Existence des constantes locales}

\paragraph{Théorème 4.1.} Il existe une et une seule fonction \(\varepsilon\), vérifiant les conditions (1) à (4) ci-dessous, qui associe un nombre
\[
\varepsilon(V,\psi,dx)\in\C^*
\]
à chaque classe d'isomorphie de sextuples \((K,\overline K,\psi,dx,V,\rho)\), formés d'un corps local \(K\), d'une clôture algébrique \(\overline K\) de \(K\), d'un caractère additif non trivial \(\psi:K\to\C^*\), d'une mesure de Haar \(dx\) sur \(K\), d'un espace vectoriel \(V\) de dimension finie sur \(\C\), et d'une représentation
\[
\rho:W(\overline K/K)\longrightarrow \GL(V).
\]
\begin{enumerate}[label=(\arabic*)]
\item Pour toute suite exacte de représentations
\[
0\to V'\to V\to V''\to0,
\]
on a
\[
\varepsilon(V,\psi,dx)=\varepsilon(V',\psi,dx)\varepsilon(V'',\psi,dx).
\]
Cette condition montre que \(\varepsilon(V,\psi,dx)\) ne dépend que de la classe de \(V\) dans le groupe de Grothendieck des représentations de \(W(\overline K/K)\), et permet de donner un sens à \(\varepsilon(V,\psi,dx)\) pour \(V\) seulement une représentation virtuelle.

\item
\[
\varepsilon(V,\psi,a\,dx)=a^{\dim V}\varepsilon(V,\psi,dx).
\]
En particulier, pour \(V\) virtuel de dimension \(0\), \(\varepsilon(V,\psi,dx)\) est indépendant de \(dx\); on le note \(\varepsilon(V,\psi)\).

\item Si \(L\) est une extension séparable finie de \(K\) dans \(\overline K\), et que \(V_K\) est la représentation virtuelle de \(W(\overline K/K)\) induite par une représentation virtuelle de dimension zéro \(V_L\) de \(W(\overline K/L)\), on a
\[
\varepsilon(V_K,\psi)=\varepsilon(V_L,\psi\circ\Tr_{L/K}).
\]

\item Si \(\dim V=1\), \(\rho\) étant définie par un quasi-caractère \(\chi\) de \(K^*\), on a
\[
\varepsilon(V,\psi,dx)=\varepsilon(\chi,\psi,dx).
\]
\end{enumerate}

\paragraph{4.2.} Soit
\[
v=\sum_{i,j}n_{ij}(H_i,\chi_{ij})\in R^+_0(W(\overline K/K))
\]
(1.10). On suppose les \(H_i\) deux à deux non conjugués. Soit \(K_i\) l'extension de \(K\) définie par \(H_i\), et notons encore \(\chi_{ij}\) le quasi-caractère de \(K_i^*\) défini par \(\chi_{ij}\). Pour un choix quelconque de mesures de Haar \(dx_i\) sur les \(K_i\), et pour \(\psi\) un caractère additif de \(K\), on pose
\[
\tag{4.2.1}
\varepsilon(v,\psi)=
\prod_{i,j}\varepsilon(\chi_{ij},\psi\circ\Tr_{K_i/K},dx_i)^{n_{ij}}.
\]
Il résulte de (3.3.2) que \(\varepsilon(v,\psi)\) est indépendant du choix des \(dx_i\).

\paragraph{Lemme 4.3.} Avec la notation 1.9.8, pour \(v\in R^+_0(W(\overline K/K))\),
\[
\varepsilon(v,\psi(ax))=\det(v)(a)\,\varepsilon(v,\psi).
\]
Il suffit de le vérifier pour \(v\) de la forme \((H,\chi)-(H,1)\), \(H\) correspondant à une extension \(L\) de \(K\). Soit \(v_0\) la représentation virtuelle \([\chi]-[1]\) de \(W(\overline K/L)\). La formule (3.3.3) fournit
\[
\varepsilon(v,\psi(ax))=\det(v_0)(a)\,\varepsilon(v,\psi).
\]
On sait ([10], XII, §4) que l'inclusion de \(K^*\) dans \(L^*\) s'identifie au transfert
\[
W(\overline K/K)^{\ab}\longrightarrow W(\overline K/L)^{\ab}.
\]
Puisque \(b(v)=\Ind(v_0)\), le lemme résulte de 1.2.

\paragraph{Lemme 4.4.} Le théorème est vrai dans le cas archimédien.

Pour \(K\simeq\C\), on définit \(\varepsilon\) par les formules (1) et (4), et (2) est automatique. Pour \(K\simeq\R\), on vérifie sur 2.9 que l'application
\[
b:R^+_0(W(\overline K/K))\longrightarrow R^0(W(\overline K/K))
\]
est surjective, de noyau engendré par les \(x_s-x_t\) pour
\[
x_s=(\overline K^*,\omega_s)-[\omega_s]-[\chi^{-1}\omega_{s+1}].
\]
Pour \(V\) une représentation virtuelle et \(v\in R^+_0(W(\overline K/K))\) tel que
\[
b(v)=V-\dim(V)[1],
\]
on pose
\[
\tag{4.4.1}
\varepsilon(V,\psi,dx)=\varepsilon(v,\psi)\,\varepsilon(1,\psi,dx)^{\dim V}.
\]
Vérifions que le second membre est indépendant du choix de \(v\), c'est-à-dire que, pour \(v\) vérifiant \(b(v)=0\), on a \(\varepsilon(v,\psi)=1\). D'après 4.3, il suffit de le vérifier pour \(\psi(x)=\exp(2\pi ix)\). Pour ce choix de \(\psi\), l'assertion résulte de la détermination ci-dessus de \(\Ker(b)\), et du fait que le second membre de (3.4.1), (3.4.2), soit indépendant de \(s\). Il est clair que \(\varepsilon\), défini par (4.4.1), vérifie (1) à (4).

On se restreint maintenant au cas non archimédien.



\paragraph{4.5.} Soit \(K\) un corps local non archimédien, \(\overline K\) une clôture séparable de \(K\), et \(K_{nr}\subset\overline K\) l'extension non ramifiée maximale de \(K\). Soit \(J\) un sous-groupe distingué ouvert du groupe d'inertie
\[
I=\Gal(\overline K/K_{nr})
\]
et \(L\) l'extension correspondante de \(K_{nr}\). Pour \(\sigma\in I/J\), on pose
\[
t_{I/J}(\sigma)=\inf\bigl(v_L(\sigma(x)-x)\mid x\text{ entier dans }L\bigr),
\]
et on définit des fonctions \(a_{I/J}\) et \(sw_{I/J}\) par les formules
\[
\tag{4.5.1}
a_{I/J}(\sigma)=-t_{I/J}(\sigma)\quad(\sigma\ne e),
\]
\[
sw_{I/J}(\sigma)=1-t_{I/J}(\sigma)\quad(\sigma\ne e),
\]
\[
\sum_\sigma a_{I/J}(\sigma)=\sum_\sigma sw_{I/J}(\sigma)=0.
\]
D'après Artin (voir [10], VI, §2, th. 1), \(a_{I/J}\) est le caractère d'une représentation \(A_{I/J}\) de \(I/J\), la représentation d'Artin. D'après loc. cit., prop. 2, \(sw_{I/J}\) est de même le caractère d'une représentation \(Sw_{I/J}\), la représentation de Swan. Ces représentations ne sont définies qu'à isomorphisme près. \(A_{I/J}\) est somme de \(Sw_{I/J}\) et de la représentation d'augmentation de \(I/J\). Pour \(I\supset J\supset K\), on a (loc. cit., prop. 3)
\[
\tag{4.5.2}
(A_{I/K})^{J/K}\iso A_{I/J},
\qquad
(Sw_{I/K})^{J/K}\iso Sw_{I/J}.
\]
Si \(d\) est la valuation du discriminant \(\delta_{L/K_{nr}}\), c'est-à-dire la valuation de la différente \(\mathfrak D_{L/K}\) ([10], III, §3, prop. 6), et que \(r\) désigne une représentation régulière, on a (loc. cit., prop. 4)
\[
\tag{4.5.3}
A_{I/K}|_{J/K}=d\,r_{J/K}+A_{J/K}.
\]
Soit \(V\) une représentation de \(W(\overline K/K)\), triviale sur \(J\subset I\), et de caractère \(\psi\). On définit les conducteurs d'Artin et de Swan de \(V\) par les formules
\[
\tag{4.5.4}
a(V)=\dim\Hom_{I/J}(A_{I/J},V)=sw(V)+\dim V-\dim V^I,
\]
\[
sw(V)=\dim\Hom_{I/J}(Sw_{I/J},V)
=\frac1{|I/J|}\sum_\sigma sw_{I/J}(\sigma)\psi(\sigma).
\]
D'après 4.5.2, ces conducteurs sont indépendants du choix de \(J\). Pour \(V\) de dimension \(1\), défini par un quasi-caractère \(\chi\), on a (loc. cit., prop. 5)
\[
\tag{4.5.5}
a(V)=\text{conducteur }a(\chi)\text{ de }\chi.
\]
Les notations précédentes sont donc compatibles avec celles introduites en 3.4.

Avec la notation 1.9.8, on a:



\paragraph{Lemme 4.5.} Pour \(v\in R^+_0(W(\overline K/K))\) et \(\chi\) un quasi-caractère non ramifié de \(K^*\), on a
\[
\varepsilon(v\cdot\chi,\psi)=\chi(\pi)^{a(v)}\varepsilon(v,\psi).
\]
Il suffit de le vérifier pour \(v\) de la forme \((H,\lambda)-(H,\mu)\), \(H\) définissant une extension \(L\) de \(K\), et \(\lambda,\mu\) étant deux quasi-caractères de \(L^*\). Soit \(d\) la valuation de la différente de l'extension \(L/K\), \(f\) le degré de l'extension résiduelle, \(I\subset W(\overline K/K)\) et \(J\subset W(\overline K/L)\) les groupes d'inertie, et \(K\subset J\) assez petit. Calculons \(a(v)\). La formule d'adjonction pour les représentations induites, (4.5.3) et (4.5.5), fournissent, pour tout quasi-caractère \(\nu\) de \(H\), c'est-à-dire de \(L^*\),
\[
\tag{4.5.1'}
\begin{aligned}
a((H,\nu))
&=\dim\Hom_{I/K}(A_{I/K},\Ind_H^W(\nu))\\
&=f\,\dim\Hom_{I/K}(A_{I/K},\Ind_J^I([\nu]|_J))\\
&=f\,\dim\Hom_{J/K}(A_{J/K}+d\,r_{J/K},[\nu]|_J)\\
&=f\,a(\nu)+fd.
\end{aligned}
\]
D'où
\[
a(v)=f\bigl(a(\lambda)-a(\mu)\bigr).
\]
Pour \(\nu\) un quasi-caractère de \(L^*\), \(\chi_1\) un quasi-caractère non ramifié de \(L^*\), et \(n=n(\psi\circ\Tr_{L/K})\) comme en 3.4, on a, d'après (3.4.3.1) ou (3.4.3.2),
\[
\frac{\varepsilon(\nu\chi_1,\psi\circ\Tr,dx)}
     {\varepsilon(\nu,\psi\circ\Tr,dx)}
=
\chi_1(\pi_L)^{a(\nu)+n}.
\]
D'où
\[
\varepsilon(v\cdot\chi,\psi)\varepsilon(v,\psi)^{-1}
=
\chi\circ N_{L/K}\bigl(\pi_L^{a(\lambda)-a(\mu)}\bigr)
=
\chi(\pi)^{f(a(\lambda)-a(\mu))},
\]
ce qui est la formule annoncée.



\paragraph{Terminologie 4.6.} Soient \(K_1\) une extension galoisienne, finie ou non, de \(K\), de groupe de Galois \(G\), \(\alpha\) un quasi-caractère de \(K^*\), et \(\psi\) un caractère additif non trivial de \(K\). On dira qu'il existe une \(\alpha,\psi\)-théorie des constantes pour \(K_1/K\) s'il existe \(\varepsilon_{\alpha,\psi}\) du type suivant:
\begin{enumerate}[label=(\arabic*)]
\setcounter{enumi}{-1}
\item Pour \(H\) un sous-groupe ouvert de \(G\), définissant une extension finie \(L\) de \(K\), \(V\) une représentation de \(H\), et \(dx\) une mesure de Haar sur \(L\),
\[
\varepsilon_{\alpha,\psi}(V,dx)\in\C^*
\]
est défini.
\item Pour une suite exacte de représentations
\[
0\to V'\to V\to V''\to0,
\]
on a
\[
\varepsilon_{\alpha,\psi}(V,dx)=
\varepsilon_{\alpha,\psi}(V',dx)\varepsilon_{\alpha,\psi}(V'',dx).
\]
\(\varepsilon_{\alpha,\psi}\) garde donc un sens pour les représentations virtuelles.
\item On a
\[
\varepsilon_{\alpha,\psi}(V,a\,dx)=a^{\dim V}\varepsilon_{\alpha,\psi}(V,dx).
\]
En particulier, pour \(V\) virtuel de dimension \(0\), \(\varepsilon_{\alpha,\psi}(V,dx)\) est indépendant de \(dx\); on le note \(\varepsilon_{\alpha,\psi}(V)\).
\item Pour \(H_1\subset H_2\subset G\), définissant \(L_1\supset L_2\), et \(V\) une représentation virtuelle de dimension \(0\) de \(H_1\), on a
\[
\varepsilon_{\alpha,\psi}(\Ind_{H_1}^{H_2}V)=\varepsilon_{\alpha,\psi}(V).
\]
\item Pour \(V\) de dimension \(1\), défini par un caractère \(\chi\) de \(L^*\),
\[
\varepsilon_{\alpha,\psi}(V,dx)
=
\varepsilon(\chi\,\alpha\circ N_{L/K},\,\psi\circ\Tr_{L/K},\,dx).
\]
\end{enumerate}



\paragraph{Lemme 4.7.} Soient \(K_1/K\) et \(\alpha\) comme plus haut, et \(\psi\) un caractère additif non trivial de \(K\). Il existe au plus une \(\alpha,\psi\)-théorie des constantes pour \(K_1/K\). Pour qu'il en existe une, il faut et il suffit que, pour tout
\[
v\in\Ker(R^+_0(G)\longrightarrow R(G)),
\]
on ait
\[
\varepsilon(v\cdot\alpha,\psi)=1.
\]
Soient \(H\subset G\), \(L\) et \(V\) comme plus haut. D'après 1.5, il existe \(v\in R^+_0(H)\) d'image \(V-\dim(V)[1]\) dans \(R(H)\). On a nécessairement
\[
\tag{4.7.1}
\varepsilon_{\alpha,\psi}(V,dx)
=
\varepsilon(v\cdot(\alpha\circ N_{L/K}),\,\psi\circ\Tr_{L/K})\,
\varepsilon(\alpha\circ N_{L/K},\,\psi\circ\Tr_{L/K},\,dx)^{\dim V},
\]
d'où l'unicité. L'hypothèse assure que \(\varepsilon_{\alpha,\psi}\), défini par cette formule, est indépendant du choix de \(v\); on vérifie aisément les axiomes (0) à (4).




\paragraph{Lemme 4.8.} 
\begin{enumerate}[label=(\roman*)]
\item Soit \(\beta\) un quasi-caractère non ramifié de \(K^*\), respectivement soit \(\psi'=\psi(ax)\) un caractère additif non trivial. Pour qu'il existe une \(\alpha,\psi\)-théorie des constantes pour \(K_1/K\), il faut et il suffit qu'il en existe une \(\alpha\beta,\psi\)-théorie, respectivement une \(\alpha,\psi'\)-théorie. Pour \(H,L,V\) comme plus haut,
\[
\psi_L=\psi\circ\Tr_{L/K},\qquad
\beta_L=\beta\circ N_{L/K},\qquad
\alpha_L=\alpha\circ N_{L/K},
\]
on a alors
\[
\tag{4.8.1}
\varepsilon_{\alpha\beta,\psi}(V,dx)
=
\beta_L\!\left(\pi^{\,a(V\alpha_L)+\dim(V)n(\psi_L)}\right)
\varepsilon_{\alpha,\psi}(V,dx),
\]
\[
\tag{4.8.2}
\varepsilon_{\alpha,\psi'}(V,dx)
=
\|a\|_L^{-\dim(V)}
\det(V\alpha_L)(a)\,
\varepsilon_{\alpha,\psi}(V,dx).
\]
\end{enumerate}
La première assertion résulte du critère 4.7 et de 4.5, respectivement de 4.3. Il suffit de prouver (4.8.1) et (4.8.2) pour \(K=L\). De plus, on peut supposer que soit \(V\in R^0(\Gal(K_1/K))\), soit \(V=[1]\). Dans le premier cas, on applique 4.5 et 4.3. Dans le second, on utilise (4.7.1) et la définition 3.4.3.

On dira qu'il existe une \(\alpha\)-théorie des constantes pour \(K_1/K\) si, pour tout caractère additif non trivial \(\psi\), ou pour un tel \(\psi\), il existe une \(\alpha,\psi\)-théorie des constantes pour \(K_1/K\). Pour \(K_1\supset L\supset K\), il existe alors une \(\alpha\circ N_{L/K}\)-théorie des constantes pour \(K_1/L\).

\paragraph{Réduction 4.9.} Le théorème 4.1 est impliqué par l'assertion suivante:
\[
(*)\quad
\text{pour tout corps local non archimédien }K\text{ et toute extension galoisienne finie }K_1/K,
\]
\[
\text{il existe un quasi-caractère non ramifié }\alpha\text{ de }K^*
\text{ et une }\alpha\text{-théorie des constantes pour }K_1/K.
\]
D'après 4.8, l'assertion \((*)\) implique que, pour \(\overline K\) une clôture séparable de \(K\), et \(\alpha\) un quasi-caractère non ramifié quelconque de \(K^*\), il existe une \(\alpha\)-théorie des constantes pour \(\overline K/K\). La discussion ci-dessous nous permettra de passer du groupe de Galois au groupe de Weil.

\paragraph{4.10.} Toute représentation \(V\) de \(W(\overline K/K)\) est triviale sur un sous-groupe d'indice fini convenable \(J\) de l'inertie \(I\). Soit \(F\) un Frobenius géométrique. Puisque \(I/J\) est fini, une puissance convenable \(F^n\) de \(F\) agit trivialement sur \(I/J\) par conjugaison, donc est centrale dans \(W(\overline K/K)/J\). On dira que \(V\) a un \emph{type} si une puissance \(F^m\) de \(F^n\) n'a qu'une valeur propre \(a\). Le type de \(V\) est alors l'élément de la limite inductive
\[
\varinjlim_{n\mid m}\{X_n,\varphi_{n,m}\},\qquad
X_n=\C^*,\quad \varphi_{n,m}(x)=x^{m/n},
\]
défini par \(a\in X_m=\C^*\). Toute représentation irréductible a un type.

Les représentations de \(W(\overline K/K)\) de type donné \(\tau\) forment une catégorie abélienne \(\mathcal R_\tau\), et la catégorie de toutes les représentations de \(W(\overline K/K)\) est somme de ces sous-catégories. Si \(R_\tau(W(\overline K/K))\) est le groupe de Grothendieck de \(\mathcal R_\tau\), on a donc
\[
\tag{4.10.1}
R(W(\overline K/K))=\bigoplus_\tau R_\tau(W(\overline K/K)).
\]
Les représentations semi-simples de type \(1\) sont exactement les représentations continues de \(\Gal(\overline K/K)\). Si \(X_\tau\) est un quasi-caractère de type \(\tau\), et il existe des quasi-caractères, même non ramifiés, de tous les types, la torsion par \(X_\tau\) définit donc un isomorphisme
\[
\tag{4.10.2}
\otimes X_\tau:R(\Gal(\overline K/K))\iso R_\tau(W(\overline K/K)).
\]

Soient \(L\) une extension séparable finie de \(K\) dans \(\overline K\), et \(f\) le degré de l'extension résiduelle. Si une représentation \(V\) de \(W(\overline K/L)\) est de type \(\sigma\), la représentation induite de \(W(\overline K/K)\) est de type \(\sigma^f\). On dit que \(V\) est de \(K\)-type \(\tau\) si \(\sigma^f=\tau\). L'ensemble des représentations de \(K\)-type donné est donc stable par induction.

\paragraph{Variante 4.10.3.} Les arguments précédents s'appliquent à tout groupe extension de \(\Z\) par un groupe fini, et aux limites projectives de telles extensions. Le cas des groupes de Weil globaux nous sera utile plus tard.

\paragraph{4.11.} Admettons \((*)\). Pour tout type \(\tau\), soit \(X_\tau\) un caractère non ramifié de type \(\tau\). Puisqu'il existe une \(X_\tau\)-théorie des constantes, on peut définir \(\varepsilon\) sur \(R_\tau(W(\overline K/K))\) par la formule
\[
\varepsilon(V,\psi,dx)=
\varepsilon_{X_\tau,\psi}(V\otimes X_\tau^{-1},dx).
\]
Il résulte de l'unicité 4.7 que cet \(\varepsilon\) est indépendant du choix de \(X_\tau\). On définit ensuite \(\varepsilon\) sur \(R(W(\overline K/K))\) par additivité (4.10.1). D'après 4.7, s'il existe une théorie 4.1, elle est donnée par cet \(\varepsilon\).

On a fait ce qu'il fallait pour que \(\varepsilon\) vérifie (1), (2), (4), et (3) lorsque \(V_L\) a un type. Il reste à vérifier (3) pour \(V_L\) de la forme
\[
[\lambda\circ N_{L/K}]-[\mu\circ N_{L/K}],
\]
avec \(\lambda\) et \(\mu\) deux quasi-caractères non ramifiés de \(K^*\), de types donnés. On a, par (3.4.3.1), (3.3.2), (3.3.3),
\[
\varepsilon\bigl([\lambda\circ N_{L/K}]-[\mu\circ N_{L/K}],
\psi\circ\Tr_{L/K}\bigr)
=
(\lambda\mu^{-1})\circ N_{L/K}
\left(\pi_L^{\,n(\psi\circ\Tr_{L/K})}\right).
\]
D'après 4.8, on a d'autre part
\[
\varepsilon\bigl(\Ind([\lambda\circ N_{L/K}]-[\mu\circ N_{L/K}]),\psi\bigr)
=
(\lambda\mu^{-1})
\left(\pi_K^{\,a(\Ind(\mu\circ N_{L/K}))+[L:K]\,n(\psi)}\right).
\]
Soient \(d\) la valuation de la différente de \(L/K\), \(e\) l'indice de ramification, et \(f\) le degré de l'extension résiduelle. La formule (4.5.1) fournit
\[
a(\Ind(\mu\circ N_{L/K}))=fd.
\]
La formule à démontrer se réduit donc à
\[
f\,n(\psi\circ\Tr_{L/K})=fd+[L:K]\,n(\psi),
\]
soit
\[
n(\psi\circ\Tr_{L/K})=d+e\,n(\psi),
\]
qui est essentiellement la définition de la différente ([10], III, §3). Le point clef de la démonstration de \((*)\) est le lemme suivant.

\paragraph{Lemme 4.12.} Soit \(K'/K\) une extension galoisienne de corps globaux. Pour chaque place \(w\) de \(K\), on choisit une place de \(K'\) au-dessus, et on la note encore \(w\). Soient \(\alpha\) un quasi-caractère du groupe des classes d'idèles de \(K\), et \(v\) une place de \(K\) qui ne se décompose pas dans \(K'\). Si, pour toute place finie \(w\ne v\) de \(K\), il existe une \(\alpha_w\)-théorie des constantes pour \(K'_w/K_w\), alors il existe une \(\alpha_v\)-théorie des constantes pour \(K'_v/K_v\).

Soient \(\psi\) un caractère non trivial de \(\mathbb A_K/K\), et \(\psi_v\) sa \(v\)-composante. Nous allons construire une \(\alpha_v,\psi_v\)-théorie des constantes pour \(K'_v/K_v\).

On a
\[
\Gal(K'/K)=\Gal(K'_v/K_v).
\]
Un sous-groupe \(H\) de \(\Gal(K'_v/K_v)\) définit donc, outre une extension \(L_v\) de \(K_v\), une extension \(L\) de \(K\) de complété \(L_v\). Soit \(dx\) la mesure de Tamagawa sur \(\mathbb A_L\), \(\mathbb A_L/L\) ayant volume \(1\), et posons
\[
dx=\otimes_w dx_w,
\]
avec \(dx_v\) choisi à l'avance. Toute représentation \(V_v\) de \(H\) définit une représentation \(V\) de \(\Gal(K'/L)\), et, par restriction, des représentations \(V_w\) des groupes de décomposition \(\Gal(K'_w/K_w)\). Une constante globale \(\varepsilon(V,\alpha)\) a été définie en 3.12.

Soit \(w\) une place de \(L\), et notons encore \(w\) la place de \(K\) image. Pour \(w\) infinie, on dispose, par 4.4, de constantes locales
\[
\varepsilon\bigl(V_w,(\alpha\circ N_{L/K})_w,
(\psi\circ\Tr_{L/K})_w,dx_w\bigr).
\]
Pour \(w\) finie autre que \(v\), on dispose aussi, par hypothèse, de constantes
\[
\varepsilon_{\alpha_w,\psi_w}(V_w,dx_w).
\]
Notant les unes et les autres \(\varepsilon_\alpha(V_w,\psi_w,dx_w)\), on pose
\[
\varepsilon_{\alpha_v,\psi_v}(V_v,dx_v)
=
\frac{\varepsilon(V,\alpha)}
{\displaystyle\prod_{w\ne v}\varepsilon_\alpha(V_w,\psi_w,dx_w)}.
\]
On vérifie que \(\varepsilon_{\alpha_v,\psi_v}\) ne dépend pas du choix de la décomposition \(dx=\otimes dx_w\), et vérifie les axiomes (0) à (4). Le point clef est 3.12(C).

\paragraph{Lemme 4.13.} Soit \(K'/K\) une extension galoisienne de corps locaux. Il existe une extension galoisienne de corps globaux \(^{1}K'/{}^{1}K\), et une place \(v\) de \(^{1}K\) qui ne se décompose pas dans \(^{1}K'\), telles que \(K\) s'identifie au complété de \(^{1}K\) en \(v\), et \(K'\) à une sous-extension de \(^{1}K'_v/{}^{1}K_v\).

C'est trivial: on prend d'abord \(^{1}K'/{}^{1}K\) avec \(^{1}K'_v\simeq K'\supset K={}^{1}K_v\), puis on remplace \(^{1}K\) par l'extension définie par un groupe de décomposition en \(v\).

\paragraph{Lemme 4.14.} Soient \(K\) un corps global, \(S\) un ensemble fini de places finies de \(K\), et \(m:S\to\mathbb N\). Il existe un caractère \(\alpha\) du groupe des classes d'idèles tel que, pour \(w\in S\), le conducteur de \(\alpha_w\) soit \(\ge m(w)\), et que, pour toute place finie hors de \(S\), \(\alpha_w\) soit non ramifié.

Si \(K\) est un corps de nombres, on peut en effet librement spécifier \(\alpha\) sur
\[
\prod_{w\ \mathrm{fini}}\mathcal O_w^*,
\]
car ce sous-groupe compact de \(\mathbb A^*\) ne rencontre \(K^*\) qu'en \(\{1\}\). Si \(K\) est un corps de fonctions, de corps de constantes \(k\), on a
\[
\prod_w\mathcal O_w^*\cap K^*=k^*,
\]
et on remarque qu'il existe des caractères de \(\mathcal O_w^*\), de conducteur arbitrairement grand, qui soient triviaux sur \(k^*\).

D'après 4.12, 4.13 et 4.14, l'assertion 4.9 \((*)\), et donc 4.1, résultent des deux lemmes suivants, où \(K'/K\) est une extension galoisienne finie de corps locaux non archimédiens, de groupe de Galois \(G\), et \(\alpha\) un quasi-caractère de \(K^*\).

\paragraph{Lemme 4.15.} Si \(\alpha\) et \(K'/K\) sont non ramifiés, il existe une \(\alpha\)-théorie des constantes pour \(K'/K\).

Choisissons pour \(\psi\) un caractère de \(K\) tel que \(n(\psi)=0\). Il suffit alors de prendre
\[
\varepsilon_{\alpha,\psi}(V,dx)=
\left(\int_{\mathcal O_L} dx\right)^{\dim(V)}.
\]

\paragraph{Lemme 4.16.} Pour \(K'/K\) donné, il existe \(m\) tel que, si le conducteur de \(\alpha\) est \(\ge m\), il existe une \(\alpha\)-théorie des constantes pour \(K'/K\).

Soit \(\alpha\) de conducteur \(m\), et
\[
n=\left[\frac{m+1}{2}\right].
\]
La fonction \(\alpha(1+a)\) est alors additive en \(a\) pour \(v(a)\ge m-n\). Si \(\psi\) est un caractère additif non trivial de \(K\), il existe donc \(y\), de valuation \(-(m+n(\psi))\), et bien défini modulo \(\pi^{-(n+n(\psi))}\), tel que
\[
\alpha(1+a)=\psi(ya)\qquad\text{pour }v(a)\ge n.
\]
Pour \(\chi\) un quasi-caractère de \(K^*\), de conducteur \(\le m-n\), on a alors
\[
\tag{4.16.1}
\varepsilon(\chi\alpha,\psi,dx)
=
\sum_{u\in\mathcal O^*/(1+(\pi^n))}
\int_{\pi^{-m-n(\psi)}u}
\chi^{-1}(x)\alpha^{-1}(x)\psi(x)\,dx.
\]
Pour \(\pi^{-m-n(\psi)}u\) non contenu dans \(y(1+(\pi^{m-n}))\), l'intégrale partielle correspondante est nulle: c'est l'intégrale d'un caractère additif non trivial. Puisque \(\chi\) est constant de valeur \(\chi(y)\) sur \(y(1+(\pi^{m-n}))\), on a donc
\[
\varepsilon(\chi\alpha,\psi,dx)
=
\chi^{-1}(y)\,\varepsilon(\alpha,\psi,dx).
\]

Nous allons montrer que, pour \(m\) assez grand, on obtient une \(\alpha,\psi\)-théorie des constantes pour \(K'/K\) en posant, pour \(H\subset G\) définissant une extension \(L\) de \(K\), et \(V\) une représentation de \(H\),
\[
\varepsilon_{\alpha,\psi}(V,dx)
=
\det(V)^{-1}(y)\,
\varepsilon(\alpha\circ N_{L/K},\psi\circ\Tr_{L/K},dx)^{\dim V}.
\]
Les axiomes (0), (1), (2) sont triviaux. L'axiome (3) résulte de ce que l'inclusion de \(K^*\) dans \(L^*\) s'identifie au transfert
\[
W(\overline K/K)^{\ab}\longrightarrow W(\overline K/L)^{\ab}
\]
([10], XIII, §4), et de 1.2. Reste l'axiome (4).

Soit \(H\subset G\) correspondant à \(L/K\), et \(e\) l'indice de ramification. Rappelons que, pour \(u\in\mathcal O_L\) et \(n\) le plus petit entier tel que \(ne\ge2v(u)\),
\[
\tag{4.16.2}
N_{L/K}(1+u)=1+\Tr_{L/K}(u)\pmod{\pi^n}.
\]
On le vérifie en écrivant la norme de \(1+u\) comme le produit de ses conjugués. Nous noterons \(O(1)\) tout nombre borné en termes de \(K'/K\) seulement. D'après 4.16.2, si \(\alpha\) est de conducteur \(m\), \(\alpha\circ N_{L/K}\) est de conducteur \(em-O(1)\). De plus, pour
\[
\frac{v(a)}e\ge \left[\frac{m+1}{2}\right]+O(1),
\]
on a
\[
\alpha\circ N_{L/K}(1+a)
=
\alpha(1+\Tr(a))
=
\psi(y\,\Tr a)
=
\psi\circ\Tr_{L/K}(ya).
\]
Pour tout caractère \(\chi\) de \(H\), de conducteur \(O(1)\), on peut refaire le raisonnement fait plus haut; on obtient la formule voulue
\[
\varepsilon(\chi\,\alpha\circ N_{L/K},\psi\circ\Tr_{L/K},dx)
=
\chi^{-1}(y)\,
\varepsilon(\alpha\circ N_{L/K},\psi\circ\Tr_{L/K},dx).
\]

\section*{5. Formulaire}

\paragraph{5. Formulaire.}

\paragraph{(5.1)} Outre la constante locale \(\varepsilon(V,\psi,dx)\) définie par 4.1, nous considérerons aussi, dans le cas non archimédien, la constante \(\varepsilon_0\) définie par
\[
\varepsilon(V,\psi,dx)
=
\varepsilon_0(V,\psi,dx)\,\det(-F,V^I)^{-1}.
\]
Les constantes \(\varepsilon\) et \(\varepsilon_0\) ont les propriétés suivantes.

\paragraph{(5.2)} \(\varepsilon\) et \(\varepsilon_0\) sont additifs en \(V\), donc gardent un sens pour \(V\) virtuel.

\paragraph{(5.3)} \(\varepsilon\) et \(\varepsilon_0\) sont indépendants de \(dx\) pour \(V\) virtuel de dimension \(0\). On les note \(\varepsilon(V,\psi)\) et \(\varepsilon_0(V,\psi)\). Plus précisément,
\[
\varepsilon(V,\psi,a\,dx)=a^{\dim V}\varepsilon(V,\psi,dx),
\]
et de même pour \(\varepsilon_0\).

\paragraph{(5.4)} On a
\[
\varepsilon(V,\psi(ax),dx)
=
\det(V)(a)\,\|a\|^{-\dim V}\varepsilon(V,\psi,dx),
\]
et de même pour \(\varepsilon_0\).

\paragraph{(5.5)} Pour \(K\) non archimédien et \(\chi\) un quasi-caractère non ramifié de \(K^*\), on a
\[
\tag{5.5.1}
\varepsilon(V\chi,\psi,dx)
=
\chi\!\left(\pi^{a(V)+\dim(V)n(\psi)}\right)\varepsilon(V,\psi,dx),
\]
et
\[
\varepsilon_0(V\chi,\psi,dx)
=
\chi\!\left(\pi^{\dim(V)+sw(V)+\dim(V)n(\psi)}\right)
\varepsilon_0(V,\psi,dx).
\]
En particulier, si on pose
\[
\varepsilon(V,\psi,dx,s)=\varepsilon(V\omega_s,\psi,dx),
\]
et de même pour \(\varepsilon_0\), on a
\[
\tag{5.5.2}
\varepsilon(V,\psi,dx,s)
=
\varepsilon(V,\psi,dx)\,
q^{-(a(V)+\dim(V)n(\psi))s},
\]
et
\[
\varepsilon_0(V,\psi,dx,s)
=
\varepsilon_0(V,\psi,dx)\,
q^{-(sw(V)+\dim(V)(n(\psi)+1))s}.
\]
Pour \(W\) une représentation non ramifiée, on a de même
\[
\tag{5.5.3}
\varepsilon(V\otimes W,\psi,dx)
=
\det W\!\left(\pi^{a(V)+\dim(V)n(\psi)}\right)
\varepsilon(V,\psi,dx)^{\dim W},
\]
\[
\varepsilon_0(V\otimes W,\psi,dx)
=
\det W\!\left(\pi^{sw(V)+\dim(V)(n(\psi)+1)}\right)
\varepsilon_0(V,\psi,dx)^{\dim W}.
\]

\paragraph{(5.6)} Pour \(L/K\) une extension séparable finie, \(V_L\) une représentation virtuelle de dimension \(0\) de \(W(\overline K/L)\), et \(V_K\) la représentation induite de \(W(\overline K/K)\), on a
\[
\tag{5.6.1}
\varepsilon(V_K,\psi)=\varepsilon(V_L,\psi\circ\Tr_{L/K}),
\]
et de même pour \(\varepsilon_0\). En d'autres termes, il existe
\[
\lambda(L/K,\psi,dx_L,dx_K)\in\C^*
\]
tel que, pour \(V_L\) de dimension quelconque,
\[
\tag{5.6.2}
\varepsilon(V_K,\psi,dx_K)
=
\lambda(L/K,\psi,dx_L,dx_K)^{\dim(V_L)}
\varepsilon(V_L,\psi\circ\Tr_{L/K},dx_L).
\]
La même formule vaut pour \(\varepsilon_0\).

Dans Langlands [9], c'est (5.6.2) qui fournit le cadre de la théorie. De plus, il prend systématiquement pour \(dx\) la mesure de Haar autoduale pour \(\psi\), et il l'omet de la notation. Les constantes sont de plus normalisées pour être de valeur absolue complexe \(1\). Tout ceci lui impose une autre formule (5.4).

\paragraph{5.7.} Soient \(V^*\) le dual de \(V\), et \(dx'\) la mesure de Haar duale de \(dx\), relativement à \(\psi\). On a
\[
\tag{5.7.1}
\varepsilon(V,\psi,dx)\,
\varepsilon(V^*\omega_1,\psi(-x),dx')=1.
\]
En particulier, si \(K\) est non archimédien et que \(V\) est unitaire, de sorte que \(V^*=\overline V\), on a
\[
\tag{5.7.2}
|\varepsilon(V,\psi,dx)|^2
=
\left(\frac{dx'}{dx}\right)^{-\dim(V)}
q^{a(V)+\dim(V)n(\psi)}.
\]

\paragraph{5.8.} \(\varepsilon\) apparaît dans l'équation fonctionnelle locale de Tate; posant \(\chi'=\omega_1\chi^{-1}\),
\[
\tag{5.8.1}
\frac{\displaystyle\int \widehat f(x)\chi'(x)d^*x}{L(\chi')}
=
\varepsilon(\chi,\psi,dx)\,
\frac{\displaystyle\int f(x)\chi(x)d^*x}{L(\chi)}
\qquad\text{(Tate).}
\]
Pour \(K\) non archimédien, \(\gamma\in K^*\) de valuation \(1+sw(\chi)+n(\psi)\), on a (3.4.3.4)
\[
\tag{5.8.2}
\varepsilon_0(\chi,\psi,dx)
=
\int_{\gamma^{-1}\mathcal O^*}\chi^{-1}(x)\psi(x)\,dx.
\]

\paragraph{5.9.} Pour \(K\) non archimédien, si \(\chi\) est non ramifié, on a
\[
\varepsilon(\chi,\psi,dx)
=
\chi(\pi^{n(\psi)})\,q^{n(\psi)}
\int_{\mathcal O}dx.
\]
En particulier, si \(n(\psi)=0\) et que \(\int_{\mathcal O}dx=1\), on a
\[
\varepsilon(\chi,\psi,dx)=1.
\]

\paragraph{5.10.} Prenons \(K\) non archimédien, \(a(\chi)\le1\), \(n(\psi)=-1\), et
\[
\int_{(\pi)}dx=1.
\]
Soient \(\chi_k\) et \(\psi_k\) les caractères de \(k^*\) et \(k\) définis par \(\chi\) et \(\psi\). On a
\[
\tag{5.10.1}
\varepsilon_0(\chi,\psi,dx)=-\tau(\chi_k,\psi_k),
\]
où \(\tau\) est la somme de Gauss
\[
\tag{5.10.2}
\tau(\chi_k,\psi_k)
=
-\sum_{x\in k^*}\chi_k^{-1}(x)\psi_k(x),
\]
égale à \(1\) pour \(\chi_k\) trivial. En particulier,
\[
\tag{5.10.3}
\varepsilon_0([\chi]-[1],\psi)=\tau(\chi_k,\psi_k).
\]




\paragraph{5.11.} Soit \(K\) un corps global. Pour \(V\) une représentation du groupe de Weil global de \(K\), \(\psi\) un caractère additif non trivial de \(\mathbb A_K/K\), et
\[
dx=\bigotimes_v dx_v
\]
la mesure de Tamagawa sur \(\mathbb A_K\), on pose
\[
\tag{5.11.1}
L(V)=\prod_v L(V_v),
\]
\[
\tag{5.11.2}
\varepsilon(V)=\prod_v\varepsilon(V_v,\psi_v,dx_v).
\]
L'équation fonctionnelle globale s'écrit
\[
\tag{5.11.3}
L(V)=\varepsilon(V)\,L(V^*\otimes\omega_1).
\]

Démonstrations. Les formules 5.2, 5.3, 5.6 et 5.8 sont 4.1, respectivement (1), (2), (3), (4), avec 3.3. Les formules 5.4 et 5.5 résultent de 4.8. Il suffit de prouver (5.7.1) pour \(V\) défini par un caractère; la formule résulte alors de (5.8.1) et de la formule d'inversion de Fourier. La formule (5.7.2) résulte de (5.7.1), (5.3) et (5.4). La formule 5.9 résulte de (3.4.3.1), et 5.10 de (5.8.2). Enfin, 5.11 résulte de ce qui précède et de (3.11.4), par des arguments connus [1], [19].

La restriction de 4.1 au cas des extensions et des caractères modérément ramifiés équivaut aux identités suivantes sur les sommes de Gauss (5.10.2). On y désigne par \(k\) un corps fini à \(q\) éléments, et par \(\psi\) un caractère additif non trivial de \(k\).

\paragraph{5.12. Hasse--Davenport.} Soient \(k'\) une extension de \(k\), de degré \(n\), et \(\chi\) un caractère de \(k^*\). On a
\[
\tag{5.12.1}
\tau(\chi\circ N_{k'/k},\psi\circ\Tr_{k'/k})
=
\tau(\chi,\psi)^n.
\]

\paragraph{5.13.} Soient \(\chi\) un caractère de \(k^*\) et \(n\) un entier premier à \(q\). Soit \(X\) un ensemble de représentants des classes d'isomorphie de couples \((k',\chi')\), formés d'une extension \(k'\) de \(k\) et d'un caractère \(\chi'\) de \(k'^*\), tels que
\[
\tag{a}
\chi'(x^n)=\chi\circ N_{k'/k}(x),
\]
et
\[
\tag{b}
\chi'\text{ ne soit pas de la forme }\chi''\circ N_{k'/k''}
\]
pour \(k''\) intermédiaire entre \(k\) et \(k'\), c'est-à-dire
\[
\chi'^{\,q^{f'}}\ne\chi'
\qquad(0<f'<[k':k]).
\]
Posons
\[
\tau(\chi,\psi^n)=
\lambda_n(\chi,\psi)\,
\prod_{(k',\chi')\in X}
\tau(\chi',\psi\circ\Tr_{k'/k}).
\]
Alors
\[
\tag{5.13.1}
\lambda_n(\chi,\psi)\text{ est indépendant de }\chi.
\]
Cette identité se dévisse en les deux suivantes.

\paragraph{5.14 (Langlands [9], 7.8).} Soient \(\ell\) un nombre premier à \(q\), \(f\) l'ordre de \(q\) dans \((\mathbb Z/\ell)^*\), \(k'\) une extension de \(k\) de degré \(f\), \(T\) un ensemble de représentants des orbites de \(\Gal(k'/k)\) dans les caractères non triviaux d'ordre \(\ell\) de \(k'^*\), et \(\chi\) un caractère de \(k^*\). On a
\[
\tau(\chi^\ell,\psi^\ell)\,
\prod_{\mu\in T}\tau(\mu,\psi\circ\Tr_{k'/k})
=
\tau(\chi,\psi)\,
\prod_{\mu\in T}\tau(\chi\circ N_{k'/k}\cdot\mu,\psi\circ\Tr_{k'/k}).
\]

\paragraph{5.15 (Langlands [9], 7.9).} Soient \(\ell\) un nombre premier divisant \(q-1\), \(T\) l'ensemble des caractères d'ordre \(\ell\) de \(k^*\), \(\chi\) un caractère de \(k^*\) qui n'est pas une puissance \(\ell\)-ième, \(k'\) une extension de degré \(\ell\) de \(k\), et \(\chi'\) un caractère de \(k'^*\) tel que
\[
\chi'(x^\ell)=\chi\circ N_{k'/k}(x).
\]
On a
\[
\tau(\chi,\psi^\ell)\,
\prod_{\mu\in T}\tau(\mu,\psi)
=
\tau(\chi',\psi\circ\Tr_{k'/k}).
\]

\paragraph{5.16. Démonstrations.} Si \(K'/K\) est une extension modérément ramifiée de corps locaux non archimédiens, et si \(\psi\) est un caractère additif de \(K\) tel que
\[
n(\psi)=-1,
\]
on a encore
\[
n(\psi\circ\Tr_{K'/K})=-1,
\]
de sorte que 5.10 s'applique tant à \(K\) qu'à \(K'\).

Soit \(V\) une représentation modérément ramifiée de \(W(\overline K/K)\). On peut toujours l'écrire comme somme de représentations induites par des caractères modérés \(\chi_i\) d'extensions non ramifiées \(K_i/K\):
\[
V=\dim(V)[1]+\sum_i\Ind([\chi_i]-[1])
+\sum_i\bigl(\Ind([1])-[K_i:K][1]\bigr).
\]
Soient \(k_i\) le corps résiduel de \(K_i\), et \(\widetilde\chi_i\) le caractère de \(k_i^*\) défini par \(\chi_i\). Pour \(dx\) une mesure de Haar sur \(K\) telle que
\[
\int_{(\pi)}dx=1,
\]
on trouve
\[
\tag{5.16.1}
\varepsilon_0(V,\psi,dx)
=
(-1)^{\dim V}
\prod_i\tau(\widetilde\chi_i,\psi_k\circ\Tr_{k_i/k}).
\]

La formule (5.12) s'obtient en prenant une extension non ramifiée \(K'/K\), et en écrivant que, pour \(\chi\) un caractère, ici modérément ramifié, de \(K^*\), on a
\[
\Ind(\chi\circ N_{K'/K}-[1])
=
\sum_\mu([\chi\mu]-[\mu]),
\]
où \(\mu\) parcourt les caractères non ramifiés d'ordre divisant \([K':K]\). Cette formule (5.12) assure que le second membre de (5.16.1) ne dépend que de \(V\).

Les identités 5.13 s'obtiennent en prenant une extension totalement ramifiée de degré \(n\), \(K'/K\), et en conjuguant 5.6.1 et (5.16.1) pour \(V\) induite par un caractère modéré \(\chi\) de \(K'\). Les formules (5.14) et (5.15) sont le cas particulier de (5.13.1) où \(n\) est un nombre premier \(\ell\). Dans 5.14, respectivement 5.15, on suppose que le caractère \(\chi\) de \(K^*\) est, respectivement n'est pas, une puissance \(\ell\)-ième.

\section*{6. Réduction modulo \(\ell\) des constantes locales}

Soient \(K\) un corps local non archimédien, \(\overline K\) une clôture séparable de \(K\), et reprenons les notations 2.2. Soit \(A\) un anneau commutatif dans lequel \(p\) soit inversible. On simplifierait l'exposé, sans perte de généralité, en supposant que \(A\) est local, voire un corps. Pour \(x\in K\), \(\|x\|\) a par hypothèse un sens dans \(A\).

\paragraph{6.1.} Une mesure de Haar \(\mu\) sur \(K\), à valeurs dans un \(\mathbb Z[1/p]\)-module \(M\), est une fonction simplement additive de parties compactes ouvertes de \(K\), à valeurs dans \(M\), et invariante par translation. La fonction additive définie par
\[
\mu_0(a+x\mathcal O)=\|x\|
\]
est l'unique mesure de Haar à valeurs dans \(\mathbb Z[1/p]\) telle que
\[
\mu_0(\mathcal O)=1,
\]
et toute mesure de Haar à valeurs dans \(M\) est de la forme
\[
m\mu_0\qquad(m=\mu(\mathcal O)\in M).
\]
On peut intégrer une fonction localement constante à support compact, à valeurs dans \(A\), contre une mesure de Haar à valeurs dans \(A\).

\paragraph{6.2.} Soient
\[
\rho:W(\overline K/K)\longrightarrow\GL_A(V)
\]
une représentation du groupe de Weil sur un \(A\)-module libre, ou projectif, \(V\), et \(J\) un sous-groupe distingué ouvert de \(I\) sur lequel \(\rho\) soit trivial. Pour tout nombre premier \(\ell\ne p\), on sait que la représentation de Swan \(Sw_{I/J}\) peut se réaliser comme un \(\mathbb Z_\ell[I/J]\)-module projectif. Si \(A\) est une \(\mathbb Z_\ell\)-algèbre,
\[
\Hom_{I/J}(Sw_{I/J},V)
\]
est donc un \(A\)-module projectif. Si \(A\) est de plus local, il est libre; son rang est le conducteur de Swan de \(V\). Plus généralement, on vérifie qu'on peut toujours définir le conducteur de Swan de \(V\) comme une fonction localement constante à valeurs entières sur \(\Spec(A)\). Ce conducteur est invariant par extension de l'anneau des scalaires \(A\).

\paragraph{6.3.} Si \(\zeta\) est une racine \(p^n\)-ième de l'unité dans \(A\), l'ordre de \(\zeta\) est une fonction localement constante sur \(\Spec(A)\). En effet, deux quelconques des polynômes cyclotomiques \(\Phi_{p^i}(X)\) engendrent l'idéal trivial de \(A[X]\), puisque \(p\in A^*\). Ainsi \(A\) est produit d'anneaux \(A_i\), \(1\le i\le n\), la \(i\)-ième coordonnée de \(\zeta\) vérifiant
\[
\Phi_{p^i}(\zeta)=0;
\]
sur \(\Spec(A_i)\), \(\zeta\) est d'ordre \(p^i\).

Si \(\psi\) est un caractère additif continu de \(K\) à valeurs dans \(A^*\), il est à valeurs dans les racines \(p^n\)-ièmes de l'unité. Ceci permet de définir \(n(\psi)\) comme une fonction localement constante sur \(\Spec(A)\): en un idéal premier \(\lambda\), \(n(\psi)\) est \(+\infty\) si \(\psi\) est trivial modulo \(\lambda\), et sinon le plus grand entier \(n\) tel que
\[
\psi|_{\pi^{-n}\mathcal O}
\]
soit trivial modulo \(\lambda\), ou dans le localisé \(A_\lambda\), ce qui revient au même. Par abus de langage, on dit que \(\psi\) est non trivial si \(n(\psi)\) n'est nulle part \(+\infty\).

\paragraph{6.4.} Pour \(\psi\) un caractère additif non trivial de \(K\) à valeurs dans \(A^*\), \(\chi\in\Hom(K^*,A^*)\), et \(dx\) une mesure de Haar à valeurs dans \(A\), on peut maintenant définir
\[
\varepsilon_0(\chi,\psi,dx)\in A^*
\]
par la formule (3.4.3.4). Plus précisément, si \(\Spec(A)\) n'est pas connexe, on commence par écrire
\[
A=\prod_i A_i
\]
avec \(n(\psi)\) et \(Sw(\chi)\) constants sur \(\Spec(A_i)\). On définit alors la \(i\)-ième projection de \(\varepsilon_0\) par (3.4.3.4).

\paragraph{Théorème 6.5.} Il existe une et une seule fonction \(\varepsilon_0\) du type suivant:
\begin{enumerate}[label=(\alph*)]
\item Pour \(K\) un corps local non archimédien, \(\overline K\) une clôture séparable de \(K\), \(A\) un corps de caractéristique \(\ne p\), \(V\) une représentation de \(W(\overline K/K)\) dans un vectoriel de dimension finie sur \(A\), \(\psi\) un caractère additif non trivial à valeurs dans \(A^*\), et \(dx\) une mesure de Haar non nulle à valeurs dans \(A\), on a
\[
\varepsilon_0(V,\psi,dx)\in A^*.
\]
\item \(\varepsilon_0\) est invariant par extension du corps \(A\).
\item Soient \(A\) un anneau de valuation discrète de caractéristique résiduelle \(\ne p\), \(F\) son corps des fractions, et \(s=A/\mathfrak m\) son corps résiduel. Pour \(V\) une représentation de \(W(\overline K/K)\) dans un \(A\)-module libre de rang fini, et pour \(\psi\) et \(dx\) à valeurs dans \(A\), on a
\[
\varepsilon_0(V_F,\psi,dx)\in A^*,
\]
et sa réduction modulo \(\mathfrak m\) est
\[
\varepsilon_0(V_s,\overline\psi,\overline{dx}).
\]
\item Les conditions (1) à (3) de 4.1 sont vérifiées.
\item Pour \(\dim V=1\), \(\varepsilon_0\) est donné par 6.4.
\end{enumerate}

Par descente galoisienne, il suffit de traiter le cas où \(A\) est un corps séparablement clos. Si \(A\) est de caractéristique \(0\), l'énoncé, étant purement algébrique, résulte du cas \(A=\C\) traité en 4.1. Supposons donc \(A\) de caractéristique \(\ell\ne0\).

Soient \(G\) un quotient fini de \(\Gal(\overline K/K)\), et \(A_0\) un anneau de valuation discrète d'inégale caractéristique, de corps résiduel \(A_0/\mathfrak m=A\). On suppose \(\psi\) relevé en un caractère additif à valeurs dans \(A_0^*\), \(dx\) en une mesure de Haar à valeurs dans \(A_0\), et le corps des fractions \(E\) de \(A_0\) assez gros, au sens de 1.4, relativement à \(G\).

Pour tout caractère \(\chi\) d'un sous-groupe de \(G\) à valeurs dans \(E^*\), en fait dans \(A_0^*\), la constante \(\varepsilon_0\) correspondante est clairement dans \(A_0\). D'après 5.7.1, elle est même dans \(A_0^*\). On en déduit que, pour toute représentation \(V_E\) de \(G\) sur \(E\), on a
\[
\varepsilon_0(V_E,\psi,dx)\in A_0^*.
\]
On sait que l'homomorphisme de décomposition
\[
d:R_E(G)\longrightarrow R_A(G)
\]
est surjectif: les représentations virtuelles se relèvent. Pour \(V\in R_A(G)\), choisissons \(\widetilde V\in R_E(G)\) tel que \(d(\widetilde V)=V\), et posons
\[
\varepsilon_0(V,\overline\psi,\overline{dx})
=
\text{réduction modulo }\mathfrak m\text{ de }
\varepsilon_0(\widetilde V,\psi,dx).
\]

D'après 1.8, pour légitimer cette définition, il suffit de vérifier que:
\[
(*)\quad
\text{pour }H\subset G\text{ définissant une extension }L\text{ de }K,
\]
\[
\psi_L=\psi\circ\Tr_{L/K},\quad
\chi',\chi''\text{ caractères de }H\text{ à valeurs dans }A_0^*,
\]
si \(\chi'\) et \(\chi''\) sont congrus modulo \(\mathfrak m\), alors
\[
\begin{aligned}
\varepsilon_0\!\left(\Ind_H^G([\chi']-[\chi'']),\psi\right)
&=\varepsilon_0\bigl([\chi']-[\chi''],\psi_L\bigr)\\
&\equiv1\pmod{\mathfrak m}.
\end{aligned}
\]
Les conducteurs de Swan de \(\chi'\) et \(\chi''\) sont tous deux égaux au conducteur de Swan de \(\chi'\) modulo \(\mathfrak m\). Pour
\[
y=\pi^{1+Sw(\chi')+n(\psi_L)},
\]
les formules de 3.4.3 donnent
\[
\varepsilon_0(\chi',\psi_L,dx_L)=
\int_{y^{-1}\mathcal O_L^*}\chi'^{-1}(x)\psi_L(x)\,dx_L,
\]
et de même pour \(\chi''\). Numérateur et dénominateur sont des unités et sont clairement congrus modulo \(\mathfrak m\). On a donc
\[
\frac{\varepsilon_0(\chi',\psi_L,dx_L)}
{\varepsilon_0(\chi'',\psi_L,dx_L)}
\equiv1\pmod{\mathfrak m}.
\]
Pour \(V\) une représentation de \(\Gal(\overline K/K)\), on obtient ainsi une fonction \(\varepsilon_0\) ayant les propriétés annoncées. Ses propriétés formelles sont les mêmes que celles explicitées dans le formulaire 5.1 dans le cas \(A=\C\). Pour passer de \(\Gal(\overline K/K)\) à \(W(\overline K/K)\), on peut donc reprendre l'argument utilisé en 4.9, à l'aide de la deuxième formule (5.5.1).

\paragraph{Remarque 6.6.} Dans tout ce numéro, il a été essentiel de travailler avec \(\varepsilon_0\) plutôt qu'avec \(\varepsilon\). D'abord pour avoir à utiliser le conducteur de Swan, invariant par réduction modulo \(\ell\), plutôt que le conducteur d'Artin; ensuite pour n'avoir pas à distinguer, dans la définition de \(\varepsilon_0(\chi,\psi,dx)\), entre les cas non ramifié et modérément ramifié. La signification de \(\varepsilon_0\) apparaîtra plus clairement au numéro suivant.

\section*{7. Fonctions \(L\) modulaires}

\paragraph{7.1.} Soit \(X\) une courbe projective non singulière sur \(\F_p\). On suppose \(X\) irréductible, on note \(K\) son corps de fonctions rationnelles, et on identifie les points fermés de \(X\) et les places de \(K\). On ne suppose pas \(X\) géométriquement irréductible, c'est-à-dire \(\F_p\) algébriquement fermé dans \(K\).

\paragraph{7.2.} Soit \(A\) un anneau noethérien. On sait ce qu'est un faisceau constructible de \(A\)-modules sur \(X_{\mathrm{\acute et}}\) (SGA 4, IX, 2.3). Pour tout \(x\in X\), et pour \(k(\overline x)\) une extension séparablement close de \(k(x)\), un tel faisceau \(\mathcal G\) a une fibre
\[
\mathcal G_{\overline x}
\]
en \(x\), qui est un \(A\)-module de type fini. Ce dernier ne dépend que de la clôture séparable de \(k(x)\) dans \(k(\overline x)\), et \(\Gal(\overline{k(x)}/k(x))\) agit sur \(\mathcal G_{\overline x}\) par transport de structure.

De plus, pour \(v\) une place de \(K\), c'est-à-dire un point fermé de \(X\), et \(\overline K_v\) une clôture algébrique de \(K_v\), on dispose d'une flèche de spécialisation
\[
\spc_v:\mathcal G_{\overline{k(v)}}\longrightarrow \mathcal G_{\overline K_v}.
\]
Cette flèche est équivariante; son image est donc contenue dans \(\mathcal G_{\overline K_v}^{I_v}\). On dit que \(\mathcal G\) est non ramifié en \(v\) si \(\spc_v\) est un isomorphisme. Un faisceau constructible est presque partout non ramifié.

Une fois choisies une clôture séparable \(\overline K\) de \(K\), et une place de \(\overline K\) au-dessus de chaque place de \(K\), on peut réciproquement définir un faisceau constructible de \(A\)-modules comme suit:
\begin{enumerate}[label=(\alph*)]
\item on se donne les fibres \(\mathcal G_{\overline K}\) et \(\mathcal G_{\overline{k(v)}}\), respectivement des \(A\)-\(\Gal(\overline K/K)\)-modules ou des \(A\)-\(\Gal(\overline{k(v)}/k(v))\)-modules, de type fini sur \(A\);
\item on se donne les flèches de spécialisation, équivariantes pour l'action des groupes de décomposition,
\[
\mathcal G_{\overline{k(v)}}\longrightarrow\mathcal G_{\overline K}.
\]
\end{enumerate}
Pour que (a) et (b) définissent un faisceau, il suffit que les flèches de spécialisation soient presque toutes des isomorphismes.

\paragraph{7.3.} Il nous sera commode de modifier la définition précédente en remplaçant les groupes de Galois par des groupes de Weil. Soit \(Y\) un schéma sur \(\F_p\). Soient \(\overline{\F}_p\) une clôture algébrique de \(\F_p\), et
\[
\overline Y=Y\otimes_{\F_p}\overline{\F}_p.
\]
Un faisceau de Weil, respectivement un faisceau de Weil de \(A\)-modules constructible, \(\mathcal G\) sur \(Y\), est la donnée d'un faisceau, respectivement d'un faisceau de \(A\)-modules constructible, \(\overline{\mathcal G}\) sur \(\overline Y_{\mathrm{\acute et}}\), plus une action de Frobenius sur \(\overline{\mathcal G}\), au-dessus de son action sur \(\overline Y\).

\emph{Remarque inutile.} Les faisceaux de Weil en ensembles sur \(Y\) forment un topos \(Y_{\mathrm{Weil}}\). On peut le décrire comme un \(2\)-produit fibré de topos: identifiant \(\Spec(\F_p)_{\mathrm{\acute et}}\) au topos classifiant \(B\widehat{\mathbb Z}\), on a
\[
Y_{\mathrm{Weil}}\simeq
Y_{\mathrm{\acute et}}\times_{B\widehat{\mathbb Z}}B\mathbb Z.
\]
Nous appellerons simplement faisceaux les faisceaux de Weil sur \(Y\); les faisceaux sur \(Y_{\mathrm{\acute et}}\) seront appelés faisceaux étales.

\paragraph{7.4.} Un faisceau 7.3 de \(A\)-modules constructible peut se décrire en termes galoisiens comme en 7.2, en remplaçant partout
\[
\Gal(\overline K/K),\qquad
\Gal(\overline K_v/K_v),\qquad
\Gal(\overline{k_v}/k_v)
\]
par
\[
W(\overline K/K),\qquad
W(\overline K_v/K_v),\qquad
W(\overline{k_v}/k_v)=\mathbb Z.
\]

\paragraph{7.5.} Soit \(K_v\) un corps local non archimédien. On définit un faisceau de Weil sur \(\Spec(\mathcal O_v)\) en paraphrasant 7.3. Soit \(\overline K_v\) une clôture séparable de \(K_v\). Il revient au même de se donner un faisceau de Weil de \(A\)-modules constructible \(\mathcal G\) sur \(\Spec(\mathcal O_v)\), ou de se donner:
\begin{enumerate}[label=\alph*)]
\item la fibre \(\mathcal G_{\overline K_v}\), un \(A\)-\(W(\overline K_v/K_v)\)-module, de type fini sur \(A\);
\item la fibre \(\mathcal G_{\overline k_v}\), un \(A\)-\(W(\overline k_v/k_v)\)-module, de type fini sur \(A\);
\item la flèche de spécialisation \(W(\overline K_v/K_v)\)-équivariante
\[
\spc_v:\mathcal G_{\overline k_v}\longrightarrow\mathcal G_{\overline K_v}.
\]
\end{enumerate}
Soit \(v\) une place de \(K\), et posons
\[
X_v=\Spec(\mathcal O_v).
\]
Par restriction, tout faisceau de Weil de \(A\)-modules constructible sur \(X\) en définit un sur \(X_v\).

Un faisceau est plat si ses fibres sont des \(A\)-modules plats.

\paragraph{7.6.} Soit \(K_v\) un corps local non archimédien. On pose
\[
\deg(v)=[k_v:\F_p],
\]
et on note \(\omega_v^t\) le quasi-caractère non ramifié de \(W(\overline K_v/K_v)\), à valeurs dans \(A(t)\), tel que
\[
\omega_v^t(F_v)=t^{\deg(v)}.
\]
Soit \(\mathcal G\) un faisceau constructible plat sur \(\Spec(\mathcal O_v)\). On pose
\[
\tag{7.6.1}
Z(\mathcal G_v,t)
=
\det(1-F_vt^{\deg(v)},\mathcal G_{\overline k_v})^{-1}
=
\det(1-F_v,\mathcal G_{\overline k_v}\otimes\omega_v^t)^{-1}.
\]
Pour \(p\) inversible dans \(A\), on définit un conducteur
\[
\tag{7.6.2}
a(\mathcal G)=sw(\mathcal G_{\overline K_v})+
\dim(\mathcal G_{\overline K_v})-\dim(\mathcal G_{\overline k_v}),
\]
un entier si \(A\) est local, une fonction localement constante sur \(\Spec(A)\) en général.

Supposons que \(A\) soit un corps de caractéristique \(\ne p\). Soient \(\psi\) un caractère additif non trivial de \(K_v\) à valeurs dans \(A^*\), et \(dx\) une mesure de Haar à valeurs dans \(A\), comme en 6.5. On pose
\[
\tag{7.6.3}
\varepsilon(\mathcal G_v,\psi,dx,t)
=
\varepsilon_0(\mathcal G_{\overline K_v}\otimes\omega_v^t,\psi,dx)\,
\det(-F_vt^{\deg(v)},\mathcal G_{\overline k_v})^{-1}.
\]
C'est un monôme. D'après (5.5.2), son degré en \(t\) est
\[
a(\mathcal G)\deg(v).
\]

\paragraph{Remarque 7.7.} Soit \(M\) une représentation de \(W(\overline K_v/K_v)\). En termes galoisiens, on définit des faisceaux \(j_!M\) et \(j_*M\) sur \(\Spec(\mathcal O_v)\) par les flèches de spécialisation
\[
0\longrightarrow M,
\qquad
M^{I_v}\longrightarrow M.
\]
Les quantités \(\varepsilon\) et \(\varepsilon_0\) des paragraphes précédents s'identifient à la quantité (7.6.3) relative respectivement à \(j_*M\) et à \(j_!M\). Au contraire de \(j_*\), le foncteur \(j_!\) commute à la réduction modulo \(\ell\). Ceci peut expliquer pourquoi, au §6, nous n'avons considéré que \(\varepsilon_0\).

\paragraph{7.8.} Soit \(\mathcal G\) un faisceau constructible de \(A\)-modules plats sur \(X\). Pour chaque place \(v\) de \(K\), \(\mathcal G\) induit un faisceau \(\mathcal G_v\) sur \(\Spec(\mathcal O_v)\). On pose
\[
\tag{7.8.1}
Z(\mathcal G,t)=\prod_v Z(\mathcal G_v,t)\in A[[t]].
\]

\paragraph{7.9.} Supposons que \(A\) soit un corps de caractéristique \(\ne p\). Il existe sur l'anneau des adèles de \(K\) une unique mesure de Haar à valeurs dans \(\mathbb Z[1/p]\), telle que
\[
\int_{\mathbb A/K}dx=1.
\]
On peut écrire
\[
dx=\bigotimes_v dx_v
\]
où \(dx_v\) est une mesure de Haar sur \(K_v\), à valeurs dans \(\mathbb Z[1/p]\), et, pour presque tout \(v\),
\[
\int_{\mathcal O_v}dx_v=1.
\]
Les \(dx_v\) définissent des mesures de Haar à valeurs dans \(A\).

S'il existe un caractère non trivial \(\psi\) de \(\mathbb A/K\), à valeurs dans \(A^*\), on pose
\[
\tag{7.9.1}
\varepsilon(\mathcal G,t)=\prod_v
\varepsilon(\mathcal G_v,\psi_v,dx_v,t).
\]
Presque tous les termes du produit sont égaux à \(1\), et le produit ne dépend pas du choix de \(\psi\) ni de la décomposition choisie de \(dx\).

Soit \(A'\) déduit de \(A\) par adjonction d'une racine primitive \(p\)-ième de \(1\). Il existe \(\psi\) à valeurs dans \(A'^*\). La constante globale correspondante, étant indépendante de \(\psi\), est invariante par \(\Gal(A'/A)\), donc dans \(A\). Ceci permet de définir \(\varepsilon\) même s'il n'existe pas de caractère \(\psi\).

\paragraph{7.10.} Notons \(i\) l'inclusion de \(\Spec(K)\) dans \(X\). Si \(V\) est une représentation de \(W(\overline K/K)\), on note \(i_*V\) le faisceau sur \(X\), de fibre en \(\overline K\) égale à \(V\), et pour lequel les applications de spécialisation sont
\[
\spc_v:V^{I_v}\hookrightarrow V.
\]

\paragraph{Théorème 7.11.} Soit \(A\) un corps de caractéristique \(\ell\ne p\). Soient \(\mathcal G\) un faisceau constructible de \(A\)-vectoriels, et \(V\) la fibre de \(\mathcal G\) en \(\overline K\).
\begin{enumerate}[label=(\roman*)]
\item \(Z(\mathcal G,t)\) est une fraction rationnelle en \(t\).
\item Pour \(t\to\infty\), \(Z(\mathcal G,t)\) est asymptotique à \(\varepsilon(\mathcal G,t)\). En d'autres termes,
\[
Z(\mathcal G,t)/\varepsilon(\mathcal G,t)
\]
est une série formelle en \(t^{-1}\), de terme constant \(1\).
\item
\[
Z(i_*V,t)=\varepsilon(i_*V,t)\,Z(i_*V^*,p^{-1}t^{-1}).\auditmark{19}
\]
\end{enumerate}

Soit \(S\) un ensemble fini de places contenant toutes les places où la représentation \(V\) est ramifiée. Soit \(j\) l'inclusion de \(X-S\) dans \(X\), et soit \(j_!V\) le faisceau sur \(X\), de fibre générale \(V\), non ramifié en dehors de \(S\), et de fibre nulle en \(v\in S\). On vérifie aussitôt que (i) et (ii) équivalent à
\[
\tag{i'}
Z(j_!V,t)\text{ est fonction rationnelle de }t,
\]
\[
\tag{ii'}
\text{pour }t\to\infty,\qquad Z(j_!V,t)\sim\varepsilon(j_!V,t).
\]

Supposons que \(\ell\ne0\). Pour toute représentation \(W\) de \(W(\overline K/K)\), non ramifiée en dehors de \(S\), et identifiée à un faisceau localement constant sur \(X-S\), posons
\[
R^0j_*W=j_*W=i_*W,
\]
et \(R^1j_*W\) le faisceau nul en dehors de \(S\), de fibre en \(v\in S\) égale à \(H^1(I_v,W)\). On pose
\[
\tag{7.11.1}
Z(Rj_*W,t)=Z(R^0j_*W,t)\,Z(R^1j_*W,t)^{-1}.
\]

Soit \(P'\) le noyau de
\[
t_\ell:I_v\longrightarrow \mathbb Z_\ell(1).
\]
Pour toute représentation \(W\) de \(I_v\) sur \(A\), on a canoniquement
\[
\tag{7.11.2}
H^i(I_v,W)=H^i(I_v/P',W^{P'}).
\]
Pour toute représentation \(W\) de \(\mathbb Z_\ell(1)\), on a canoniquement
\[
\tag{7.11.3}
H^0(\mathbb Z_\ell(1),W)=W^{\mathbb Z_\ell(1)}
\quad\text{(invariants),}
\]
\[
H^1(\mathbb Z_\ell(1),W)=W_{\mathbb Z_\ell(1)}(-1)
\quad\text{(coinvariants tordus),}
\]
\[
H^i(\mathbb Z_\ell(1),W)=0\qquad(i\ge2).
\]
Puisque la dualité échange invariants et coinvariants, on a, cf. 2.3:

\paragraph{Lemme 7.11.4.} Pour \(\ell\ne0\), \(H^0(I_v,W)\) et \(H^1(I_v,\omega_1W^*)\) sont canoniquement en dualité.

Nous n'utiliserons pas que cette dualité peut s'interpréter comme un cup-produit à valeurs dans
\[
H^1(I_v,A(1))=\Hom(I_v,A(1))\simeq A,
\]
engendré par \(t_\ell\) modulo \(\ell\).

Pour \(\ell=0\), dans le courant de cette démonstration, nous poserons
\[
H^1(I_v,W):=H^0(I_v,\omega_1W^*)^*,
\]
et définirons \(Z(Rj_*W,t)\) par (7.11.1). Cette définition est justifiée par 7.11.5 et 7.11.6 ci-dessous. Pour \(W\) une représentation de \(W(\overline K_v/K_v)\), nous poserons encore
\[
Z_v(Rj_*W,t)=
\det(1-Ft^{\deg(v)},H^0(I_v,W))^{-1}
\det(1-Ft^{\deg(v)},H^1(I_v,W)).
\]

Soient \(A\) un anneau de valuation discrète d'inégale caractéristique \(\ell\ne p\), et \(W\) un \(A\)-module libre sur lequel \(W(\overline K_v/K_v)\) agit. Soient \(\eta\) le corps des fractions de \(A\), et \(s=A/\mathfrak m\) le corps résiduel.

\paragraph{Lemme 7.11.5.}
\[
Z_v(Rj_*W_s,t)=Z_v(Rj_*W_\eta,t)\pmod{\mathfrak m}.
\]

La démonstration ci-dessous perdra son mystère au §10. Soient
\[
P'=\Ker(t_\ell:I\to\mathbb Z_\ell(1)),
\]
\(\sigma\) un générateur de \(\mathbb Z_\ell(1)\), \(F\) un Frobenius géométrique, et \(K^\bullet\) le complexe concentré en degrés \(0\) et \(1\):
\[
K^\bullet:\quad W^{P'}\xrightarrow{\ \sigma-1\ } W^{P'}.
\]
L'endomorphisme
\[
\sum_{i=0}^{q-1}\sigma^i
\]
de \(W^{P'}\) est inversible: il suffit de le voir après réduction modulo l'idéal maximal de \(A\). Après réduction, puisqu'une puissance \(\sigma^{\ell^n}\) est l'identité, \(\sigma\) est unipotent, et \(q\) est l'unique valeur propre de \(\sum\sigma^i\).

Soit \(\underline F=(F_0,F_1)\) l'endomorphisme de \(K^\bullet\) de composantes
\[
F_0=F,\qquad
F_1=F\circ\left(\sum_{i=0}^{q-1}\sigma^i\right).\auditmark{18}
\]
La quantité
\[
Z=\det(1-F_0t^{\deg(v)},W^{P'})^{-1}
\det(1-F_1t^{\deg(v)},W^{P'})
\]
est de formation compatible à l'extension des scalaires à \(\eta\) ou à \(s\). Après extension des scalaires à \(\eta\), respectivement \(s\), on a
\[
Z=\det(1-Ft^{\deg(v)},H^0(K^\bullet))^{-1}
\det(1-Ft^{\deg(v)},H^1(K^\bullet)).
\]
On conclut en notant que
\[
(H^i(K_\eta^\bullet),H^i(\underline F))\simeq(H^i(I_v,W_\eta),F),
\]
et
\[
(H^i(K_s^\bullet),H^i(\underline F))\simeq(H^i(I_v,W_s),F).
\]

\paragraph{Variante 7.11.6.} Supposons \(A\) complet, et soit \(W\) un \(A\)-module sur lequel \(W(\overline K_v/K_v)\) agit continûment pour la topologie \(\ell\)-adique, et non plus, comme auparavant, pour la topologie discrète. On pose encore
\[
H^1(I_v,W_\eta)=(W_\eta^{P'})_{\mathbb Z_\ell(1)}(-1)
\quad\text{(coinvariants tordus),}
\]
et on définit \(Z_v(Rj_*W,t)\) comme auparavant. Le \(H^1\) défini plus haut peut s'interpréter comme un \(H^1\) continu:
\[
\tag{7.11.6.1}
H^1(I_v,W_\eta)=
\varprojlim_n H^1(I_v,W/\ell^n)\otimes_A\eta.
\]
La formule de réduction 7.11.5 reste valable, avec la même démonstration.

\paragraph{Lemme 7.11.7.} La formule (iii) équivaut à
\[
\tag{7.11.8}
Z(j_!V,t)=\varepsilon(j_!V,t)\,Z(Rj_*V^*,p^{-1}t^{-1}).\auditmark{20}
\]
Le lemme 7.11.7 résulte de l'identité locale
\[
\det(1-F_vt^{\deg(v)},H^0(I_v,V))^{-1}
=
\det(-F_vt^{\deg(v)},H^0(I_v,V))^{-1}
\det(1-F_v(p^{-1}t^{-1})^{\deg(v)},H^1(I_v,V^*))^{-1}.\auditmark{19}
\]
Si on remplace \(V\) par la représentation \(V\otimes\omega^t\) de \(W(\overline K_v/K_v)\) sur \(A((t))\), cette formule se simplifie en
\[
\det(1-F_v,H^0(I_v,V))\auditmark{21}
=
\det(-F_v,H^0(I_v,V))\,
\det(1-F_v,H^1(I_v,\omega_1\otimes V^*)),
\]
qui résulte de la dualité locale 7.11.4.

Prouvons 7.11. Si \(A=\C\), le produit infini \(Z(j_*V,t)\) converge pour \(|t|\) assez petit, et
\[
L(V,s)=Z(j_*V,q^{-s}).
\]
Ainsi \(Z\) est une fonction méromorphe de \(t\); l'équation fonctionnelle montre qu'elle est à croissance modérée pour \(t\to\infty\), donc une fonction rationnelle de \(t\). L'équation fonctionnelle de la fonction \(L\) fournit (iii).

Les énoncés (i) et (iii) sont purement algébriques. Étant vrais pour \(A=\C\), ainsi que (i') et (7.11.8), ils restent vrais pour \(A\) de caractéristique \(0\). Supposons maintenant \(A\) de caractéristique \(\ell\ne0\).

Soit \(J\) un sous-groupe ouvert distingué de \(W(\overline K/K)\), contenant \(I_v\) pour \(v\notin S\).\auditmark{22} Si une représentation \(V\) de \(W(\overline K/K)/J\) sur \(A\) se relève en caractéristique \(0\), les assertions (i') et (7.11.8) pour \(V\) s'obtiennent par réduction modulo \(\ell\), compte tenu, pour (7.11.8), de 7.11.5. Les quantités
\[
Z(j_!V,t),\qquad Z(Rj_*V,t),\qquad \varepsilon(j_!V,t)
\]
dépendent additivement de \(V\), une extension donnant un produit. Pour obtenir (i) et (7.11.8) en général, il suffit donc de savoir que, après extension des scalaires, toute représentation de \(W(\overline K/K)/J\) se relève virtuellement en caractéristique \(0\). La méthode de 4.10, cf. 4.10.3, ramène cette question au cas des groupes finis, c'est-à-dire au théorème de Brauer.

L'énoncé (ii') résulte de (7.11.8) et de ce que, pour \(t=0\), \(Z=1\).

\paragraph{Proposition 7.12.} Si \(A\) est un corps de caractéristique \(p\), et \(\mathcal G\) un faisceau constructible de \(A\)-vectoriels, \(Z(\mathcal G,t)\) est encore une fraction rationnelle en \(t\).

Dans la démonstration de 7.11(i), nous n'avons en effet pas fait usage de l'hypothèse \(\ell\ne p\).

\section*{8. Représentations \(\ell\)-adiques des groupes de Weil locaux non archimédiens; structures de Hodge et groupes de Weil locaux archimédiens}

\paragraph{8.1.} Soit \(K\) un corps local non archimédien. Nous utiliserons les notations 2.1 et 2.2.

Soient \(\ell\) un nombre premier \(\ne p\), et \(E_\lambda\) une extension finie de \(\mathbb Q_\ell\); on pourrait prendre pour \(E_\lambda\), plus généralement, un corps valué extension de \(\mathbb Q_\ell\). Une représentation \(\lambda\)-adique \(V\) de \(W(\overline K/K)\) est une représentation continue pour la topologie \(\lambda\)-adique
\[
\rho:W(\overline K/K)\longrightarrow\GL(V)
\]
de \(W(\overline K/K)\) sur un \(E_\lambda\)-vectoriel de dimension finie.

Le résultat suivant est démontré dans [13].

\paragraph{Théorème 8.2 (Grothendieck).} Soit \((V,\rho)\) une représentation \(\lambda\)-adique de \(W(\overline K/K)\). Il existe
\[
N\in\End(V)(-1),
\]
nilpotent, tel que, pour \(\sigma\) dans un sous-groupe d'indice fini de \(I\), on ait
\[
\rho(\sigma)=\exp(N\,t_\ell(\sigma)).
\]
Ce \(N\in\End(V)(-1)\) est unique, donc invariant par Galois; pour \(w\in W(\overline K/K)\), on a
\[
\rho(w)N\rho(w)^{-1}=q^{v'(w)}N.
\]
Nous allons utiliser 8.2 pour décrire les représentations \(\lambda\)-adiques de \(W(\overline K/K)\) en termes indépendants de la topologie de \(E_\lambda\). Ce sera fait deux fois: 8.3 est plus intrinsèque, et 8.4, qui nous suffit, plus élémentaire.

\paragraph{8.3.1.} À isomorphisme unique près, il existe un et un seul groupe algébrique \(G\) sur \(\mathbb Q_\ell\), isomorphe à \(\mathbb G_a\), et muni d'un isomorphisme
\[
\mathbb Q_\ell(1)\iso G(\mathbb Q_\ell).
\]
On le note \(\mathbb G_a(1)\).

\paragraph{8.3.2.}\auditmark{24} Soit \(J\subset I\) un sous-groupe invariant ouvert de \(W(\overline K/K)\), et soit
\[
J'=J\cap\Ker(t_\ell).
\]
Le quotient \(W(\overline K/K)/J'\) est une extension du groupe discret \(W(\overline K/K)/J\) par \(J/J'\). En appliquant à cette extension
\[
t_\ell:J/J'\longrightarrow\mathbb Z_\ell(1)\longrightarrow\mathbb G_a(1),
\]
on obtient une extension \(W_J^\ell\) de \(W(\overline K/K)/J\) par \(\mathbb G_a(1)\), voir Serre [12], II, 1.3:
\[
\begin{tikzcd}[column sep=2.7em,row sep=2.3em]
0 \arrow[r] & J/J' \arrow[r] \arrow[d,"t_\ell"'] &
W(\overline K/K)/J' \arrow[r] \arrow[d,"\alpha"] &
W(\overline K/K)/J \arrow[r] \arrow[d,equal] & 0\\
0 \arrow[r] & \mathbb G_a(1) \arrow[r,"i"] &
W_J^\ell \arrow[r] &
W(\overline K/K)/J \arrow[r] & 0 .
\end{tikzcd}
\]

\paragraph{8.3.3.} Les \(W_J^\ell\) forment un système projectif, et 8.2 signifie que toute représentation \(\lambda\)-adique de \(W(\overline K/K)\) se factorise par une représentation algébrique d'un \(W_J^\ell\), c'est-à-dire par une représentation du schéma en groupes
\[
W^\ell(\overline K/K)=\varprojlim_J W_J^\ell.
\]

\paragraph{8.3.4.} Soient \(I_J^\ell\) l'image réciproque de \(I/J\subset W(\overline K/K)/J\), et
\[
I^\ell=\varprojlim_J I_J^\ell.
\]
Associons à \(\bar\sigma\in I/J\), image de \(\sigma\in I\), l'élément \(\alpha(\sigma)\,i(-t_\ell(\sigma))\) de \(W_J^\ell\). Cette section définit des scindages
\[
I_J^\ell=\mathbb G_a(1)\times I/J,
\qquad
I^\ell=\mathbb G_a(1)\times I.
\]
On a un diagramme\auditmark{23}
\[
\begin{tikzcd}[column sep=1.8em,row sep=1.8em]
& 0 & 0 & 0 & \\
0 \arrow[r] & I \arrow[r] \arrow[u] & W(\overline K/K) \arrow[r] \arrow[u] & \Z \arrow[r] \arrow[u] & 0\\
0 \arrow[r] & I\times\mathbb G_a(1) \arrow[r] \arrow[u] & W^\ell(\overline K/K) \arrow[r] \arrow[u] & \Z \arrow[r] \arrow[u,equal] & 0\\
0 \arrow[r] & \mathbb G_a(1) \arrow[r,equal] \arrow[u] & \mathbb G_a(1) \arrow[u] & &\\
& 0 \arrow[u] & 0 \arrow[u] & &
\end{tikzcd}
\]

\paragraph{8.3.5.} Soit \({}'W^\ell\) le produit semi-direct de \(W(\overline K/K)\) par \(\mathbb G_a(1)\), sur lequel \(W(\overline K/K)\) agit par
\[
w x w^{-1}=q^{v'(w)}x.
\]
Il existe des isomorphismes \(\beta\) entre \(W^\ell(\overline K/K)\) et \({}'W^\ell(\overline K/K)\) qui rendent commutatif
\[
\begin{tikzcd}[column sep=2.8em,row sep=2.0em]
I\times\mathbb G_a(1)\arrow[r]\arrow[d,equal]&
W^\ell(\overline K/K)\arrow[r]\arrow[d,"\beta"]&
W(\overline K/K)\arrow[d,equal]\\
I\times\mathbb G_a(1)\arrow[r]&
{}'W^\ell(\overline K/K)\arrow[r]&
W(\overline K/K),
\end{tikzcd}
\]
et on vérifie facilement que deux quelconques sont conjugués par un élément de \(\mathbb G_a(1)(\mathbb Q_\ell)\). L'indétermination est dans le choix d'un relèvement d'un Frobenius \(F\in W(\overline K/K)\), et on utilise que \(q-1\ne0\); voir 8.4.3 ci-dessous.

\paragraph{8.3.6.} Soit \({}'W\) le schéma en groupes sur \(\mathbb Q\), produit semi-direct de \(W(\overline K/K)\) par \(\mathbb G_a\), sur lequel \(W(\overline K/K)\) agit par
\[
w x w^{-1}=q^{v'(w)}x.
\]
Chaque isomorphisme
\[
\mathbb G_a\iso\mathbb G_a(1),
\qquad\text{c'est-à-dire }\mathbb Q_\ell\iso\mathbb Q_\ell(1),
\]
définit un isomorphisme
\[
{}'W_{\mathbb Q_\ell}\iso{}'W^\ell,
\]
d'où une classe de tels automorphismes modulo \(\mathbb Q_\ell^*\) agissant sur \({}'W^\ell\) via son action sur \(\mathbb G_a(1)\). On vérifie, voir 8.4.3 ci-dessous, que ce groupe d'automorphismes agit trivialement sur l'ensemble des classes d'isomorphisme de représentations de \({}'W^\ell\). Au total, on a obtenu que:

\noindent\textbf{(8.3.7).} Les classes d'isomorphie de représentations \(\lambda\)-adiques de \(W(\overline K/K)\) sont en correspondance bijective naturelle avec les classes d'isomorphie de représentations sur \(E_\lambda\) du schéma en groupes sur \(\mathbb Q\), \({}'W\).

\paragraph{Définition 8.4.1.}\auditmark{24} Soit \(E\) un corps de caractéristique \(0\). Une représentation de \({}'W(\overline K/K)\) sur \(E\) est un couple
\[
\rho=(\rho',N')
\]
consistant en:
\begin{enumerate}[label=(\alph*)]
\item une représentation de dimension finie
\[
\rho':W(\overline K/K)\longrightarrow \GL(V)
\]
de \(W(\overline K/K)\) sur \(E\);
\item un endomorphisme nilpotent \(N'\) de \(V\), tel que
\[
\tag{8.4.1.1}
\rho'(w)N'\rho'(w)^{-1}=q^{v'(w)}N'.
\]
\end{enumerate}

\paragraph{8.4.2.} Choisissons un Frobenius géométrique \(F\in W(\overline K/K)\) et un isomorphisme
\[
\tau:\mathbb Q_\ell(1)\iso\mathbb Q_\ell.
\]
On associe comme suit une représentation \((\rho',N')\) de \({}'W(\overline K/K)\) sur \(E_\lambda\) à une représentation \(\lambda\)-adique \((V,\rho)\) de \(W(\overline K/K)\):
\begin{enumerate}[label=(\alph*)]
\item Pour \(\sigma\in I\), on pose, \(N\) étant défini par 8.2,
\[
\rho'(F^n\sigma)=\rho(F^n\sigma)\exp(-t_\ell(\sigma)N).
\]
\item \(N'\in\End(V)\) correspond à \(N\), via \(\tau\).
\end{enumerate}

\paragraph{Lemme 8.4.3.} La classe d'isomorphie de \((\rho',N')\) ne dépend que de \(\rho\), non des choix de \(F\) et \(\tau\).

Changeons \(F\) en \(F\varpi\), pour obtenir \((\rho'',N')\). On a
\[
\rho''(F^n\sigma)=
\exp\left(\frac{1-q^{-n}}{q-1}\,t_\ell(\varpi)N\right)\,\rho'(F^n\sigma)
=
A\,\rho'(F^n\sigma)\,A^{-1},
\]
avec
\[
A=\exp((q-1)^{-1}t_\ell(\varpi)N).\auditmark{54}
\]
Changer \(\tau\) en \(a\tau\) change \((\rho',N')\) en \((\rho',aN')\). Nous montrerons que, pour toute représentation \((\rho',N')\) de \({}'W(\overline K/K)\) sur un corps \(E\), et tout \(a\in E^\times\),
\[
(\rho',N')\simeq(\rho',aN').
\]
Soient \(I'=\Ker(\rho')\), et \(n\) tel que \(F^n\) soit central dans \(W/I'\), par exemple le cardinal de \(\Aut(I/I')\). Pour \(\alpha\) une classe de conjugaison dans une clôture algébrique de \(E\), soit \(V^\alpha\) le plus grand sous-espace de \(V\), stable par \(F^n\), tel que les valeurs propres de \(F^n|V^\alpha\) soient dans \(\alpha\).

On vérifie que:
\begin{enumerate}[label=(\alph*)]
\item \(V^\alpha\) est stable sous \(\rho'(W(\overline K/K))\), et \(V=\bigoplus V^\alpha\);
\item \(N' V^\alpha\subset V^{q^{-n}\alpha}\).
\end{enumerate}
Si \(A\) est la transformation linéaire laissant stables les \(V^\alpha\), et telle que
\[
A|_{V^\alpha}=\lambda(\alpha),
\]
avec
\[
\lambda(q^{-n}\alpha)=a\,\lambda(\alpha),
\]
on a donc
\[
A(\rho',N')A^{-1}=(\rho',aN').
\]
De 8.4.3 on déduit aussitôt la conclusion 8.3.7.

\paragraph{8.5.} Soit \((\rho',N')\) une représentation de \({}'W(\overline K/K)\) sur \(E\). Soit \(u\) la composante unipotente de \(\rho'(F)\):
\[
\rho'(F)=\rho'(F)^{ss}\,u.
\]
Puisque \(N'\in\End(V)\) est un vecteur propre pour l'action par conjugaison de \(\rho'(F)\), \(u\) et \(N'\) commutent. Puisqu'une puissance de \(F\) est centrale dans \(\im(\rho')\), une puissance de \(u\), donc \(u\), centralise \(\im(\rho')\). Si
\[
w_1,w_2\in W(\overline K/K)-I,
\]
\(\rho'(w_1)\) et \(\rho'(w_2)\) ont une puissance commune. On en déduit que \(u^n\) est la composante unipotente de \(\rho'(F^n\sigma)\), pour \(\sigma\in I\), et, pour \(n\ne0\), aussi celle de
\[
\rho'(F^n\sigma)\exp(aN'),
\]
qui en est le conjugué par
\[
\exp((1-q^{-n})^{-1}aN').
\]
Posons
\[
\tag{8.5.1}
\rho'^{ss}(F^n\sigma)=\rho'(F^n\sigma)u^{-n}.
\]

\paragraph{Définition 8.6.}
\begin{enumerate}[label=(\roman*)]
\item \((\rho'^{ss},N')\) est la \(F\)-semi-simplifiée de \((\rho',N')\).
\item On dit que \((\rho',N')\) est \(F\)-semi-simple si \(\rho'=\rho'^{ss}\), c'est-à-dire si, pour un, respectivement pour tout,
\[
w\in W(\overline K/K)-I,
\]
et un, respectivement tout, \(a\in E\),
\[
\rho'(w)\exp(aN')
\]
est semi-simple.
\end{enumerate}
Pour une représentation \(\lambda\)-adique, on définit encore sa \(F\)-semi-simplifiée par (8.5.1), et la \(F\)-semi-simplicité comme en (ii).

\paragraph{8.7.} Le formalisme 8.3 ou 8.4 permet de comparer des représentations \(\lambda\)-adiques pour différents corps \(E_\lambda\). Pour ce faire, introduisons la terminologie classique suivante.

Soient \(F\) une extension de \(E\), et \(\rho\) une représentation de \({}'W(\overline K/K)\) sur \(F\). On dit que \(\rho\) est \emph{définie sur} \(E\) si, pour toute, ou pour une, extension algébriquement close \(\overline F\) de \(F\), \(\rho_{\overline F}\) est isomorphe à ses conjuguées sous \(\Aut(\overline F/E)\).

Soient \(F_1\) et \(F_2\) deux extensions de \(E\), et \(\rho_i\) une représentation de \({}'W(\overline K/K)\) sur \(F_i\). On dit que \(\rho_1\) et \(\rho_2\) sont compatibles si elles sont définies sur \(E\) et si, pour une extension algébriquement close commune
\[
F\supset F_1,F_2\supset E,
\]
les représentations \((\rho_1)_F\) et \((\rho_2)_F\) sont isomorphes.

\paragraph{Définition 8.8.} Soient $E$ un corps de nombres, $\lambda_1$ et $\lambda_2$ deux places finies de $E$ ne divisant pas $p$, et $\rho_1,\rho_2$ des représentations $\lambda_1$- et $\lambda_2$-adiques de $W(\overline K/K)$. On dit que $\rho_1$ et $\rho_2$ sont compatibles si les représentations correspondantes de ${}'W(\overline K/K)$ le sont au sens de 8.7.

On laisse au lecteur la vérification du plaisant exercice suivant.

\paragraph{Proposition 8.9.} Supposons $(V_1,\rho_1)$ et $(V_2,\rho_2)$ $F$-semi-simples, et notons $N_i$ l'opérateur de monodromie associé à $(V_i,\rho_i)$. Soit $W$ l'unique filtration finie croissante de $V_i$ telle que
\[
N_i W^k\subset W^{k-2}
\]
et que $N_i^j$ induise un isomorphisme
\[
\Gr_W^j(V_i)\xrightarrow{\;N_i^j\;}\Gr_W^{-j}(V_i).
\]
Pour que $\rho_1$ et $\rho_2$ soient compatibles, il faut et il suffit que les caractères des représentations de $W(\overline K/K)$ sur les $\Gr_W^j(V_1)$ et $\Gr_W^j(V_2)$ soient à valeurs dans $E$ et respectivement égaux.

\paragraph{Exemple 8.10} On fait $E=\Q$. Soit $A$ une variété abélienne sur $K$. Les représentations de
\[
W(\overline K/K)\subset \Gal(\overline K/K)
\]
sur les
\[
V_\ell(A)=T_\ell(A)\otimes\Q_\ell
\]
sont $F$-semi-simples et compatibles.

\paragraph{Variante 8.11.} Soit $G$ un groupe algébrique linéaire sur $E$, de caractéristique $0$. Une représentation de ${}'W$ dans $G$, sur $E$, est un couple $(\rho',N')$ où
\[
\rho':W(\overline K/K)\longrightarrow G(E)
\]
est trivial sur un sous-groupe ouvert, et $N'\in\Lie(G)$, la relation (8.4.1.1) étant vérifiée. Pour $E=E_\lambda$, les formules (8.4.2) associent une représentation de ${}'W$ dans $G$ à toute représentation continue, pour la topologie $\lambda$-adique,
\[
\rho:W\longrightarrow G(E_\lambda).
\]
Sa classe d'isomorphie géométrique, après extension des scalaires à $\overline E_\lambda$, ne dépend pas des choix arbitraires faits: la première partie de la démonstration 8.4.3 s'applique encore; il reste à montrer que
\[
(\rho',N')\simeq(\rho',aN').
\]
Pour ce faire, on note que le sous-groupe algébrique $H$ de $G$ qui centralise $\im(\rho')$ et envoie $N'$ sur un multiple contient un $\rho'(F^m)$, donc un élément $g$ tel que
\[
\ad(g)(N')=q^mN'.
\]
Le caractère $\chi:H\to\mathbb G_m$ donnant l'action sur $N'$ ne peut donc qu'être surjectif.

Si $G$ est défini sur un corps de nombres $E$, on définit la compatibilité de représentations $\lambda$-adiques
\[
W(\overline K/K)\longrightarrow G(E_\lambda)
\]
comme en 8.7 et 8.8. Le sorite 8.5, 8.6 se généralise tel quel.

\paragraph{8.12. Facteurs locaux.} Soit $(\rho',N')$ une représentation de ${}'W(\overline K/K)$ sur un $E$-vectoriel $V$. Par abus de notation, on pose
\[
V^I=\Ker(N')^{\rho'(I)};
\]
c'est une représentation de $W(\overline k/k)$. On définit un conducteur, un facteur $L$ local et des constantes locales par les formules
\[
\tag{8.12.1}
a(V)=sw(V,\rho')+\dim V-\dim V^I,
\]
\[
\tag{8.12.2}
Z(V,t)=\det(1-Ft^d,V^I)^{-1},
\]
\[
\tag{8.12.3}
\varepsilon_0(V,t)=\varepsilon_0(\rho',t),
\]
\[
\tag{8.12.4}
\varepsilon(V,t)=\varepsilon_0(\rho',t)\,\det(-Ft,V^I)^{-1}.
\]
Pour $(\rho',N')$ définis par une représentation $\lambda$-adique $\rho$ sur $V$, on a
\[
V^I=V^{\rho(I)},
\qquad
Z(V,t)=\det(1-Ft,V^{\rho(I)})^{-1},\auditmark{55}
\]
et
\[
\varepsilon_0(V,t)=\varepsilon_0(\text{semi-simplifiée de }V,t).
\]

\paragraph{Le cas archimédien.} Dans le cas non archimédien, ce que fournit la géométrie, c'est-à-dire la cohomologie $\ell$-adique, ce sont des représentations $\lambda$-adiques, donc des représentations de ${}'W(\overline K/K)$, non des représentations de $W(\overline K/K)$.

Dans le cas archimédien, de même, la géométrie fournit non des représentations de $W(\overline K/K)$, mais des structures de Hodge. Voici la règle à utiliser pour passer des structures de Hodge aux représentations.

\emph{(A) $K$ complexe.} La géométrie fournit des vectoriels $V$ sur $K$, munis d'une bigraduation de Hodge
\[
V=\bigoplus V^{pq}
\]
(par exemple, pour $X$ une variété algébrique sur $K$, $H^n(X,K)$ est muni d'une telle décomposition). On en déduit une représentation semi-simple $\rho$ de
\[
W(\overline K/K)=K^*
\]
sur $V$ en posant
\[
\rho(z)v^{pq}=z^{-p}\overline z^{-q}v^{pq}\qquad(v^{pq}\in V^{pq}).
\]

\emph{(B) $K$ réel.} La géométrie fournit des vectoriels $V$ sur $\overline K$, munis d'une bigraduation de Hodge
\[
V=\bigoplus V^{pq}
\]
et d'une involution $F_\infty$ telle que
\[
F_\infty(V^{pq})=V^{qp}.
\]
Par exemple, pour $X$ une variété algébrique sur $K$, $\Gal(\overline K/K)=\Z/2\Z$ agit sur $X(\overline K)$ et $H^n(X(\overline K),\overline K)$ est muni d'une telle structure. On en déduit une représentation semi-simple $\rho$ de $W(\overline K/K)$ sur $V$ en posant
\[
\rho(z)v^{pq}=z^{-p}\overline z^{-q}v^{pq}\qquad(z\in\overline K^*\subset W(\overline K/K),\ v^{pq}\in V^{pq}),
\]
et
\[
\rho(F)v^{pq}=i^{p+q}F_\infty(v^{pq}).
\]
Ces formules diffèrent de celles de [3], 1.12, défigurée en outre par une faute d'impression p. 21, l. 2, parce qu'aux places finies j'utilisais dans loc. cit. l'inverse de l'isomorphisme de la théorie du corps de classe local utilisé ici.

Aux places finies,
\[
H_2(\mathbb P^1,\Z_\ell)\simeq\Z_\ell(1)
\]
correspond au quasi-caractère $\omega_1$. De même,
\[
H_2(\mathbb P^1(\overline K),\overline K),
\]
structure de Hodge de type $(-1,-1)$, avec $F_\infty=-1$ si $K=\R$, correspond au quasi-caractère $\omega_1$ de $W(\overline K/K)$.

\section*{9. Systèmes compatibles de représentations $\ell$-adiques}

Pour simplifier l'exposé, nous ne considérerons dans ce paragraphe que des représentations de groupes de Galois, plutôt que des représentations de groupes de Weil. Comme en 7.1, on note $X$ une courbe projective et lisse sur $\F_p$, irréductible sur $\F_p$, de corps de fonctions $K$.

\paragraph{9.1.} Soit $E_\lambda$ un corps local, extension finie de $\Q_\ell$, avec $\ell\ne p$. Soit $W$ une représentation $\lambda$-adique de $\Gal(\overline K/K)$, c'est-à-dire une représentation continue, pour la topologie de $E_\lambda$, et presque partout non ramifiée de $\Gal(\overline K/K)$ dans un vectoriel de dimension finie sur $E_\lambda$. Le groupe de Galois étant compact, il existe dans $W$ un réseau invariant $W^0$, un réseau étant un sous-$\mathcal O_\lambda$-module libre de $W$ tel que $E_\lambda W^0=W$. Les facteurs locaux
\[
Z_v(W,t)=\det(1-F_vt^{\deg(v)},W_v^{I_v})^{-1}
\]
sont donc dans $\mathcal O_\lambda[[t]]$, ainsi que leur produit $Z(W,t)$. Soit $\Gal^0(\overline K/K)$ le noyau de l'application de $\Gal(\overline K/K)$ dans $\Gal(\overline{\F}_p/\F_p)$. D'après un théorème de Grothendieck, $Z(W,t)$ est une fraction rationnelle; c'est même un polynôme si
\[
W^{\Gal^0(\overline K/K)}=W_{\Gal^0(\overline K/K)}=0,
\]
et il vérifie une équation fonctionnelle
\[
\tag{9.1.1}
Z(W,t)=\varepsilon_{\mathrm{Gr}}(W,t)Z(W^*\omega_1,t^{-1}),
\]
où $\varepsilon_{\mathrm{Gr}}$ est un monôme de degré
\[
\sum_v\deg(v)\,a(i_*W_v)
\]
(8.12.1). Le lecteur trouvera au §10 un résumé de la théorie de Grothendieck.

\paragraph{9.2.} Soient $E$ un corps de nombres, $\LLL$ un ensemble infini de places finies ne divisant pas $p$ de $E$, et $\mathcal O\subset E$ l'anneau des éléments de $E$ $\lambda$-entiers pour tout $\lambda\in\LLL$. On désigne encore par $\lambda$ l'idéal premier de $\mathcal O$ défini par $\lambda$. Supposons donnée, pour chaque $\lambda\in\LLL$, une représentation $\lambda$-adique $V_\lambda$ de $\Gal(\overline K/K)$. On suppose que, pour toute place $v$ de $K$, les représentations $\lambda$-adiques $V_{\lambda v}$ de $W(\overline K_v/K_v)$ soient compatibles, au sens strict (8.8). Il nous suffirait en fait de savoir que leurs $F$-semi-simplifiées (8.6) sont compatibles, cf. aussi 9.8. La fraction rationnelle
\[
Z(V_\lambda,t)\in\mathcal O(t)
\]
et la constante, c'est-à-dire le monôme,
\[
\varepsilon_\lambda(V,t)=\prod_v\varepsilon(V_{\lambda v},\psi_v,dx_v,t)\in\mathcal O(t)
\]
sont donc indépendantes de $\lambda$; comme d'habitude, les $\psi_v$ sont les composantes d'un caractère non trivial de $\mathbb A/K$, et $\otimes dx_v$ est une mesure de Tamagawa. On note ces fonction et constante
\[
Z(V,t),\qquad \varepsilon(V,t).
\]

\paragraph{Théorème 9.3.} Sous les hypothèses précédentes,
\[
\tag{9.3.1}
Z(V,t)=\varepsilon(V,t)Z(V^*,p^{-1}t^{-1}).\auditmark{19}
\]
Puisque $\LLL$ est infini, l'application
\[
\mathcal O\longrightarrow\prod_{\lambda\in\LLL}\mathcal O/\lambda
\]
est injective: il suffit de prouver 9.3.1 après réduction modulo $\lambda$, pour tout $\lambda\in\LLL$. Nous allons pour cela mettre 9.3.1 sous une forme compatible à la réduction modulo $\lambda$, et justifiable, modulo $\lambda$, par (7.11.8).

Soient $S$ un ensemble fini de places de $K$, contenant toutes les places où $V$ se ramifie, et $j$ l'inclusion de $X-S$ dans $X$. Pour toute représentation $\lambda$-adique $W$, non ramifiée en dehors de $S$, nous poserons
\[
Z(j,W,t)=\prod_{v\notin S}\det(1-F_vt^{\deg(v)},W_v)^{-1},
\]
\[
Z(Rj_*W,t)=
\prod_v\det(1-F_vt^{\deg(v)},W_v^{I_v})^{-1}
\Big/
\prod_{v\in S}\det(1-F_vt^{\deg(v)},W_{v,I_v}(-1))^{-1},
\]
\[
\varepsilon(j,W,t)=
\prod_{v\notin S}\varepsilon(W_v,\psi_v,dx_v,t)
\prod_{v\in S}\varepsilon_0(W_v,
\psi_v,dx_v,t).
\]
Un calcul facile, déjà fait en 7.11.7, montre que (9.3.1) équivaut à
\[
\tag{9.3.2}
Z(j,W,t)=\varepsilon(j,W,t)Z(Rj_*W^*,p^{-1}t^{-1}).\auditmark{20}
\]
On conclut par 7.11.6 et 7.11.8.

\paragraph{Stabilité 9.4.} Soient $F$ une extension de $E$ et $\LLL_F$ l'image réciproque de $\LLL$. Un système compatible de représentations $\lambda$-adiques, $\lambda\in\LLL$, définit par extension des scalaires un système compatible de représentations $\lambda$-adiques, $\lambda\in\LLL_F$. Si $\chi$ est un caractère du groupe des classes d'idèles, à valeurs dans les racines de l'unité de $F$, les $V_\lambda\chi$ forment encore un système compatible, justiciable de 9.3.

\paragraph{Corollaire 9.5.} Soit $\chi$ un caractère du groupe des classes d'idèles, à valeurs dans les racines de l'unité d'une extension finie de $E$ (cf. 9.4). Soient $S(\chi)$, respectivement $S(V)$, l'ensemble des places de $K$ où $\chi$, respectivement $V$, se ramifie. Supposons que
\[
S(V)\cap S(\chi)=\varnothing.
\]
On a
\[
\varepsilon\bigl((V-\dim V[1])([\chi]-[1]),t\bigr)
\]
égal, par définition, à
\[
\varepsilon(V\chi,t)\varepsilon(V,t)^{-1}
\varepsilon(\chi,t)^{-\dim V}\varepsilon([1],t)^{\dim V}
\]
et donc à
\[
\prod_{v\in S(V)}\chi_v(\pi_v^{a(V_v)})
\cdot
\prod_{v\in S(\chi)}\det V_v(\pi_v^{a(\chi_v)}).
\]
C'est un corollaire de 5.5.3.

\paragraph{Exemple 9.6.} Pour $E=\Q$ et $\LLL$ l'ensemble des nombres premiers autres que $p$, si $F$ est une courbe elliptique sur $K$, les $H^1(F,\Q_\ell)$ forment un système compatible de représentations $\ell$-adiques. Si l'invariant modulaire $j$ n'est pas constant, alors, pour tout caractère $\chi$ comme en 9.5, les fractions rationnelles $Z(F\chi,t)$ correspondantes sont des polynômes et vérifient 9.5.

\paragraph{9.7.} On renvoie à Weil [16] pour la relation que 9.5 implique, lorsque $\dim(V)=2$, entre $Z(V,t)$ et les «formes modulaires» sur $\GL(2,\mathbb A)/\GL(2,K)$. Le cas le plus intéressant est fourni par 9.6.

On obtient des résultats plus précis en invoquant Jacquet--Langlands. Soit par exemple $E$ une courbe elliptique sur $K$, d'invariant modulaire non constant. D'après le dictionnaire de Langlands, voir [4], §3, pour chaque place $v$ de $K$, la représentation de
\[
\Gal(\overline K_v/K_v)\supset W(\overline K_v/K_v)
\]
sur $H^1(E,\Q_\ell)$ définit une représentation admissible irréductible $\pi_v$ de $\GL_2(K_v)$. Les représentations $H^1(E,\Q_\ell)$, pour $\ell$ variable, étant compatibles, $\pi_v$ peut être défini sur $\Q$ et est indépendant de $\ell$. Nous considérerons la représentation complexe correspondante. La représentation $\pi_v$ est de dimension infinie, de poids $2$, c'est-à-dire que le centre $K_v^*$ de $\GL(2,K_v)$ agit par $\omega_2$, et a presque toujours un vecteur $\GL(2,\mathcal O_v)$-invariant. D'après 9.3, 9.6 et [8], 11.5, le produit tensoriel restreint
\[
\pi(E)=\bigotimes_v \pi_v
\]
figure comme facteur direct dans la représentation admissible
\[
L^0_{\omega_2}\bigl(\GL(2,\mathbb A)/\GL(2,K)\bigr)
\]
de $\GL(2,\mathbb A)$ dans l'espace des fonctions localement constantes sur $\GL(2,\mathbb A)/\GL(2,K)$, se transformant sous l'action du centre par le quasi-caractère $\omega_2$, à support compact modulo le centre et cuspidales.

\paragraph{Théorème 9.8.} Soient $\lambda,\mu$ des places de $E$, et $V_\lambda,V_\mu$ des représentations $\lambda$- et $\mu$-adiques de $\Gal(\overline K/K)$. Soit $S$ un ensemble fini de places de $K$, et supposons que, pour $v\notin S$, les semi-simplifiées de $V_{\lambda v}$ et $V_{\mu v}$ soient compatibles. Alors, pour toute place $v$, les semi-simplifiées de $V_{\lambda v}$ et $V_{\mu v}$ sont compatibles.

Si $S$ est pris assez grand pour contenir toutes les places ramifiées pour $V_\lambda$ ou $V_\mu$, l'hypothèse se réduit à:
\[
(*)\quad
\text{pour }v\notin S,\text{ les polynômes caractéristiques de }F_v
\text{ agissant dans }V_{\lambda v}\text{ et }V_{\mu v}
\]
\[
\text{sont dans }E[t]\text{ et égaux.}
\]
Le théorème 9.8 est donc une réponse partielle, dans le cas le moins intéressant des corps de fonctions, à la question 1 de Serre [12], I, 12.

On prendra garde aux points suivants.
\begin{enumerate}[label=(\alph*)]
\item La conclusion porte seulement sur les semi-simplifiées des représentations locales, non sur leur $F$-semi-simplifiée. En d'autres termes, avec les notations de 8.4.2, on considère seulement la semi-simplifiée de $\rho'$, et non $N$.
\item Pour $\lambda=\mu$, $V_\lambda=V_\mu$, le théorème affirme que les semi-simplifiées des représentations locales $V_{\lambda v}$ de $W(\overline K_v/K_v)$ sont définies sur $E$, c'est-à-dire ont un caractère à valeurs dans $E$.
\end{enumerate}
On peut supposer, et on suppose, que $S$ contient toutes les places où $I_v$ agit via un groupe infini. Pour $v\notin S$, il n'y a dès lors plus à distinguer entre semi-simplifiée de $V_{\lambda v}$ et $F$-semi-simplifiée.

Appliquons le théorème de Grothendieck rappelé en 9.1:
\[
Z(V_\lambda,t)=\varepsilon_{\mathrm{Gr}}(V_\lambda,t)Z(V_\lambda^*\omega_1,t^{-1}).
\]
Cette identité se récrit
\[
\prod_{v\notin S}Z_v(V_\lambda,t)\prod_{v\notin S}Z_v(V_\lambda^*\omega_1,t^{-1})^{-1}
=
\varepsilon_{\mathrm{Gr}}
\frac{\prod_{v\in S}Z_v(V_\lambda^*\omega_1,t^{-1})}
{\prod_{v\in S}Z_v(V_\lambda,t)}.\auditmark{56}\auditmark{57}
\]
Le premier membre est dans $E(T)$, et de même pour $V_\mu$. Le second membre est donc dans $E(T)$, indépendant de $\lambda$. Décomposant le second membre en le produit d'un monôme et d'une fraction rationnelle valant $1$ en $t=0$, on trouve par le calcul local de 7.11.7 que
\[
\prod_{v\in S}
\frac{\det(1-F_vt^{\deg(v)},V_{\lambda v}^{I_v})}
{\det(1-F_vt^{\deg(v)},V_{\lambda v,I_v}(-1))}
\]
est le même pour $V_\lambda$ et $V_\mu$. Les facteurs de ce produit sont additifs en $V_{\lambda v}$, une extension donnant un produit, ce qui résulte de (7.11.6.1); ce sont donc les mêmes pour $V_{\lambda v}$ et sa semi-simplifiée $V_{\lambda v}^{ss}$. Pour celle-ci, invariants de $I_v$ et coinvariants de $I_v$ sont pareils, et on trouve que
\[
\tag{9.8.1}
\prod_{v\in S}
\frac{\det(1-F_vt^{\deg(v)},V_{\lambda v}^{ss,I_v})}
{\det(1-F_v(pt)^{\deg(v)},V_{\lambda v}^{ss,I_v})}
\]
est le même pour $V_\lambda$ et $V_\mu$.

Soit $v_0\in S$, et tordons $V_\lambda$ et $V_\mu$ par un caractère $\chi$ du groupe des classes d'idèles non ramifié en $v_0$ et très ramifié en les $v\in S$ autres que $v_0$; $\chi$ est à valeurs dans une extension de $E$, et, pour l'existence, cf. 4.14. Les hypothèses faites sur $(V_\lambda,V_\mu)$ sont stables par torsion. Pour $\chi$ convenable, après torsion, les facteurs de 9.8.1 relatifs aux $v\ne v_0$ deviennent $1$. Celui relatif à $v_0$ subit une modification triviale, qui se lit sur $\chi(F_v)$. Posons
\[
P_{\lambda v}(T)=\det(1-F_vT,(V_{\lambda v}^{ss})^{I_v}).
\]
Appliquant l'invariance de (9.8.1) à $(V_\lambda\chi,V_\mu\chi)$, on trouve que, pour chaque $v\in S$, on a, dans $E(T)$,
\[
\tag{9.8.2}
\frac{P_{\lambda v}(T)}{P_{\lambda v}(qT)}=
\frac{P_{\mu v}(T)}{P_{\mu v}(qT)}.
\]
La transformation qui, à une fraction rationnelle $R$ valant $1$ en $0$, associe
\[
R_q(T)=\frac{R(T)}{R(qT)}
\]
est injective: pour les diviseurs, on a
\[
\operatorname{div}(R)=\lim_N \operatorname{div}\left(\prod_{0}^{N}R_q(q^iT)\right)
\]
(limite simple). De 9.8.2, on déduit donc que $P_{\lambda v}(T)\in E[T]$ et que
\[
\tag{9.8.3}
P_{\lambda v}(T)=P_{\mu v}(T).
\]
Les hypothèses faites sur $V_\lambda$ et $V_\mu$ sont respectées par passage de $K$ à une extension finie $K'$, en remplaçant $S$ par son image réciproque. Toute extension finie locale $K'_v/K_v$ étant induite par une extension globale, car $\Gal(\overline K_v/K_v)$ est un sous-groupe fermé de $\Gal(\overline K/K)$, on trouve que 9.8.3 vaut aussi pour les représentations de $W(\overline K_v/K'_v)$ définies par $V_\lambda$ et $V_\mu$. Prenant pour $K'_v$ les diverses extensions telles que les semi-simplifiées de ces représentations soient non ramifiées, on trouve que, pour tout $F\in W(\overline K_v/K_v)$, $F\notin I_v$, on a
\[
\det(1-Ft,V_\lambda)=\det(1-Ft,V_\mu).\auditmark{58}
\]
En dehors de $I_v$, les caractères des représentations $V_\lambda$ et $V_\mu$ coïncident donc.

Les mêmes arguments, où on ne retient de 9.8.3 que l'égalité des degrés, donnent que, pour tout sous-groupe ouvert $J$ de $I_v$ et tout caractère $\chi$ de $J$,
\[
\dim\bigl((V_{\lambda v}^{ss}\otimes\chi)^J\bigr)
=
\dim\bigl((V_{\mu v}^{ss}\otimes\chi)^J\bigr).
\]
D'après le théorème de Brauer, ceci implique que
\[
V_{\lambda v}^{ss}|_{I_v}\quad\text{et}\quad V_{\mu v}^{ss}|_{I_v}
\]
ont même caractère. Les caractères de $V_\lambda$ et $V_\mu$ coïncident donc sur tout $W(\overline K_v/K_v)$, et ceci achève la démonstration.

\paragraph{Proposition 9.9.} Soit $(V_\lambda)_{\lambda\in\LLL}$ un système infini de représentations $\lambda$-adiques. On suppose que, pour toute place $v$ de $K$ (ou, par 9.8, pour presque toute place), les semi-simplifiées des $V_{\lambda v}$ sont compatibles. Alors, pour chaque $\lambda$,
\[
\tag{9.9.1}
Z_\lambda(V,t)=\varepsilon_\lambda(V,t)Z_\lambda(V^*,p^{-1}t^{-1}).\auditmark{19}
\]
Posons
\[
Z_\lambda^{ss}(V,t)=
\prod_v\det(1-F_vt^{\deg(v)},(V_{\lambda v}^{ss})^{I_v})^{-1},
\]
\[
\varepsilon_\lambda^{ss}(V,t)=
\prod_v\varepsilon(V_{\lambda v}^{ss},\psi_v,dx_v,t).
\]
Un calcul local montre que (9.9.1) équivaut à
\[
\tag{9.9.2}
Z_\lambda^{ss}(V,t)=
\varepsilon_\lambda^{ss}(V,t)Z_\lambda^{ss}(V^*,p^{-1}t^{-1}).\auditmark{19}
\]
Le quotient des facteurs locaux en $v$ des deux membres de (9.9.1) ou (9.9.2) est additif en $V_{\lambda v}$, et ces quotients coïncident dans le cas semi-simple. L'identité à prouver (9.9.2) est indépendante de $\lambda$. On la prouve modulo $\lambda$ pour chaque $\lambda$, et on conclut comme en 9.3.




\section*{10. La théorie de Grothendieck}

\paragraph{10.1.} Soient \(X\) une courbe projective non singulière sur le corps fini \(\F_p\), \(\overline{\F}_p\) une clôture algébrique de \(\F_p\), et
\[
\overline X=X\otimes_{\F_p}\overline{\F}_p
\]
déduite de \(X\) par extension des scalaires de \(\F_p\) à \(\overline{\F}_p\). Soient \(\ell\) un nombre premier premier à \(p\), et \(\Lambda\) un anneau noethérien tué par une puissance de \(\ell\).

Soit \(G\) un faisceau de Weil, au sens de 7.3 et 7.4, constructible de \(\Lambda\)-modules plats. Notons encore \(G\) le faisceau étale sur \(\overline X\) qui s'en déduit.

\paragraph{10.2.} Les groupes de cohomologie étale
\[
H^i(\overline X,G)
\]
sont des \(\Lambda\)-modules de type fini. N'étant pas projectifs, ils ne nous seront d'aucun usage. Nous devrons utiliser un objet plus fin
\[
R\Gamma(\overline X,G)\in\operatorname{Ob}\Dparf(\Lambda)
\]
qui leur donne naissance. Rappelons que \(\Dparf(\Lambda)\) peut se définir comme la catégorie des complexes finis de \(\Lambda\)-modules projectifs de type fini, les morphismes étant les morphismes de complexes pris à homotopie près. \(R\Gamma(\overline X,G)\) est un complexe fini de \(\Lambda\)-modules projectifs de type fini, concentré en degrés \(0,1,2\) si l'on veut, bien défini à homotopie près, et dont les groupes de cohomologie sont les \(H^i(\overline X,G)\).

Par transport de structure, la substitution de Frobenius
\[
\varphi\in\Gal(\overline{\F}_p/\F_p)
\]
agit sur \(R\Gamma(\overline X,G)\): c'est un endomorphisme du complexe \(R\Gamma(\overline X,G)\), bien déterminé à homotopie près, et induisant sur la cohomologie l'automorphisme de Frobenius.

\paragraph{10.3.} Pour \(K\) un complexe fini de \(\Lambda\)-modules projectifs et \(u\) un endomorphisme de \(K\), on pose
\[
\tag{10.3.1}
\det(1-ut,K)=
\prod_i\det(1-ut,K^i)^{(-1)^i}\in\Lambda(t).
\]
Si \(u\) et \(v\) sont des endomorphismes homotopes, on a
\[
\det(1-ut,K)=\det(1-vt,K),
\]
de sorte que \(\det(1-ut,K)\) est bien défini pour \(K\in\operatorname{Ob}\Dparf(\Lambda)\) et \(u\in\End(K)\). Prenant l'opposé du coefficient de \(t\) dans le développement en série formelle de \(\det(1-ut,K)\), on définit la trace
\[
\tag{10.3.2}
\Tr(u,K)=\sum_i(-1)^i\Tr(u,K^i).
\]

Si \(u\) est une auto-équivalence d'homotopie, c'est-à-dire induit un automorphisme sur la cohomologie, on peut trouver \((K',u')\), isomorphe dans \(\Dparf(\Lambda)\) à \((K,u)\), avec \(u'\) un automorphisme du complexe \(K'\). On pose alors
\[
\tag{10.3.3}
\det(u,K)=\prod_i\det(u,K^i)^{(-1)^i},
\]
et cette quantité ne dépend pas des choix arbitraires faits. Si \(u^{-1}\) est un inverse homotopique de \(u\), on a
\[
\tag{10.3.4}
\det(1-ut,K)=\det(-ut,K)\det(1-u^{-1}t^{-1},K).
\]
Si \(D(K)\) est le dual de \(K\), complexe de composantes
\[
D(K)^i=\Hom(K^{-i},\Lambda),
\]
et si \(\check u\) est un contragrédient de \(u\), on a aussi
\[
\tag{10.3.5}
\det(1-ut,K)=\det(-ut,K)\det(1-\check u\,t^{-1},D(K)).
\]

Une esquisse de la démonstration du théorème suivant est donnée en 10.11.\auditmark{01}

\paragraph{Théorème 10.4 (Grothendieck).} On a
\[
\tag{10.4.1}
Z(G,t)=\det(1-Ft,R\Gamma(\overline X,G))^{-1}.
\]

Dans les applications, les hypothèses faites sur \(G\) sont trop restrictives. Il y a lieu de prendre comme coefficients un élément de \(\Dparf(X,\Lambda)\), c'est-à-dire un complexe fini \(G^\bullet\) de faisceaux de Weil de \(\Lambda\)-modules sur \(X\), qui, localement pour la topologie étale sur \(\overline X\), soit homotope à un complexe fini de faisceaux constructibles de \(\Lambda\)-modules plats. Utilisant 10.3, on peut, pour \(x\) un point fermé de \(X\), définir
\[
Z_x(G^\bullet,t)=
\det(1-F_x t^{\deg(x)},G_{\overline x}^\bullet)^{-1}.
\]
On pose
\[
Z(G^\bullet,t)=\prod_x Z_x(G^\bullet,t).
\]
On définit encore \(R\Gamma(\overline X,G^\bullet)\), de groupes de cohomologie les groupes d'hypercohomologie de \(\overline X\) à valeurs dans \(G^\bullet\), et on a
\[
\tag{10.4.2}
Z(G^\bullet,t)=
\det(1-Ft,R\Gamma(\overline X,G^\bullet))^{-1}.
\]

\paragraph{10.5.} Tout \(G^\bullet\in\Dparf(X,\Lambda)\) a un dual à valeur dans le complexe dualisant \(\Lambda(1)\), placé en degré \(-2\):
\[
D(G^\bullet)=R\Hom_\Lambda(G^\bullet,\Lambda(1)[2])\in\Dparf(X,\Lambda).
\]
Son \(i\)-ième faisceau de cohomologie est le \((2+i)\)-ième hyper-Ext local de \(G^\bullet\) avec \(\Lambda(1)\). La dualité de Poincaré prend la forme
\[
\tag{10.5.1}
R\Gamma(\overline X,D(G^\bullet))=D(R\Gamma(\overline X,G^\bullet)).
\]
Appliquant 10.3.5, on obtient l'équation fonctionnelle\auditmark{17}
\[
\tag{10.5.2}
Z(G,t)=
\det(-Ft,R\Gamma(\overline X,G))^{-1}\,
Z(D(G),t^{-1}).
\]

\paragraph{10.6.} Soient
\[
j:U\hookrightarrow X
\]
un ouvert de \(X\), complément d'un ensemble fini \(S\) de points fermés, et \(G^\bullet\in\Dparf(U,\Lambda)\). On note
\[
j_!:\Dparf(U,\Lambda)\longrightarrow\Dparf(X,\Lambda)
\]
le foncteur dérivé du foncteur exact prolongement par zéro, et
\[
Rj_*:\Dparf(U,\Lambda)\longrightarrow\Dparf(X,\Lambda)
\]
le foncteur dérivé du foncteur image directe par \(j\). Le théorème de dualité locale, prouvé notamment en dimension un, donne
\[
\tag{10.6.1}
D(j_!G^\bullet)=Rj_*D(G^\bullet).
\]
Pour \(G\) un faisceau localement constant de \(\Lambda\)-modules projectifs sur \(U\), de dual \(\check G\), on a
\[
D(G)=\check G(1)[2],
\]
et (10.5.2) prend la forme
\[
\tag{10.6.2}
Z(j_!G,t)=
\det(-Ft,R\Gamma(\overline X,j_!G))^{-1}
Z(Rj_*\check G(1),t^{-1}).
\]

\paragraph{10.7.} Nous allons expliciter cette formule. Les fonctions \(Z\) des deux membres sont produits de facteurs locaux. Pour le membre de gauche, le facteur local en \(x\), point fermé de \(X\), est
\[
Z_x(j_!G,t)=
\begin{cases}
\det(1-F_x t^{\deg(x)},G_{\overline x})^{-1},&x\notin S,\\
1,&x\in S.
\end{cases}
\]
Pour le membre de droite, c'est
\[
Z_x(Rj_*\check G(1),t)=
\begin{cases}
\det(1-F_x t^{\deg(x)},\check G_{\overline x}(1))^{-1},&x\notin S,\\
\det(1-F_x t^{\deg(x)},R\Gamma(I_x,\check G_{\overline K}(1)))^{-1},&x\in S.
\end{cases}
\]
Cette dernière formule se lit ainsi: \(\check G_{\overline K}(1)\) définit une représentation de \(W(\overline K_x/K_x)\), à laquelle on applique le foncteur dérivé du foncteur invariants sous \(I\). On obtient ainsi un complexe de \(\Lambda\)-modules projectifs, sur lequel Frobenius agit, et on effectue la construction 10.3. Explicitons.

\begin{enumerate}[label=(\alph*)]
\item Le foncteur invariants sous \(I\) est composé du foncteur invariants sous
\[
P'=\Ker(t_\ell):
\quad I\text{-modules}\longrightarrow I/P'\text{-modules},
\]
et du foncteur invariants sous \(\Z_\ell(1)\). Le premier foncteur est exact et transforme un \(\Lambda\)-module projectif en \(\Lambda\)-module projectif. Pour tout \(I\)-module \(V\), on a
\[
R\Gamma(I,V)=R\Gamma(\Z_\ell(1),V^{P'}).
\]
Pour tout \(\Z_\ell(1)\)-module \(W\), on peut calculer \(R\Gamma(\Z_\ell(1),W)\) en prenant une résolution de \(W\) par des modules coinduits
\[
\mathfrak F(\Z_\ell(1),M),
\]
module des fonctions localement constantes de \(\Z_\ell(1)\) dans \(M\), sur lequel \(\Z_\ell(1)\) agit par \(f(x)\mapsto f(\tau x)\). Pour \(\sigma\) un générateur de \(\Z_\ell(1)\), on a une suite exacte
\[
0\longrightarrow W\longrightarrow\mathfrak F(\Z_\ell(1),W)
\xrightarrow{\alpha_\sigma}
\mathfrak F(\Z_\ell(1),W)\longrightarrow0,
\]

% D016_SOURCE_PHYSICAL_PAGE_085
flèches \(m\mapsto(i\mapsto im)\) et
\[
f\xmapsto{\alpha_\sigma}(i\mapsto f(\sigma i)-\sigma(f(i))).
\]
Passant aux invariants, on identifie \(R\Gamma(\Z_\ell(1),W)\) au complexe
\[
\tag{10.7.1}
W\xrightarrow{\,1-\sigma\,}W.
\]

Quand on change \(\sigma\) en \(\sigma^\nu\), \(\nu\in\Z_\ell^*\), l'unique diagramme commutatif
\[
\begin{tikzcd}[column sep=3.5em,row sep=2.5em]
\mathfrak F(\Z_\ell(1),W)\arrow[r,"\alpha_\sigma"]\arrow[d,equal]&
\mathfrak F(\Z_\ell(1),W)\arrow[r]\arrow[d,"\beta"]&0\\
\mathfrak F(\Z_\ell(1),W)\arrow[r,"\alpha_{\sigma^\nu}"']&
\mathfrak F(\Z_\ell(1),W)\arrow[r]&0
\end{tikzcd}
\]
fournit par passage aux invariants
\[
\tag{10.7.2}
\begin{tikzcd}[column sep=3.8em,row sep=2.6em]
W\arrow[r,"1-\sigma"]\arrow[d,"1"']&
W\arrow[d,"\frac{1-\sigma^\nu}{1-\sigma}"]\\
W\arrow[r,"1-\sigma^\nu"']&W ,
\end{tikzcd}
\]
où, pour \(n\) entier convergeant dans \(\Z_\ell\) vers \(\nu\),
\[
\tag{10.7.3}
\frac{1-\sigma^\nu}{1-\sigma}\,w
=
\lim_n\sum_{i=0}^{n-1}\sigma^iw.
\]
On conclut que \(R\Gamma(\Z_\ell(1),W)\) se représente comme étant n'importe lequel des complexes (10.7.1), ces complexes, pour \(\sigma\) variable, étant identifiés par les isomorphismes (10.7.2).

Quand on part d'un \(W(\overline K/K)\)-module \(V\), Frobenius agit par transport de structure sur
\[
R\Gamma(I,V)=
\bigl[V^{P'}\xrightarrow{\,1-\sigma\,}V^{P'}\bigr].
\]
Quel que soit le Frobenius géométrique \(F\) et \(\sigma\) engendrant \(\Z_\ell(1)\), \(F\) définit
% D016_SOURCE_PHYSICAL_PAGE_086
\[
\begin{tikzcd}[column sep=4.0em,row sep=2.6em]
V^{P'}\arrow[r,"1-\sigma"]\arrow[d,"F"']&
V^{P'}\arrow[d,"F"]\\
V^{P'}\arrow[r,"1-\sigma^{1/q}"']&
V^{P'},
\end{tikzcd}
\]
\auditmark{03}
qui, via (10.7.2), s'identifie à l'automorphisme du complexe
\[
\bigl[V^{P'}\xrightarrow{\,1-\sigma\,}V^{P'}\bigr]
\]
de composantes
\[
\left(F,\ \frac{1-\sigma}{1-\sigma^{1/q}}\circ F\right),
\]
où
\[
\frac{1-\sigma}{1-\sigma^{1/q}}
=
\frac{1-(\sigma^{1/q})^q}{1-\sigma^{1/q}}=\sum_{i=0}^{q-1}\sigma^{i/q}
\]
est défini comme en (10.7.3). On a donc
\[
\det(1-Ft,R\Gamma(I,V))
=
\det(1-Ft,V^{P'})\,
\det\!\left(1-\frac{1-\sigma}{1-\sigma^{1/q}}\circ Ft,V^{P'}\right)^{-1},
\]
où \(F\) désigne un quelconque Frobenius géométrique et \(\sigma\) un quelconque générateur de \(\Z_\ell(1)\).

\par\smallskip\noindent\textit{Notation éditoriale pour (10.7.4).} Pour \(x\in S\), posons
\[
W_x=\bigl(\check G_{\overline K}(1)\bigr)^{P'_x},\qquad
q_x=\#k_x,\qquad
B_x=\left(\sum_{i=0}^{q_x-1}\sigma_x^{i/q_x}\right)F_x.
\]
La représentation générique est ici restreinte au groupe de Weil local en \(x\);
\(F_x\) est un relèvement du Frobenius géométrique et \(\sigma_x\) engendre
\(I_x/P'_x\).\auditmark{04}

La formule 10.6.2 se récrit:

\[
\tag{10.7.4}
\begin{aligned}
&\prod_{x\notin S}\det(1-F_xt^{\deg(x)},G_{\overline x})^{-1}
\\
&\quad=\det(-Ft,R\Gamma(\overline X,j_!G))^{-1}
\\
&\qquad\cdot\prod_{x\notin S}
\det(1-F_xt^{-\deg(x)},\check G_{\overline x}(1))^{-1}
\\
&\qquad\cdot\prod_{x\in S}
\left[\det(1-F_xt^{-\deg(x)},W_x)^{-1}
\det(1-B_xt^{-\deg(x)},W_x)\right].
\end{aligned}
\]

\end{enumerate}

\paragraph{10.8.} Par passage à la limite projective, les résultats précédents se généralisent aux faisceaux \(\ell\)-adiques, ou \(\lambda\)-adiques, et on obtient les énoncés de 9.1. La raison, cf. 10.5.2, de la formule simple (9.1.1) est le résultat local que, pour \(G\) localement constant sur
\[
U\xrightarrow{j}X,
\]
on a
\[
j_*D(G)=D(j_*G).
\]\auditmark{05}
Toutefois, la formule obtenue pour la constante, c'est-à-dire le déterminant de \(-Ft\) agissant sur la cohomologie, n'est pas très explicite, et on n'en a pas de théorie locale.

\paragraph{Remarque 10.9.} Soient \(X\) une courbe projective non singulière sur un corps algébriquement clos \(\overline k\), et \(V\) un faisceau localement constant de \(\Lambda\)-modules libres sur \(X_{\mathrm{\acute et}}\). Pour simplifier, on suppose \(X\) connexe de genre \(g\ge1\) et \(\Lambda\) local. Voici deux méthodes pour attacher à \(V\) un \(\Lambda\)-module libre de rang \(1\).

\begin{enumerate}[label=(\alph*)]
\item Écrivons \(R\Gamma(X,V)\) comme un complexe fini \(K\) de \(\Lambda\)-modules libres. Alors
\[
\det(R\Gamma(X,V))=
\bigotimes_i \det(K^i)^{(-1)^i},
\]
où \(\det\) désigne la puissance extérieure maximale, ne dépend, à isomorphisme canonique près, que de \(R\Gamma(X,V)\). On note \(h(V)\) son dual.

\item Soit
\[
K=\sum_i n_iP_i
\]
un diviseur canonique, \(\mathcal O(K)\simeq\Omega_X^1\), et
\[
N_K(V)=\bigotimes_i \det(V_{P_i})^{\otimes n_i}.
\]
L'ensemble des diviseurs canoniques positifs forme un espace projectif, et les \(N_K(V)\) forment un système local nécessairement trivial sur cet espace projectif. Les \(N_K(V)\), pour \(K\) variable, sont donc canoniquement isomorphes, et on pose
\[
\langle\Omega^1,V\rangle=N_K(V).
\]
\end{enumerate}

Si \((X,V)\) est défini sur \(k\subset\overline k\), c'est-à-dire provient par extension des scalaires de \((X_0,V_0)\) sur \(\Spec(k)\), le groupe de Galois \(\Gal(\overline k/k)\) agit sur \(h(V)\) et \(\langle\Omega^1,V\rangle\) par des caractères \(\chi_h(V)\) et \(\chi_{\Omega}(V)\), à valeurs dans \(\Lambda^*\). Ces constructions se généralisent au cas \(\lambda\)-adique.

\paragraph{Proposition 10.9.1.}\auditmark{06} Supposons vérifiée l'une des conditions suivantes:
\[
\tag{a}
\Lambda=\mathcal O_\lambda\quad\text{et }V\text{ se trivialise sur un revêtement fini de }X;
\]
\[
\tag{b}
V\text{ est de rang }1.
\]
Alors
\[
\tag{10.9.1}
\chi_h(V)=\chi_{\Omega}(V)\,\omega_{(1-g)\dim(V)},
\]
où \(\omega_n\) est le caractère donnant l'action de Galois sur \(\Z_\ell(n)\).

Un argument standard utilisant Tchebotarev nous ramène au cas où \(k\) est un corps fini et où \(\overline k\) est sa clôture algébrique. Il s'agit de montrer que \(\chi_h\) et \(\chi_{\Omega}\omega_{(1-g)\dim V}\) prennent la même valeur sur le Frobenius. Dans le cas (a), le groupe de Weil \(W(\overline K_0/K_0)\) de \(X_0\) agit continûment, pour la topologie discrète, sur \((V_0)_{\overline K_0}\), de sorte qu'on peut appliquer 9.3.

Comparons (9.3) à (10.7.4), pour \(S=\varnothing\). Ce sont des équations fonctionnelles reliant les mêmes fonctions \(Z\). Exprimant qu'elles ont la même constante, on trouve (10.9.1). Il peut être plus commode de ne faire ce calcul que pour \(V\) virtuel de dimension \(0\), et de vérifier (10.9.1) directement pour \(V=\Z_\ell\). Dans le cas (b), on procède de même en utilisant 10.12.1 ci-dessous.

Il serait très intéressant de préciser (10.9.1), en définissant un isomorphisme canonique entre
\[
h(V)
\quad\text{et}\quad
\langle\Omega^1,V\rangle((1-g)\dim(V)),
\]
et de généraliser (10.9.1) au cas ramifié.

\paragraph{Remarque 10.10.} Le théorème 10.4 est sans doute encore valable pour \(\Lambda\) un anneau noethérien de caractéristique \(p\), c'est-à-dire tel que \(p\Lambda=0\). La formule de traces requise devrait résulter de la formule des traces en cohomologie cohérente, la Woods Hole trace formula. La démonstration n'a toutefois pas été rédigée.

\paragraph{10.11. Démonstration de 10.4.} Dans [7], 10.4 n'est démontré que pour des faisceaux \(\ell\)-adiques, car la méthode de démonstration, basée sur la formule des traces de Lefschetz sur \(\overline X\) pour les puissances de l'endomorphisme de Frobenius, requérait que l'on divise par un entier arbitraire pour passer de Lefschetz à (10.4.1), passage par la dérivée logarithmique de \(Z\).

Pour obtenir 10.4 tel quel, il faut utiliser SGA 4, XVII, 5.5, la formule des traces de Lefschetz sur les \(\Sym^n(X)\), et développer en série les deux membres de 10.4.1, le second se développant comme
\[
\det(1-Ft,R\Gamma(\overline X,G))^{-1}
=
\sum_n \Tr(F,\Sym^n(R\Gamma(\overline X,G)))\,t^n,
\]
où \(\Sym^n\) est le foncteur dérivé du foncteur non additif \(n\)-ième puissance symétrique. Par ailleurs, la méthode de démonstration de [7], par fibrations par courbes successives, ramène la formule des traces à démontrer au cas des courbes.

\paragraph{Formule des traces 10.11.1.} Soient \(X\) une courbe projective et lisse sur \(\F_q\), \(U\) un ouvert de \(X\), et \(G\) un faisceau de Weil localement constant de \(\Lambda\)-modules projectifs. On a
\[
\sum_{x\in U(\F_q)}\Tr(F_x,G_{\overline x})
=
\Tr(F,R\Gamma(\overline X,j_!G)).
\]
Dans cette formule, les Frobenius sont relatifs à \(\F_q\), et
\[
\overline X=X\otimes_{\F_q}\overline{\F}_q.
\]
On peut supposer que \(\F_q\) est le sous-corps des constantes du corps des fonctions \(K\) de \(X\), sans quoi les deux membres sont \(0\).

La méthode de démonstration exposée par Grothendieck dans le séminaire oral SGA 5 consiste à se ramener à la formule de Lefschetz sur le nombre de points fixes d'endomorphismes de courbes complètes, déjà prouvée dans [17], en étudiant la cohomologie de revêtements de \(\overline X\) comme représentations du groupe du revêtement.

Soit \(\overline K\) une clôture algébrique de \(K\), contenant le corps des fonctions de \(\overline X\). Le faisceau \(G\) définit une représentation de \(W(\overline K/K)\) sur le \(\Lambda\)-module projectif \(G_{\overline K}\). L'intersection \(J\) de \(\Gal^0(\overline K/K)\) et du noyau de cette représentation définit un revêtement étale connexe
\[
f:\overline U'\longrightarrow \overline U
\]
de \(\overline U\), sur lequel \(G\) devient trivial. Posons
\[
H=W(\overline K/K)/J.
\]
Alors \(H\) est une extension de \(\Z\) par le groupe \(H_0\) du revêtement \(\overline U'/\overline U\), et \(G_{\overline K}\) est une représentation de \(H\).

Soit \(j'\) l'inclusion de \(\overline U'\) dans une courbe projective non singulière \(\overline X'\). La cohomologie \(\ell\)-adique à support propre de \(\overline U'\) est une représentation de \(H\). Des arguments standards permettent de faire mieux, et de définir
\[
R\Gamma_c(\overline U',\Z_\ell)
=
R\Gamma(\overline X',j'_!\Z_\ell)
\in\Dparf(\Z_\ell[H_0]),
\]
sur lequel \(H\) agit de façon semi-linéaire relativement à son action sur \(\Z_\ell[H_0]\), l'action de \(H_0\subset H\) étant l'action évidente. On a canoniquement
\[
\tag{10.11.2}
R\Gamma(\overline X,j_!G)
=
\bigl[
R\Gamma(\overline X',j'_!\Z_\ell)
\otimes_{\Z_\ell}G_{\overline K}
\bigr]^{H_0},
\]
isomorphisme compatible à l'action de Frobenius sur les deux membres.

Écrivons \(R\Gamma(\overline X',j'_!\Z_\ell)\) comme un complexe fini de \(\Z_\ell[H_0]\)-modules projectifs, sur lequel Frobenius agit. Les composantes de ce complexe apparaissent comme des \(\Z_\ell[H^+]\)-modules, pour \(H^+\) le sous-monoïde de \(H\), extension de \(\mathbb N\subset\Z\) par \(H_0\). On leur applique le lemme suivant.

\paragraph{Lemme 10.11.3.}\auditmark{08} Soient \(H^+\) un produit semi-direct de \(\mathbb N\) par un groupe fini \(H_0\), \(M\) un \(\Z_\ell[H^+]\)-module, projectif de rang fini en tant que \(\Z_\ell[H_0]\)-module, \(\Lambda\) une \(\Z_\ell\)-algèbre et \(V\) un \(\Lambda[H^+]\)-module, projectif de rang fini comme \(\Lambda\)-module. Soient \(\chi_M\) et \(\chi_V\) les caractères de \(M\) et \(V\).
\begin{enumerate}[label=(\roman*)]
\item Il existe une fonction \(\chi_M^*\) sur \(H^+\), à valeurs dans \(\Z_\ell\), telle que, pour \(h\in H^+\), et \(H_0(h)\) le centralisateur de \(h\) dans \(H_0\), on ait
\[
\chi_M(h)=\Tr(h,M)=|H_0(h)|\,\chi_M^*(h).
\]
\item Pour \(F\in H^+/H_0\),
\[
\Tr\bigl(F,(M\otimes_{\Z_\ell}V)^{H_0}\bigr)
=
\sum_{\substack{h\mapsto F\\ \mathrm{mod}\ H_0\text{-conjugaison}}}
\chi_M^*(h)\chi_V(h).
\]
\end{enumerate}

Appliquant Lefschetz, on calcule \(\chi_M^*\) pour \(R\Gamma(\overline X',j'_!\Z_\ell)\), virtuellement la différence de \(R\Gamma(\overline X',\Z_\ell)\) et d'une quantité élémentaire; 10.11.2 et 10.11.3 fournissent la formule voulue.





% D016_SOURCE_PHYSICAL_PAGE_091
\paragraph{10.12. Le cas abélien, et une application farfelue.} Soit
\[
\chi:W(\overline K/K)\longrightarrow \Lambda^*
\]
un quasi-caractère à valeurs dans \(\Lambda^*\), non ramifié en dehors d'un ensemble fini \(S\) de places de \(K\). Soit \(j\) l'inclusion de \(U=X-S\) dans \(X\), et soit \(j_![\chi]\) le faisceau de rang un sur \(U\) défini par \(\chi\) et prolongé par \(0\). Posons
\[
\varepsilon(j_![\chi],t)=
\prod_{v\in S}\varepsilon_0(\chi_v\omega_v^t,\psi_v,dx_v)
\prod_{v\notin S}\varepsilon(\chi_v\omega_v^t,\psi_v,dx_v).
\]
\(\omega^t\) est comme en 7.6; \(\varepsilon_0\) est défini en 6.4; \(\varepsilon\), pour un caractère non ramifié, est défini par 3.4.3.3.

\paragraph{Proposition 10.12.1.}\auditmark{09} On a
\[
\det(-Ft,\RGamma(\overline X,j_![\chi]))^{-1}=\varepsilon(j_![\chi],t).
\]
Cette proposition précise 10.7.4.

Soit \(H\) le quotient de \(W(\overline K/K)\) par \(\Ker(\chi)\). Il suffit de traiter le cas universel où
\[
\Lambda=\mathbb Z/\ell^n[H].
\]
Ce cas se relève en caractéristique \(0\), et 10.12.1 s'obtient par réduction.

\paragraph{10.13.} Prenons par exemple pour \(\Lambda\) les nombres duaux sur \(\mathbb Z/\ell^n\):
\[
\Lambda=\mathbb Z/\ell^n[\epsilon]\qquad(\epsilon^2=0),
\]
pour \(\alpha\) un homomorphisme continu du groupe des classes d'idèles \(\bbA^*/K^*\) dans le groupe additif \(\mathbb Z/\ell^n\), et posons
\[
\chi=1+\alpha\epsilon.
\]
La ramification de \(\alpha\) est automatiquement modérée. Soit \(S\) un ensemble fini de places en dehors desquelles \(\alpha\) est non ramifié. Appliquant 10.4 pour \(\chi\) et pour le caractère trivial, et faisant le rapport, on trouve que

% D016_SOURCE_PHYSICAL_PAGE_092
\[
\tag{10.13.1}
L_S^+(\alpha,t)=
\sum_{v\notin S}\alpha_v(\pi_v)\frac{t^{\deg(v)}}{1-t^{\deg(v)}}
=
\sum_{r\ge1}t^r\sum_{\substack{v\notin S\\\deg(v)\mid r}}\alpha_v(\pi_v)
\]
est dans \(\mathbb Z/\ell^n(t)\).

Pour toute place \(v\) de \(K\), choisissons une uniformisante \(\pi_v\). Adjoignons à \(\mathbb Z/\ell^n\) une racine primitive \(p\)-ième de l'unité, et soit \(\psi\) un caractère additif non trivial de \(\bbA/K\), à valeurs dans les racines \(p\)-ièmes de l'unité. Prenons le quotient de (10.7.4) pour \(\chi\) et de (10.7.4) pour le caractère trivial, et utilisons (10.12.1). Après simplifications, on obtient l'identité suivante. Dans le premier membre, les deux premières sommes infinies sont sommées par (10.13.1). Les sommes portent toutes sur toutes les places de \(K\):
\par\smallskip\noindent\textit{Notation éditoriale.} Ici \(q_v=\#k_v\), \(n_v=n(\psi_v)\);
\(\widetilde x\in\mathcal O_v\) relève \(x\in k_v^*\), et
\(\alpha_v(x)\) désigne le caractère sur les unités modulo les unités principales.
L'argument \(\pi_v^{-1-n_v}\widetilde x\) définit un caractère additif non trivial
du corps résiduel, indépendant du relèvement.\auditmark{10}

\[
\tag{10.13.2}
\begin{aligned}
&\sum_v\alpha_v(\pi_v)\frac{t^{\deg(v)}}{1-t^{\deg(v)}}
+\sum_v\alpha_v(\pi_v)\frac{q_v^{-1}t^{-\deg(v)}}{1-q_v^{-1}t^{-\deg(v)}}
\\
&\quad+\sum_v\alpha_v(-1)\frac{1}{1-t^{\deg(v)}}
\\
&=\sum_v n_v\alpha_v(\pi_v)
+\sum_v\sum_{x\in k_v^*}\alpha_v(x)\psi_v(\pi_v^{-1-n_v}\widetilde x).
\end{aligned}
\]
On vérifie de façon élémentaire que la validité de (10.13.2) ne dépend pas du choix des uniformisantes \(\pi_v\), en se rappelant que, pour \(x\in k_v^*\), on a \((q_v-1)\alpha_v(x)=0\). Dans le premier membre, les \(\alpha_v(-1)\) sont d'ordre \(2\), donc nuls si \(\ell\ne2\).


% D016_SOURCE_PHYSICAL_PAGE_093
\section*{11. Appendice. Le calcul de \(\varepsilon\) modulo les racines de l'unité}

Cet appendice est une variation sur un thème de Lakkis et Dwork. Mon attention a été attirée sur leurs résultats par J.-P. Serre.

Nous noterons \(a\sim b\) la relation suivante entre nombres complexes inversibles:
\[
a\sim b\quad:\quad a/b\text{ est produit d'une racine de l'unité par une puissance de }q^{1/2}.
\]
Soit \(K\) un corps local non archimédien et reprenons les notations de 2.2.1, 2.2.2 et 3.1. La mesure de Haar \(dx\) sera supposée telle que
\[
\int_{\OO}dx=q^{n/2},
\]
avec \(n\) entier. Soient \(V\) une représentation complexe, d'un quotient fini, de \(\Gal(\overline K/K)\), et \(V^P\) les invariants sous l'inertie sauvage. C'est une sous-représentation de \(V\).

\paragraph{Théorème 11.1.} On a
\[
\varepsilon(V,\psi,dx)\sim\varepsilon(V^P,\psi,dx).
\]

\paragraph{Lemme 11.2.} Modulo l'équivalence \(\sim\), \(\varepsilon(V,\psi,dx)\) est indépendant de \(\psi\) et de \(dx\). Cela résulte aussitôt de 5.3 et 5.4. Ce lemme nous permettra d'abréger \(\varepsilon(V,\psi,dx)\) en \(\varepsilon(V)\).

\paragraph{Lemme 11.3.} Si le caractère de \(V\) est réel, c'est-à-dire si \(V\) est isomorphe à \(V^*\), alors
\[
\varepsilon(V)\sim1.
\]
On peut supposer que \(\int_{\OO}dx=1\) et que \(n(\psi)=0\). On applique alors 5.7.1, 5.4 et 5.5.

Soient \(L\) une extension finie de \(K\) dans \(\overline K\), \(W\) une représentation de \(\Gal(\overline K/L)\), et \(\Ind(W)\) la représentation induite de \(\Gal(\overline K/K)\).

\paragraph{Lemme 11.4.} On a
\[
\varepsilon_L(W)\sim\varepsilon_K(\Ind(W)).
\]
Soit \(W'\) la représentation triviale de même dimension que \(W\). Pour \(\psi_L=\psi_K\circ\Tr_{L/K}\), on a par (5.6.1)
\[
\varepsilon(W-W',\psi_L)=\varepsilon(\Ind(W)-\Ind(W'),\psi_K),
\]
et on conclut en appliquant 11.3 à \(W'\) et à \(\Ind(W')\).

\paragraph{11.5.} Soient \(L\) et \(W\) comme ci-dessus, \(P_L\) le groupe d'inertie sauvage de \(L\), et regardons \(W^{P_L}\) comme une représentation de
\[
H=P\cdot\Gal(\overline K/L),
\]
triviale sur \(P\). On a
\[
\Ind(W)^P=\Ind_H^{\Gal(\overline K/K)}(W^{P_L}).
\]
D'après 11.4, le théorème pour \(W\) équivaut donc au théorème pour \(\Ind(W)\); par Brauer, ceci nous ramène au cas où \(V\) est de dimension un, défini par un caractère d'ordre fini \(\chi\) de \(K^*\). Si \(\chi\) est modéré, \([\chi]^P=[\chi]\), et 11.1 est trivial. Si \(\chi\) est sauvagement ramifié, \([\chi]^P=0\), et 11.1 résulte du lemme suivant.

\paragraph{Lemme 11.6.} Pour \(\chi\) sauvagement ramifié,
\[
\varepsilon(\chi,\psi,dx)\sim1.
\]
Soient \(m\) le conducteur de \(\chi\),
\[
n=\left[\frac{m+1}{2}\right],
\]
et \(y\) l'élément de \(K^*\), de valuation \(-m-n(\psi)\), bien défini modulo \((\pi^{-n-n(\psi)})\), tel que, cf. 4.16,
\[
\chi(1+a)=\psi(ay)\qquad(v(a)\ge n).
\]
Décomposons \(\chi\) en \(\chi=\chi_1\chi_2\), avec \(\chi_1\) d'ordre premier à \(p\), et \(\chi_2\) d'ordre une puissance de \(p\). Comme en 4.16, on prouve que
\noindent\auditmark{11}
\[
\tag{11.6.1}
\varepsilon(\chi,\psi,dx)
=
\int \chi^{-1}(x)\psi(x)\,dx
=
\int_{y(1+(\pi^{m-n}))}\chi^{-1}(x)\psi(x)\,dx
\]
\[
=
\chi_1^{-1}(y)
\int_{y(1+(\pi^{m-n}))}\chi_2^{-1}(x)\psi(x)\,dx.
\]
Si \(m\) est pair, \(m=2n\), la fonction intégrée est constante et 11.6 est clair. Pour \(m\) impair, \(m=2n-1\), voici deux façons de procéder.

(a) \(\chi^{-1}(x)\psi(x)\) est un caractère quadratique non dégénéré sur
\[
y(1+(\pi^{m-n}))/(1+(\pi^n));
\]
d'après Weil, la somme de ses valeurs est \(\sim1\).

(b) D'après (11.6.1), il suffit de prouver que \(\varepsilon(\chi_2,\psi,dx)\sim1\). Le carré de ce nombre, multiplié par une puissance convenable de \(q\), appartient à un corps de racines \(p^N\)-ièmes de l'unité, et est de valuation \(1\) en toutes les places de ce corps, sauf peut-être en l'unique place divisant \(p\). C'est donc une racine de l'unité.

\paragraph{Corollaire 11.7.} Prenons \(\psi\) et \(dx\) tels que \(n(\psi)=-1\) et que
\[
\int_{(\pi)}dx=1.
\]
Alors \(\varepsilon(V,\psi,dx)\) est un entier algébrique.

Décomposant \(V\) en \(V^P+W\), on voit qu'il suffit de traiter les deux cas suivants.

(a) \(V\) est modérément ramifiée. Dans ce cas, \(V\) est une somme de représentations induites de représentations de dimension un de sous-groupes, et \(\varepsilon\) est, au signe près, un produit de sommes de Gauss, cf. 5.10.

(b) \(V^P=0\). Dans ce cas, \(\varepsilon\sim1\), et il suffit de prouver que la valeur absolue complexe \(|\varepsilon|\) est une puissance positive de \(q\). Ceci résulte de 5.7.2, où \(dx'/dx=q^{-1}\)\auditmark{13}.


% D016_SOURCE_PHYSICAL_PAGE_096
\begin{center}
\bfseries BIBLIOGRAPHIE
\end{center}

\begin{enumerate}[label={[\arabic*]},leftmargin=2.8em]
\item E. Artin, \emph{Zur Theorie der L-Reihen mit allgemeinen Gruppencharakteren}, Hamb. Abh. \textbf{8} (1930), 292--306; Collected Papers, 165--179.
\item E. Artin and J. Tate, \emph{Class field theory}, Benjamin, New York.
\item P. Deligne, \emph{Les constantes des équations fonctionnelles}, Séminaire Delange--Pisot--Poitou 1969/70, 19 bis.
\item P. Deligne, \emph{Formes modulaires et représentations de \(GL(2)\)}, ce volume.
\item H. Davenport and H. Hasse, \emph{Die Nullstellen der Kongruenzzetafunktionen in gewissen zyklischen Fällen}, J. reine angewandte Math. \textbf{172} (1935), 151--182.
\item B. Dwork, \emph{On the Artin root number}, Amer. J. Math. \textbf{78} (1956), 444--472.
\item A. Grothendieck, \emph{Formule de Lefschetz et rationalité des fonctions L}, Séminaire Bourbaki 279, décembre 1964.
\item H. Jacquet and R. P. Langlands, \emph{Automorphic forms on GL(2)}, Lecture Notes in Math. 114, Springer-Verlag, 1970.
\item R. P. Langlands, \emph{On the functional equation of the Artin L-functions}, notes polycopiées, Yale University, incomplet.
\item J. P. Serre, \emph{Corps locaux}, Hermann, Paris, 1962.
\item J. P. Serre, \emph{Représentations linéaires des groupes finis}, 2\textsuperscript{e} éd. refondue, Hermann, Paris, 1971.
\item J. P. Serre, \emph{Abelian \(\ell\)-adic representations and elliptic curves}, New York, 1958.\auditmark{59}
\item J. P. Serre and J. Tate, \emph{Good reduction of abelian varieties}, Ann. of Math. \textbf{88}, 3 (1968), 492--517.

% D016_SOURCE_PHYSICAL_PAGE_097
\item J. Tate, \emph{Fourier analysis in number fields and Hecke's zeta functions}, thesis, in \emph{Algebraic Number Theory}, edited by Cassels and Fröhlich, Academic Press, 1967.
\item A. Weil, \emph{Basic Number Theory}, Springer-Verlag, 1967.
\item A. Weil, \emph{Dirichlet series and automorphic forms}, Lecture Notes in Math. 189, Springer-Verlag, 1971.
\item A. Weil, \emph{Variétés abéliennes et courbes algébriques}, Hermann, Paris, 1948.
\item A. Weil, \emph{Number of solutions of equations in finite fields}, Bull. Amer. Math. Soc. \textbf{55} (1949), 497--508.
\item A. Weil, \emph{Sur la théorie du corps de classes}, J. of the Math. Soc. of Japan \textbf{3} (1951).
\item[SGA] Séminaire de géométrie algébrique du Bois-Marie. SGA 4 est publié dans les Lecture Notes 269, 270, 305. SGA 5 est diffusé par l'I.H.E.S.
\item P. X. Gallagher, \emph{Determinants of representations of finite groups}, Abh. Math. Seminar Univ. Hamburg \textbf{28}, 3/4 (1965), 162--167.
\end{enumerate}

\medskip
\noindent\textit{[Note postérieure reproduite dans le témoin.]}\auditmark{15}

\noindent Le problème mentionné dans l'introduction, p. 4, de prouver la conclusion de 9.9 pour \emph{une} représentation \(\lambda\)-adique a été résolu:

\noindent G. Laumon, \emph{Les constantes des équations fonctionnelles des fonctions L sur un corps global de caractéristique positive}, C.R. Acad. Sci. Paris \textbf{298}, 8 (1984), 181--184.

\medskip
\noindent Un bel exposé du théorème d'existence des constantes locales est donné dans:

\noindent J. Tate, \emph{Local constants}, in \emph{Algebraic Number Fields}, edited by A. Fröhlich, Academic Press, 1977, 89--137.


% D016_SOURCE_EOF_PHYSICAL_097


% BEGIN EDITORIAL APPARATUS -- OUTSIDE SOURCE EXTENT
\begingroup
\fontsize{10pt}{11pt}\selectfont
\setlength{\parskip}{0pt}
\setlength{\abovedisplayskip}{6pt plus 2pt minus 2pt}
\setlength{\belowdisplayskip}{6pt plus 2pt minus 2pt}
\setlist[description]{parsep=0pt,itemsep=0pt}
\clearpage
\section*{Notes éditoriales de cette révision}
\noindent\textit{Appareil éditorial, hors texte de l'auteur.}
Les notes distinguent les lectures du témoin, les réparations de transcription,
les précisions de notation et les émendations mathématiques déduites ici.
Ces dernières ne sont pas des errata attribués à Deligne. Les renvois de page
ci-dessous désignent les pages physiques du témoin (ajouter 500 pour la pagination imprimée).
La comparaison cumulée des textes français et anglais au témoin couvre les 97 pages,
y compris l'appendice, la bibliographie et la note postérieure. Les émendations
restent identifiées séparément des lectures imprimées ; le témoin est conservé inchangé.

\begin{description}[style=nextline,leftmargin=2.6em,labelwidth=2.2em,font=\normalfont\bfseries]
\item[E01] \textbf{Source p.~82 : renvoi à la démonstration.} Le témoin renvoie à 10.8 ; la démonstration porte le titre 10.11, p.~88. Le renvoi est donc corrigé en 10.11. Le déterminant déjà réparé dans (10.3.4) et la fibre géométrique après (10.4.1) sont conservés. Ce dernier affichage n'est pas numéroté : il ne faut pas le confondre avec l'identité globale (10.4.1).

\item[E02] \textbf{Source p.~84 : fibre générique et condition de branche.} Le maître hérité omettait la barre sur \(K\) dans les deux occurrences de \(\check G_{\overline K}(1)\) ; la lecture du témoin est rétablie. La condition \(x\in S\) de la branche ramifiée était déjà correcte dans le maître livré au Prompt 15 et est conservée. Le témoin répète \(x\notin S\), ce qui ne convient pas à la branche de cohomologie d'inertie. Aucune nouvelle modification de cette condition n'est revendiquée ici.

\item[E03] \textbf{Source p.~86 : Frobenius géométrique. Émendation mathématique.} Le témoin porte \(1-\sigma^q\) dans la différentielle inférieure et \((1-\sigma)/(1-\sigma^q)\) dans l'opérateur de degré un. Les définitions 2.2.3--2.2.4 fixent cependant le Frobenius géométrique, donc \(F\sigma F^{-1}=\sigma^{1/q}\). L'opérateur retenu est
\[
 B_\sigma=\frac{1-\sigma}{1-\sigma^{1/q}}F=\left(\sum_{i=0}^{q-1}\sigma^{i/q}\right)F.
\]
Le quotient désigne la somme finie et n'inverse pas \(1-\sigma\). Les puissances ont un sens puisque \(q\in\Z_\ell^*\). L'identité \(F(1-\sigma)=(1-\sigma^{1/q})F\) donne \(B_\sigma(1-\sigma)=(1-\sigma)F\), condition de morphisme de complexes. Pour l'inertie triviale, \(B_\sigma=qF\), conformément au twist de \(H^1(I,V)\). Le diagramme, l'identité d'opérateurs, le déterminant et le facteur de (10.7.4) sont corrigés ensemble. Il s'agit d'une déduction éditoriale, non d'un erratum attribué à l'auteur.

\item[E04] \textbf{Source p.~86 : formule (10.7.4).} Le témoin omet l'exposant \(-1\) dans le produit eulérien de gauche ; il est rétabli d'après (10.6.2) et la définition de \(Z_x\). Le témoin et le maître hérité emploient une notation de fibre fermée aux points \(x\in S\), alors que le faisceau est donné seulement sur \(U\). Les notations explicitement éditoriales \(W_x\) et \(B_x\) précisent la représentation générique restreinte au groupe de Weil local et ses invariants sous \(P'_x\). Elles évitent aussi de lire les \(q\), \(F\) et \(\sigma\) locaux comme des choix globaux communs. La base géométrique de \(R\Gamma(\overline X,j_!G)\) est explicitée.

\item[E05] \textbf{Source p.~87 : dualité dans 10.8.} Le maître hérité porte \(j_*D(G)=D(j_!G)\). Le témoin porte \(j_*D(G)=D(j_*G)\), lecture rétablie. Les deux occurrences de \(j_*\) sont non dérivées dans cet affichage. La formule distincte \(D(j_!G)=Rj_*D(G)\) de 10.6.1 n'est pas substituée à celle-ci. Le contexte est le cas \(\ell\)-adique introduit en 10.8.

\item[E06] \textbf{Source pp.~87--88 : proposition 10.9.1.} L'hypothèse du témoin est \(\Lambda=\mathcal O_\lambda\), non \(\Q_\lambda\) comme dans le maître hérité. L'anneau des coefficients est rétabli. Les noms \(k\subset\overline k\) du témoin remplacent le renommage \(k_0\subset k\). Le numéro d'équation (10.9.1), ajouté par la composition héritée, est conservé : le texte cite à plusieurs reprises l'égalité sous ce numéro, absent à côté de l'affichage imprimé. Les notations \(\chi_\Omega\), \((V_0)_{\overline K_0}\) et le twist de Tate sont conservés.

\item[E07] \textbf{Source p.~89 : notation de la trace.} Le témoin affiche \(G_x\) dans la formule des traces 10.11.1. La notation conservée \(G_{\overline x}\) explicite la fibre géométrique. C'est une précision de notation, non la restauration littérale d'une barre visiblement imprimée. Le précédent compte rendu classait ce changement de façon trop forte. La puissance symétrique de l'affichage précédent reste le foncteur dérivé du foncteur non additif puissance symétrique.

\item[E08] \textbf{Source p.~90 : lemme 10.11.3. Qualification de domaine.} Les hypothèses donnent des \(H^+\)-modules, tandis que les conclusions imprimées emploient \(H\), \(h\in H\) et \(F\in H/H_0\). Les conclusions sont ici restreintes à \(H^+\) et à \(H^+/H_0\). L'action du générateur d'un monoïde sur un module projectif peut ne pas être inversible : elle peut être nulle si \(H_0=1\). Son prolongement au groupe \(H\) serait donc une hypothèse supplémentaire au niveau des termes du complexe. Lorsque les actions se prolongent, la formulation sur le groupe s'applique aussi. Cette restriction est une qualification éditoriale explicite. Le témoin énonce le lemme puis l'applique ; aucune démonstration séparée n'a été ajoutée ou attribuée à Deligne.

\item[E09] \textbf{Source p.~91 : proposition 10.12.1. Émendation mathématique.} Le témoin imprime le déterminant sans exposant \(-1\). Les équations (10.5.2) et (10.6.2) identifient le multiplicateur de l'équation fonctionnelle à l'inverse du déterminant ; cet exposant est retenu ici. La base géométrique est explicitée. Pour le faisceau constant sur \(\mathbf P^1\), le déterminant cohomologique est \(qt^2\), tandis que le multiplicateur est \((qt^2)^{-1}\).

\item[E10] \textbf{Source p.~92 : formules (10.13.1)--(10.13.2).} Dans (10.13.1), le \(v\in S\) hérité est corrigé en \(v\notin S\), lecture imprimée. La même restriction est explicitée dans la somme développée, dont l'indice muet devient \(r\) pour le distinguer de l'exposant du module des coefficients. Dans (10.13.2), le témoin porte \(+\sum_v n_v\alpha_v(\pi_v)\) ; le maître hérité avait un signe moins. Ce signe plus est une restauration de la source. Trois autres changements sont des émendations mathématiques : \(q_v\) devient \(q_v^{-1}\) dans la deuxième somme ; le signe moins devant la somme de Gauss devient plus ; \(\pi_v^{1+n_v}\) devient \(\pi_v^{-1-n_v}\).

Voici le calcul local. Posons \(a_v=\alpha_v(\pi_v)\), \(\overline\psi_v(x)=\psi_v(\pi_v^{-1-n_v}\widetilde x)\) et \(T_v=\sum_{x\in k_v^*}\alpha_v(x)\overline\psi_v(x)\). L'exposant résulte de la convention de conducteur de 3.4. Sur la couronne de valuation \(-1-n_v\), 3.4.3.4 donne
\[
 \frac{\varepsilon_0(1+\alpha_v\epsilon,\psi_v,dx_v)}{\varepsilon_0(1,\psi_v,dx_v)}=1+\epsilon\bigl((n_v+1)a_v+T_v\bigr),
\]
car \(\sum_{x\in k_v^*}\overline\psi_v(x)=-1\). Le signe moins du caractère inverse est divisé par \(-1\), d'où \(+T_v\). Hors de \(S\), 5.9 donne le coefficient \(n_va_v\). Le Frobenius géométrique sur le twist dual a pour scalaire \(q_v^{-1}\). Pour \(z=t^{\deg(v)}\), \(r=q_v^{-1}\) et \(c_v=\alpha_v(-1)\), le quotient des facteurs duaux ramifiés a pour coefficient
\[
 -a_v\frac{rz^{-1}}{1-rz^{-1}}+(a_v+c_v)\frac{z^{-1}}{1-z^{-1}}.
\]
L'adjonction des places omises aux deux sommes eulériennes retranche \(\sum_{v\in S}a_v\) et transporte \(\sum_v c_v/(1-z)\) à gauche. On obtient l'identité sur toutes les places. La valeur \(c_v\) est annulée par \(2\), sans être nécessairement d'ordre exactement \(2\). Si \(\pi'_v=u\pi_v\), alors \(a'_v=a_v+\alpha_v(u)\) et \(T'_v=T_v-(n_v+1)\alpha_v(u)\) ; les deux membres varient de \(-\alpha_v(u)\), en utilisant \((q_v-1)\alpha_v(u)=0\). Les contrôles algébriques vérifient aussi le caractère degré sur \(\mathbf P^1\), pour lequel l'identité corrigée est \(L(t)+L(q^{-1}t^{-1})=-2\). Ces calculs sont éditoriaux, non des passages supplémentaires de l'auteur.

\item[E11] \textbf{Source p.~94 : lemme 11.6.} Le témoin détermine \(y\) modulo \((\pi^{-n-n(\psi)})\) ; l'expression héritée \(\pi^{n-m-n(\psi)}\) est corrigée. Pour un conducteur impair, ces idéaux sont distincts. Dans (11.6.1), le témoin imprime \(\chi_1(y)\), alors que \(\chi^{-1}=\chi_1^{-1}\chi_2^{-1}\) impose \(\chi_1^{-1}(y)\) sur la classe d'intégration. Ce dernier changement est une émendation mathématique identifiée. Le premier symbole intégral conserve la convention d'intégrale oscillante de 3.4.3.2 et n'affirme pas la convergence absolue sur \(K^*\).

\item[E12] \textbf{Source p.~95 : vocabulaire de valuation.} L'expression « de valuation 1 » est lue au sens multiplicatif : valeur absolue locale \(1\). Une lecture en valuation additive n'exprimerait pas l'argument des unités et de la formule du produit. La traduction anglaise explicite ce sens ; le français conserve la formulation du témoin avec la présente qualification.

\item[E13] \textbf{Source p.~95 : corollaire 11.7. Émendation mathématique.} Le témoin porte \(dx'/dx=q\). Avec \(n(\psi)=-1\) et \(\operatorname{vol}_{dx}(\pi\mathcal O)=1\), on a \(\operatorname{vol}_{dx}(\mathcal O)=q\) et \(\mathcal O^\perp=\pi\mathcal O\). L'inversion de Fourier impose
\[
 \frac{dx'}{dx}=\frac{1}{\operatorname{vol}_{dx}(\mathcal O)\operatorname{vol}_{dx}(\mathcal O^\perp)}=q^{-1}.
\]
La substitution dans (5.7.2) donne \( |\varepsilon(V,\psi,dx)|^2=q^{a(V)}\). Le rapport des mesures est corrigé, non l'énoncé d'intégralité.

\item[E14] \textbf{Source pp.~96--97 : bibliographie.} Les vingt références numérotées sont conservées. Le titre de [18] porte « Number of solutions of equations in finite fields », non « Numbers ». Le volume de [19] est \(3\), non le \(1\) hérité. L'ordre du témoin est [19], l'entrée SGA non numérotée, puis [20]. L'intervalle de six pages du Prompt 16 s'arrête après [13], p.~96 ; [14]--[20], SGA et la note postérieure constituent l'intervalle terminal distinct d'une page. La bibliographie n'est pas modernisée silencieusement.

\item[E15] \textbf{Source p.~97 : note postérieure du témoin.} Le témoin contient une note citant Laumon (1984) et Tate (1977). Le titre entre crochets qui l'identifie comme postérieure est éditorial ; il ne figure pas dans le témoin. L'accent mis sur « une » représentation et le \(\lambda\) minuscule de la source sont rétablis. La note appartient au témoin consulté, mais n'est pas présentée comme contemporaine de la publication de 1973. La source se termine après la référence à Tate, pp.~89--137. Le présent appareil éditorial commence seulement après ce point terminal.

\item[E16] \textbf{Source p.~28 : formule (3.4.3.3). Correction de transcription.} Le facteur imprimé est \((-\chi(\pi))^{-1}\). Le signe moins, omis dans le texte hérité, est rétabli après comparaison directe avec le témoin, p.~528. Il ne s'agit pas d'une émendation de la source.
\item[E17] \textbf{Source pp.~82--84 : notations conservées de l'édition de travail.} À la p.~82, après le choix de \((K',u')\), le texte imprime « du complexe \(K\) » ; l'édition lit \(K'\), le complexe sur lequel agit \(u'\). La formule (10.3.3) conserve la notation abrégée \((u,K)\) du témoin. Dans (10.5.2), p.~83, le témoin imprime \(R\Gamma(G,t)\). La lecture conservée \(R\Gamma(\overline X,G)\) est une correction éditoriale des arguments, conforme à (10.4.1), (10.4.2) et (10.5.1), non une transcription littérale de cet affichage. Dans (10.6.2), la même base géométrique est explicitée là où le témoin abrège en \(R\Gamma(j_!G)\). Le titre imprimé « 0.7 » à la p.~84 est normalisé en 10.7. Ces lectures figuraient déjà dans la révision~02 ; la révision~03 les documente sans en modifier le contenu mathématique.
\item[E18] \textbf{Source p.~63 : lemme 7.11.5. Émendations mathématiques.} Les deux sommes dont les bornes sont explicitement imprimées vont de \(1\) à \(q-1\) ; l'édition les fait commencer à \(0\). Après réduction, \(\sigma\) est unipotent : la somme corrigée a pour seule valeur propre \(q\), comme l'affirme le témoin, tandis que la somme imprimée a pour seule valeur propre \(q-1\), qui n'est pas nécessairement inversible. L'opérateur de degré un imprimé fait en outre intervenir l'inverse de la somme. Le Frobenius géométrique agit cependant sur un homomorphisme croisé par \((Fc)(\sigma)=F c(F^{-1}\sigma F)=F c(\sigma^q)\). Or \(c(\sigma^q)=\sum_{i=0}^{q-1}\sigma^i c(\sigma)\). Son opérateur de degré un est donc \(F_1=F\sum_{i=0}^{q-1}\sigma^i\), conformément à E03. Lorsque l'inertie agit trivialement, on obtient \(qF\), comme l'exige la formule \(H^1=W_I(-1)\) du témoin, en (7.11.3). Les changements de borne inférieure et la suppression de l'inverse sont des émendations mathématiques, non des réparations de transcription.

\item[E19] \textbf{Source pp.~4, 61, 64, 75, 80 : changement de variable du twist. Émendation mathématique.} L'argument imprimé \(pt^{-1}\) de l'introduction (B), de 7.11(iii), (7.11.8), de son identité locale de déterminants, de (9.3.1), (9.9.1) et (9.9.2) est remplacé par \(p^{-1}t^{-1}\). Les définitions 2.2.3 et 2.3 fixent le Frobenius géométrique, dont la valeur propre sur le twist de Tate est \(q_v^{-1}=p^{-\deg(v)}\). Ainsi \(Z(W^*(1),t^{-1})=Z(W^*,p^{-1}t^{-1})\), conformément à (9.1.1) et (10.6.2). La lecture du témoin est conservée dans le PDF source joint.

\item[E20] \textbf{Source pp.~64, 76 : données de dualité omises. Émendation mathématique.} Le facteur de droite imprimé en (7.11.8) est \(Z(Rj_*V,pt^{-1})\) ; le facteur retenu est \(Z(Rj_*V^*,p^{-1}t^{-1})\). Celui de (9.3.2) est \(Z(Rj_*W,t)\) ; le facteur retenu est \(Z(Rj_*W^*,p^{-1}t^{-1})\). Tous deux résultent de la formulation explicite avec dual et twist de Tate en (10.6.2), jointe à E19. Il ne s'agit pas de lectures littérales de ces deux formules imprimées.

\item[E21] \textbf{Source p.~64 : identité locale simplifiée de déterminants. Émendation mathématique.} Le membre de gauche imprimé porte \(F_v^{-1}\), remplacé ici par \(F_v\). D'après le lemme 7.11.4, \(H^1(I_v,\omega_1 V^*)\) est dual de \(H^0(I_v,V)\). Si \(A\) désigne Frobenius sur ce dernier espace, l'identité requise est \(\det(1-A)=\det(-A)\det(1-A^{-1})\). Le membre de droite est inchangé. C'est l'identité de déterminants sous-jacente à la formule locale précédente, non un résultat nouveau attribué au témoin.
\item[E22] \textbf{Source p.~65 : sous-groupe de passage au quotient. Émendation mathématique.} Le témoin porte \(v\in S\) dans la condition \(I_v\subset J\) ; l'édition lit \(v\notin S\). À la p.~62, \(S\) contient les places ramifiées de \(V\). Si \(\rho\) est la représentation correspondante, l'absence de ramification hors de \(S\) donne \(I_v\subset\ker\rho\) pour \(v\notin S\). La continuité pour la topologie discrète et la dimension finie rendent \(\ker\rho\) ouvert ; il est distingué. On peut donc prendre \(J=\ker\rho\). La condition imprimée annulerait au contraire l'inertie aux places où la ramification est permise. La réduction aux groupes finis reste celle de 4.10 ; aucune finitude de \(W(\overline K/K)/J\) n'est déduite ici.
\item[E23] \textbf{Source p.~67 : diagramme de 8.3.4. Réparation de transcription.} Le diagramme hérité inversait les flèches verticales et omettait les cinq objets nuls terminaux. Le sens montant du témoin est rétabli : de \(\mathbb G_a(1)\) vers les extensions, puis vers \(I\) et \(W(\overline K/K)\), et enfin vers \(0\). Les deux lignes \(\Z\) sont reliées par une égalité, de même que les deux copies inférieures de \(\mathbb G_a(1)\). Les objets et les flèches horizontales restants sont conservés. Il ne s'agit pas d'une émendation mathématique du témoin.
\item[E24] \textbf{Source pp.~67--68 : notations de 8.3--8.4.} L'inclusion \(i\) du diagramme et \(i(-t_\ell(\sigma))\), mal transcrit avec \(t\), sont rétablis. Les précisions héritées suivantes sont conservées et signalées : \(\mathbb Z_\ell(1)\) pour \(\mathbb Z_\ell\), conformément à 2.2.2, p.~22 ; un représentant \(\sigma\in I\) de \(\bar\sigma\in I/J\), là où le témoin répète \(I/J\) ; \(N'\) dans tout 8.4.1, là où le témoin alterne \(N'\) et \(N\). L'étiquette verticale \(t_\ell\), ajoutée dans le diagramme hérité, explicite l'application du texte. Ces précisions ne sont pas des lectures littérales.
\item[E25] \textbf{Source p.~8 : transfert et trace.} Les \([G:H]\) feuillets, l'application \(f:B_H\to B_G\), le transfert en homologie et sa transposée, omis ou confondus dans le texte hérité, sont rétablis. \textit{Émendation mathématique distincte :} le témoin imprime \(f_!:H^1(B_G,X)\to H^1(B_H,X)\) et \(\eta\in H^1(B_G,K^*)\). La transposée de \(H_1(B_G)\to H_1(B_H)\) va en sens inverse : l'édition lit \(H^1(B_H,X)\to H^1(B_G,X)\) et \(\eta\in H^1(B_H,K^*)\). Cela place \(L\) sur \(B_H\), et sa norme sur \(B_G\), comme l'exige la construction. L'attribution à Gallagher, imprimée en anglais dans le témoin français, est conservée.
\item[E26] \textbf{Source pp.~11--12 : formules omises.} L'identité du noyau dans la preuve de 1.8, le domaine et le codomaine de (1.9.5), et l'inflation (1.9.7) sont rétablis. Les indices \(K\) et \(k\), ainsi que les crochets \([\chi']-[\chi'']\), distinguent les corps de coefficients et les classes de représentations. Ce sont des réparations de transcription, non des émendations de ces formules. La convention typographique héritée \(R^+\) pour le \(R_+\) du témoin est conservée dans toute l'édition.
\item[E27] \textbf{Source pp.~14--18 : relations de groupes finis.} La classe \([1]\) du type I, la restriction \(\chi|_{H\cap A_\mu}\) du type III, les égalités de 1.13.1, \((a_i)\), (b) et (4), le quotient \(G_\mu/\Ker(\mu)\), l'application induite \(\bar u\), ainsi que la définition de \(S\) et la décomposition de \(R\) en 1.14 sont rétablis d'après les pages imprimées. Le texte hérité omettait ou déformait ces formules ; il ne s'agit pas d'émendations mathématiques du témoin.
\item[E28] \textbf{Source p.~18 : domaine de \(\chi\).} Le témoin imprime « \(\chi\) un caractère de \(G\) » au début de la preuve de 1.14.1 ; cette lecture est conservée. Pour établir que \(N(G)\) est un idéal en faisant varier les générateurs \((H,\chi)\) de \(R^+(G)\), \(\chi\) doit pouvoir être un caractère de \(H\). La preuve qui suit utilise cette généralité. La tension est signalée sans correction silencieuse de la phrase imprimée.
\item[E29] \textbf{Source p.~19 : relation induite du cas (B2).} L'affichage \(\sum_i n_i\Ind_{H_i}^{G_\mu}(\chi_i)=0\) et les conditions \(H_i\supset C\), \(\chi_i|_C=\mu\) sont restitués. Le texte hérité réduisait cette donnée à une « identité dans \(G_\mu\) ». L'inclusion \(H_i\subset G_\mu\), nécessaire pour l'induction affichée, demeure implicite comme dans le témoin ; elle n'est pas ajoutée silencieusement au texte de l'auteur.
\item[E30] \textbf{Source p.~20 : racines de l'unité et twist.} La notation \(\Z/n(1)(\overline K)\), la base séparablement close \(\overline K\), le \(\widehat\Z\)-module \(\Q/\Z\) et l'identification \(\frac{n}{m}\otimes x\mapsto\bar x^n\) sont rétablis. La phrase héritée sur \(K=\C\) ne figure pas dans le témoin et effaçait la construction de \(\Q/\Z(1)\). Le texte conserve la notation abrégée \(\Z/n(1)\) du témoin, sans choisir de racine primitive pour identifier ce module à \(\Z/n\Z\).
\item[E31] \textbf{Source pp.~21--22 : congruence de l'inertie.} Le signe de congruence et la cible imprimée \(t(\sigma)\,v(x)\in\Q/\Z(1)(\overline k)\simeq\overline k^*\) sont rétablis. Le produit avec \(v(x)\) est l'action sur le module tordu; il n'est pas remplacé silencieusement par une multiplication ordinaire dans \(\overline k^*\).
\item[E32] \textbf{Source p.~22 : suite exacte du groupe de Weil local.} Le diagramme omis dans la transcription antérieure est restitué. Les deux lignes ont pour noyau \(I\), et la flèche de droite est l'inclusion \(\Z\hookrightarrow\widehat\Z\) induite par \(W(\overline k/k)\subset\Gal(\overline k/k)\). Les §§2.2.5--2.4 et les diagrammes (2.3.1)--(2.3.2), relus sur les pages 23--24, restent inchangés.
\item[E33] \textbf{Source p.~24 : premier objet inférieur de (2.3.1). Émendation héritée, maintenant signalée.} Le témoin imprime \(W(\overline K/L)\) en bas à gauche, comme en haut. L'édition héritée porte \(W(\overline K/K)\), exigé par la flèche verticale d'inclusion et par la ligne inférieure \(W(\overline K/K)^{\ab}\simeq K^*\). Cette correction mathématique ne doit pas être prise pour une lecture littérale du témoin.
\item[E34] \textbf{Source p.~25 : abélianisation dans le diagramme de réciprocité globale. Émendation héritée, maintenant signalée.} Le témoin imprime \(W(\overline K/K)\) à droite de la ligne supérieure, tandis que l'édition antérieure portait déjà \(W(\overline K/K)^{\ab}\). Le domaine \(W(\overline K_v/K_v)^{\ab}\), la flèche vers le groupe des classes d'idèles et l'isomorphisme affiché juste au-dessus imposent ce quotient abélien. Les définitions, constantes et formules des pages 26--28 sont conservées après nouvelle comparaison directe.
\item[E35] \textbf{Source p.~29 : domaine des valeurs de (3.5.2).} Le témoin écrit \(L(V)=Z(V;1)\ (\in E^*\cup\{\infty\})\). La transcription héritée écrivait à tort \(E\in E^*\cup\{\infty\}\), proposition sans type correct. Le sujet de la parenthèse est bien la valeur \(L(V)\); le champ \(E\) n'est pas modifié.
\item[E36] \textbf{Source pp.~31--32 : caractère additif global.} Le témoin demande un caractère \emph{non trivial} de \(\mathbb A/K\), de composantes locales \(\psi_v\). La transcription héritée disait « trivial ». La nouvelle lecture conserve la trivialité sur \(K\), qui est implicite dans le fait de descendre au quotient, mais n'annule pas le caractère sur \(\mathbb A/K\).
\item[E37] \textbf{Source pp.~33--34 : argument d'induction de 3.12.} La relation sur les facteurs \(L\) est replacée dans la démonstration de (C), avec le renvoi imprimé à 2.10. (C) énonce la relation sur \(\varepsilon\); (D) vient après la démonstration, et son signe \(+\) est conservé. L'ancienne édition déplaçait la relation \(L\) dans (C), puis remplaçait la justification imprimée par une allusion à la norme et au transfert.
\pagebreak
\item[E38] \textbf{Source p.~36 : identité intermédiaire du lemme 4.3.} L'affichage \(\varepsilon(v,\psi(ax))=\det(v_0)(a)\varepsilon(v,\psi)\), qui relie (3.3.3) au transfert, est rétabli. Il avait été remplacé par une phrase générale dans les deux éditions héritées.
\item[E39] \textbf{Source p.~36 : domaine réduit et relation archimédienne de 4.4.} Le témoin porte \(b:R^0_+\to R^0\) et \(x_s=(\overline K^*,\omega_s)-[\omega_s]-[\chi^{-1}\omega_{s+1}]\). L'édition emploie \(R^+_0\) pour \(R^0_+\), conformément à sa notation en 1.10; les qualifications de dimension zéro sont rétablies. La barre sur \(\overline K^*\) et le symbole \(x_s\) étaient perdus. Pour \(K\simeq\mathbb R\), ce terme correspond à l'extension complexe de degré deux, et non au groupe multiplicatif réel.
\item[E40] \textbf{Source p.~39 : étape d'induction du lemme 4.5.}
L'égalité intermédiaire avec
\(\Ind_J^I([\nu]|_J)\), omise dans les deux textes hérités, est rétablie.
Le témoin répète le numéro (4.5.1) pour ce calcul; le numéro éditorial
(4.5.1') déjà présent dans l'édition est conservé pour distinguer les deux affichages.

\item[E41] \textbf{Source p.~38 : précisions de notation conservées.}
Dans (4.5.4), le témoin imprime \(\dim V-V^I\); l'édition conserve
\(\dim V-\dim V^I\), qui explicite la dimension des invariants.
La différente est imprimée \(\mathfrak D_{L/K}\), après le discriminant
\(\delta_{L/K_{nr}}\); cette différence d'indices est signalée, non corrigée silencieusement.
Le symbole \(K\) est aussi réutilisé pour un sous-groupe d'inertie; cette notation du témoin est conservée.

\item[E42] \textbf{Source p.~42 : catégorie et induction.}
La catégorie \(\mathcal R_\tau\) est à nouveau distinguée de son groupe
de Grothendieck \(R_\tau\), comme dans le témoin.
La dernière phrase de 4.10 concerne les représentations de \(K\)-type donné,
sans les restreindre à \(W(\overline K/L)\) : l'induction change précisément ce groupe.
La définition du type, ses exposants et les formules (4.10.1)--(4.10.2) sont conservés.
\item[E43] \textbf{Source p.~45 : domaine d'intégration du lemme 4.15.}
Le témoin imprime \(\int_{\mathcal O_L}dx\), le volume de l'anneau des entiers.
L'édition héritée portait \(\int_L dx\), volume infini du corps local.
Le domaine imprimé est rétabli dans les deux éditions; ce n'est pas une émendation mathématique.
\item[E44] \textbf{Source p.~44 : places et groupes de décomposition.}
« Ne se décompose pas » signifie ici qu'il n'y a qu'une place au-dessus de \(v\),
comme l'indique l'égalité des groupes de Galois du texte. L'anglais explicite ce sens
dans les lemmes 4.12 et 4.13, sans imposer l'absence de ramification.
Le groupe imprimé \(\Gal(K'_w/K_w)\), après la représentation de \(\Gal(K'/L)\),
est conservé; l'action de cette dernière sur un sous-groupe ne doit pas être
confondue avec une extension automatique à ce groupe plus grand.
\item[E45] \textbf{Source p.~50 : nom de la constante en 5.8.}
Le témoin écrit \(\varepsilon\), non \(\varepsilon^{-1}\), avant l'équation de Tate.
Cette lecture est rétablie; la constante de l'affichage (5.8.1) est inchangée.
\item[E46] \textbf{Source p.~53 : corps du caractère à la fin de 5.16.}
Le témoin écrit \(K^*\), non le \(K'^*\) hérité, dans la phrase qui compare 5.14 et 5.15.
La lecture imprimée est rétablie dans les deux éditions.
\item[E47] \textbf{Source p.~54 : lettre de l'anneau des coefficients.}
La source emploie \(\Lambda\); l'édition héritée le renomme \(A\) dans la section 6.
Ce renommage est conservé et signalé; les hypothèses et les formules ne changent pas.
\item[E48] \textbf{Source p.~55 : raccourcis de 6.3.}
L'intervalle imprimé \(1\le i\le n\) est conservé, mais omet le facteur d'ordre \(1\).
La phrase sur les racines \(p^n\)-ièmes n'affirme pas ici de borne uniforme pour tout le caractère.
Le témoin dit ensuite « \(n(\psi)\) est non trivial »; l'édition conserve l'émendation héritée « \(\psi\) est non trivial ».
\item[E49] \textbf{Source p.~55 : valeurs du caractère en 6.4.}
Le témoin imprime \(\Lambda\); l'édition conserve \(A^*\), comme en 6.3 et dans le théorème 6.5.
C'est une précision héritée, non une lecture littérale de cette phrase.
\item[E50] \textbf{Source p.~59 : nom du topos en 7.3.}
Le témoin nomme le topos de Weil \(Y_{w\mathrm{\acute et}}\); \(Y_{\mathrm{Weil}}\) n'est qu'un renommage.
Le facteur imprimé du produit fibré est \(Y_{\mathrm{\acute et}}\), non \(\overline Y_{\mathrm{\acute et}}\): cette barre héritée est supprimée.
La première mention imprimée du topos classifiant porte \(B\mathbb Z\), tandis que la base du produit porte \(B\widehat{\mathbb Z}\); l'édition conserve cette dernière notation dans les deux mentions.
\item[E51] \textbf{Source p.~60 : fibre du second déterminant en (7.6.1).}
Le témoin porte \(\mathcal G_{\overline k_v}\) dans les deux déterminants.
Le second \(\mathcal G_{\overline K_v}\) hérité est remplacé par cette fibre fermée dans les deux éditions.
\item[E52] \textbf{Source p.~63 : indices du complexe en 7.11.5.}
Le déterminant imprimé après extension des scalaires porte \(H_0(K)\) et \(H_1(K)\), tandis que la page suivante emploie \(H^i(K)\).
L'édition conserve les indices cohomologiques \(H^0(K^\bullet)\) et \(H^1(K^\bullet)\), pour le complexe en degrés \(0\) et \(1\).
\item[E53] \textbf{Source p.~64 : deux notations dans l'ouverture de la preuve de 7.11.}
La phrase imprimée de convergence porte \(Z(j_*V,t)\), mais l'affichage qui suit porte \(L(V,s)=Z(j_!V,q^{-s})\).
L'édition conserve \(j_*\) dans les deux occurrences. Cette précision héritée n'est pas une lecture littérale du second affichage.
\item[E54] \textbf{Source p.~69 : changement de Frobenius dans le lemme 8.4.3. Émendation mathématique.}
Le témoin imprime \(n\,t_\ell(\varpi)\) dans la première exponentielle et \((1-q^{-1})^{-1}t_\ell(\varpi)\) dans la seconde, sans \(N\) dans ces deux affichages. Pour le changement à droite \(F\mapsto F\varpi\), le Frobenius géométrique donne \(\rho'(F)N\rho'(F)^{-1}=q^{-1}N\). L'édition écrit donc le coefficient \((1-q^{-n})/(q-1)\) dans la première exponentielle, et \((q-1)^{-1}\) dans le conjugateur, avec \(N\) explicite. Pour \(n=1\), le premier coefficient est \(q^{-1}\). C'est une correction éditoriale déduite de ces conventions, non une lecture littérale ni un erratum attribué à l'auteur.
\item[E55] \textbf{Source p.~72 : inverse du déterminant dans 8.12. Émendation héritée, maintenant signalée.}
Le dernier affichage imprimé de \(Z(V,t)\) ne porte pas l'exposant \(-1\), alors que (8.12.2) définit le facteur local par le déterminant inverse. L'édition conserve l'inverse déjà présent dans les deux textes hérités. Le remplacement local de \(t^d\) par \(t\) dans ce dernier affichage est conservé comme dans le témoin; aucune uniformisation silencieuse du paramètre n'est effectuée.
\item[E56] \textbf{Source p.~78 : ensembles d'indices dans la preuve de 9.8. Restauration du témoin.}
Les deux produits du membre gauche de l'équation récrite portent \(v\notin S\) dans le témoin, et non \(v\in S\) comme dans l'édition héritée. Ces deux ensembles d'indices sont restaurés. Les produits du membre droit restent sur \(S\).
\item[E57] \textbf{Source p.~78 : torsion de Tate dans cette même équation. Émendation mathématique.}
Le témoin imprime \(V_\lambda^*\) dans le second produit gauche et \(V_\lambda\omega_1\) au dénominateur droit. L'édition adopte respectivement \(V_\lambda^*\omega_1\) et \(V_\lambda\). En effet, pour \(D=V_\lambda^*\omega_1\), décomposer \(Z(V_\lambda,t)=\varepsilon_{\mathrm{Gr}}Z(D,t^{-1})\) en produits sur \(S\) et hors de \(S\) donne exactement l'équation corrigée. Cette correction est déduite de (9.1.1), non présentée comme une lecture littérale ou un erratum de l'auteur. L'exposant \(-1\) du second produit gauche est conservé.
\item[E58] \textbf{Source p.~79 : élément du groupe de Weil dans la fin de la preuve de 9.8. Émendation de notation.}
La phrase quantifie sur tout \(F\in W(\overline K_v/K_v)\), hors de \(I_v\), mais le déterminant imprimé utilise \(F_v\). L'édition emploie \(F\) dans les deux déterminants, conformément à ce quantificateur et à la conclusion sur les caractères hors de l'inertie. Le témoin reste inchangé; cette correction n'est pas une lecture littérale du déterminant imprimé.
\item[E59] \textbf{Source p.~96 : date de la référence [12]. Lecture du témoin restaurée.}
Le témoin imprime 1958, et non 1968 comme dans l'édition héritée. La date imprimée est conservée ici pour distinguer transcription et correction bibliographique. La référence [7] de J. Tate, \emph{Classes d'isogénie des variétés abéliennes sur un corps fini}, Séminaire Bourbaki, exposé 352 (1968/69), indique W. A. Benjamin, New York, 1968. Cette attestation explique la correction possible ; elle n'est pas présentée comme la lecture du témoin.
\end{description}
\par\medskip
\noindent\textbf{Témoin.} P. Deligne, \emph{Les constantes des équations fonctionnelles des fonctions L},
\emph{Modular functions of one variable II}, Lecture Notes in Mathematics 349 (1973), pp.~501--597.
Le PDF de comparaison de 97 pages est fourni dans le dossier SOURCE.
Son SHA-256 commence par \texttt{F151A9A6B71CCFED}; l'identité complète et l'absence des octets
exactement sélectionnés (empreinte \texttt{B03F483C4EEC...}) sont consignées dans les contrôles.

% END EDITORIAL APPARATUS
\endgroup
\end{document}
